% =====================================================================
%  Poul Martin Møller, TANKER OVER MULIGHEDEN AF BEVISER FOR MENNESKETS UDØDELIGHED, med Hensyn til den nyeste derhen
%  hørende Literatur.  Maanedsskrift for Litteratur 17 (Kjøbenhavn 1837): I, pp. 1-72; II („Fortsat og sluttet“),
%  pp. 422-453.  Fraktur, with antiqua for Latin, French and some German words and for numerals in the notes.
%
%  A DIPLOMATIC transcription of the first printing: spelling, capitals, punctuation, word division at the page turn
%  (marked %), misprints (marked PRINTED AS IS), emphasis and notes as printed.  \opage{N} = the page of MfL 17.
%
%  GENERATED; DO NOT EDIT BY HAND.  Rebuild from the repo root:
%    python3 ebook-method/udoedelighed/vote2.py && python3 ebook-method/udoedelighed/fromreads.py
%  (the readings ebook-method/udoedelighed/reads/reads-R1..R22.txt, the corrections vote-fixes.tsv, this template).
%
%  SOURCE.  Google Books scan of MfL 17, 104 PDF pages, 600 ppi bitonal, with Google's OCR as a text layer:
%  ~/bibliotek/Møller, Poul Martin/1837-maanedsskrift-17-udoedelighed.pdf (SCANS.tsv, sha256 0a62304a...).  PAGE MAP:
%  MfL pp. 1-72 = PDF 1-72; MfL pp. 422-453 = PDF 73-104.  The scan holds only the essay's pages.  P. 72 also carries
%  the opening of the next piece in the journal (a review of Brunius, Lunds Domkyrka); it is not transcribed.
%  This text REPLACES an earlier hybrid (Art. I from the modernised Skrifter i Udvalg 1930, Art. II from this print
%  respelled 1930-style), which also lacked the note on p. 62.
%
%  METHOD (2026-10-18; TRANSCRIPTION-PLAYBOOK §00 and §6, the begrebet-noiagtighed route).
%  1. Every page was transcribed BLIND at the image by 22 readers (5 pages each, the page cut into bands of ~7 lines,
%     600 ppi), who saw neither OCR.  Two CONTROLS were planted in each batch (a word's image replaced by another word's):
%     44 of 44 were read as planted, then undone.
%  2. Every word of the readings (30,437) was voted against Google's layer, tesseract (Fraktur) and a collated draft;
%     a word that agreed with none of them, or that the reader marked uncertain, went to the image: 418 questions, with
%     6 possible omissions, the 38 rule decisions of the draft that the 1856 resetting (Efterladte Skrifter V) contradicts,
%     the note count of p. 431 and 38 questions of type face, were cut as labelled crops (both halves of a word divided at
%     a line end) and read blind by six further readers, with two hidden controls a sheet (172 of 174 right; 10 more
%     that a script first scored wrong were controls cut without the word's second half).
%  3. What a crop reader and a machine reading both supported over the page reader was corrected (vote-fixes.tsv,
%     93 lines including the control undo lines); what the page reader and a machine reading supported stood;
%     the rest was settled by a further crop round or, where the readings already decided it, by review (reasons in
%     .parts/vote/settled-hand.tsv).  The reader errors found were mostly MODERNISATIONS (the print's skee read ske,
%     Beskeed read Besked, Eenhed read Enhed, Conseqventser read Consequentser; vist read viist, saadant read saabant):
%     1837 spelling is kept throughout.
%  4. A last round of 15 crops (11 words, 4 faces; 6 hidden controls, 5 right) settled what the readings still left open.
%  ACCURACY.  Every word of the text was read at the image by a blind reader and agrees with at least one independent
%  machine reading, or was read again at the image by a second blind reader; band controls 44/44, crop-sheet controls
%  177/180.  Nine sites the readings left doubtful were decided by Hans Halvorson at the image (2026-10-18): Leibnitz
%  p. 7, tumlet p. 15, skader p. 24, Voxne p. 56, Udtryk p. 62 (the d worn), Beviser: p. 442 (a smear at the colon),
%  der p. 446 (no mark: a speck), skoleret p. 450, and the faces.  He also found the crop reader's „Leibniz“ (five
%  places) a reader fault: the print has Leibnitz with the tz ligature at all six.
%  What this method cannot see: a misprint that the page reader, the crop reader and both machines all read alike as
%  the right word (the floor measured on comparable books is about one in a thousand words).
%  CONVENTIONS.  Type faces: PROPER NAMES are set in Schwabacher throughout (Frederich, Ferdinand, Fichte, Leibnitz …
%  -- the Schwabacher F of every F-name against the Fraktur F of ordinary words; also Göschel, Marheineckes) and are NOT
%  marked here (Stewart italicises them).  The rule holds in Møller's own text, not inside the German quotation in the
%  note on p. 451 („Freder. Christ. Sibbern“ there is Fraktur).  EMPHASIS is Schwabacher on an ordinary word -> \textbf;
%  the essay has two: „virkelig Tilværelse“ p. 70 and „noget Udødeligt“ p. 425 (across the line break).  This was
%  established by checking, against the 1837 image (judged by Hans Halvorson), every word letterspaced in the 1856
%  resetting (Efterladte Skrifter V; ebook-method/udoedelighed/emph.py: 12 hits, 2 real) and every emphatic italic in
%  Jensen & Stewart's translation (Brill 2022, made from the 1837 print: 3, 2 real).  Weight in this scan is ink and
%  unreliable: the readers' „heavy“ tags were checked and were plain.
%  Letterspacing -> \emph: Magdeburg, Breslau (p. 1), Tanken (p. 449).  Antiqua -> \textit (numerals and § in antiqua
%  are not marked).  The Fraktur I/J is one sort: I before a consonant and as the word „I“, J before a vowel.  The
%  quotations (Weiße, Fichte, Göschel) are set in the body type in the print and are run in with their quotation marks.
%  Footnotes are marked *), **) per page, as printed.  The ø rule (CLAUDE.md): an o-form that is no Danish word is set
%  with ø; here 2 such changes.  The owner's pencil and ink underlining and marginal numerals (pp. 48, 52-54, 57,
%  59-61, 67, 440) are not print.  Sheet signatures and running heads are not transcribed.
% =====================================================================
\documentclass[12pt,a4paper,openany]{book}
\usepackage[utf8]{inputenc}
\usepackage[T1]{fontenc}
\usepackage{libertinus}
\usepackage{graphicx}
\DeclareUnicodeCharacter{0254}{\reflectbox{c}}  % ɔ (U+0254) = reversed c
\usepackage[danish]{babel}
\usepackage[protrusion=true,expansion=true]{microtype}
\usepackage{setspace}
\usepackage[margin=1.3in]{geometry}
\usepackage{fancyhdr}
\usepackage{hyperref}
\hypersetup{pdftitle={Tanker over Muligheden af Beviser for Menneskets Udødelighed}, pdfauthor={Poul Martin Møller}, hidelinks}
\setlength{\headheight}{15pt}
\pagestyle{fancy}
\fancyhf{}
\fancyhead[LE,RO]{\thepage}
\fancyhead[C]{\textit{P. M. Møller: Om Menneskets Udødelighed (1837)}}
\renewcommand{\thefootnote}{\ifcase\value{footnote}\or *)\or **)\or ***)\else
  \arabic{footnote})\fi}
\setstretch{1.3}
\usepackage{marginnote}
\newcommand{\opage}[1]{\marginnote{\footnotesize\textit{[#1]}}}

\begin{document}

% ===== ARTICLE I: MfL 17 pp. 1--72 =====
% --- p. 1 ---
\opage{1}\setcounter{footnote}{0}%

\begin{center}{\Large Tanker over Muligheden af Beviser for\\
Menneskets Udødelighed, med Hensyn\\
til den nyeste derhen hørende Literatur.}\end{center}

\begin{center}\rule{3em}{0.4pt}\end{center}

\begin{center}I.\end{center}

Den Mening, at en Regjering handler rigtig, naar den tilstæder
Udbredelsen af Skrifter i videnskabelig Form, selv om disse gaae
ud paa en aabenbar Bestridelse af den christelige Religions Hovedsætninger, er nylig bleven understøttet ved en mærkelig Kjendsgjerning. Frederich Richter fra \emph{Magdeburg} har nemlig,
(hvad sandsynligviis er en stor Deel af Maanedskriftets Læsere
bekjendt), for et Par Aar siden i \emph{Breslau} udgivet et Skrift,
der har gjort stor Opsigt, fordi det har til Formaal at godtgjøre,
at den christelige Tro paa den personlige Udødelighed er en Fordom, der nu har varet sin længste Tid. Han mener nemlig, at
denne Sandhed forlængst har været en afgjort Sag blandt
Mænd, der staae-paa et videnskabeligt Standpunct, og at Tiden
nu er kommen, da hiin Overtro ogsaa hos Mængden bør forsvinde for den consequent fremskridende Protestantismes Lys.
Hans Afhandling røber en alvorlig, varm Overbeviisning. Uagtet
han nærer en til Sværmeri grændsende Forvisning om, at han
har faaet et høiere Kald til at optræde mod en dybt rodfæstet
% --- p. 2 ---
\opage{2}\setcounter{footnote}{0}%
Vildfarelse, vogter han sig dog i Almindelighed for den Kaadhed,
som en saadan Mening let kunde friste til, og holder sig saa temmelig indenfor den videnskabelige Anstændigheds Grændser. Dette
troer jeg kun en Fanatiker vil nægte ham. Et saadant Skrift
maa naturligviis forekomme mangen frygtsom Tilhænger af Statsreligionen saare farligt, og han vilde maaskee ansee det for raadeligt, at Udbredelsen af en saadan Vranglære blev hindret ved Foranstaltninger fra Øvrighedens Side. Alligevel har Udfaldet viist,
at den preussiske Regjering har gjort vel i, at den ikke ved udvortes Midler har hindret fri Meningsyttring om hiint væsentlige
Punct i Troeslæren. Richters Yttringer have nemlig fremkaldt
en stor Mængde Modskrifter og have været Aarsag i, at den
hemmelighedsfulde Taushed over dette Problem, som ikke blot
blandt Philosopher, men selv blandt Theologer i lang Tid har
hørt til Dagens Orden, er bleven afbrudt ved aabenhjertige Meddelelser derom. Der vil neppe for Tiden komme nogen Dogmatik for Lyset, der, saaledes som Marheineckes og Fleres, kaster
et forsigtigt Slør over det hele Capitel om Udødeligheden. I det
Mindste maa nu hver Philosoph og Theolog være saa kraftig opfordret til uforbeholdne Meddelelser, at forblommet Yttring herefter vil blive betragtet som ligegjældende med aabenbar Fornægtelse. Neppe vil man da see et saadant Særsyn gjentage sig
som det, at en Philosoph indhyller Sagen i et saa dunkelt Skolesprog, at der selv blandt hans i Systemets Mysterier meest indviede Tilhængere skal kunne føres en literær Strid om, hvad han
egentlig har lært om et af Menneskehedens vigtigste Spørgsmaal:
saaledes som Tilfældet nylig har været i den hegelske Skole. Den
afdøde Mußmann, „den høitelskede, der havde ligget ved sin Mesters Bryst,“ erklærede det i een af sine sidste literære Yttringer
for en væsentlig Mangel ved Systemet, at den personlige Udødelighed var fornægtet deri; medens Göschel, der af Mange be%
% --- p. 3 ---
\opage{3}\setcounter{footnote}{0}%
tragtes som Skolens Formand, har paataget sig udførlig at bevise den modsatte Paastand.

Det gaaer med Udødelighedslæren som med enhver problematisk Sætning om det Usandselige eller Oversandselige. Man
kan umulig komme til et fast Resultat, hvis man ikke i det
Mindste forsøgsviis gjennemfører Thesis eller Antithesis, og derved enten paa hiin eller denne Side støder paa forkastelige Følgesætninger. Men hvorledes kan da Opløsningen af Menneskelivets
Gaader kraftigere fremmes, end naar eet Individ føler Kald
til at forfægte det ene Led af Dilemmaet og forfølge det i dets
vigtigste Consequentser, og en Anden derved bliver opfordret til
at stille sig paa den modsatte Side? Og dette vil altid blive Tilfældet, da enhver ret afgjørende Benægtelse eller Bekræftelse af
omtvistede Sætninger bestandig ægger til Modsigelse. Derved
fremkommer en Personification af Begrebet, hvorunder det vel ikke,
som mange af Hegels Eftersnakkere sige, bevæger sig selv, men
dog træder frem som det bevægende Princip i flere tænkende
Subjecter, til, hvilke dets forskjellige Momenter ere ligesom fordeelte. Denne i Sandhed levende Dialectik, i hvilken Menneskehedens Aand taler med sig selv, gjentager i større og stærkere
Stiil den samme Virksomhed, som foretages af en alsidig eenlig
Tænker, der veier de modsatte Verdensanskuelser mod hinanden.
Hiin Verdensdialectik, som Ingen kan standse, hindres og besværes meget i sin Udvikling, naar Stemmen fra de Partier, der
stride mod Statsreligionens Hovedsætninger, ved udvortes Magt
aldeles undertrykkes\footnote{Efter Forordningen af 26 September 1799 § 5 er det her i Landet forbudt at udgive Skrifter, som have til Formaal, at „nedbryde Læren om den menneskelige Sjæls Udødelighed.“ --- Dersom
Hegels Skrifter vare udkomne hertillands, vilde det have faldet
endogsaa Høiesteret meget vanskeligt at afgjøre, om deri fandtes en
Overtrædelse af hiint Forbud eller ikke.}. Statsreligionen taber ogsaa derved nødven%
% --- p. 4 ---
\opage{4}\setcounter{footnote}{0}%
digviis meget af sin Anseelse, naar Mængden erfarer, at Indvendinger mod dens Lærdom ei tør komme for Dagens Lys. Ogsaa
i en anden Henseende lider Religionens Sag ved saadanne Forholdsregler; thi den Opposition, der ikke har noget andet Element end den mundtlige Conversation, yttrer sig der med desto
mere sammentrængt Kraft og desto mere tøileløs Eensidighed,
ligesom visse Planter med større Frodighed udbrede deres Rod i
Jordens mørke Skjød, naar de hindres i at udskyde deres Top
efter deres naturlige Drift i Lysets Regioner. Den blot mundtlige Polemik faaer endogsaa paa sin Side den store Fordeel, at
den ei behøver at udsætte sig for nogen Kritik, som den ei
skjøtter om.

Jeg haaber nedenfor at lægge klart for Dagen, at det efter
Sagens særegne Natur gjælder meer om Udødelighedslæren, end
om nogen anden thetisk Sætning, at den, naar den engang er
rokket i sin Grundvold, ikke kan komme til det Herredom i den
menneskelige Bevidsthed, som tilkommer den, med mindre Negationen kommer til at udvikle sig frit i alle sine Følger, for at det
kan vise sig, om Menneskeheden formaaer at leve med den Verdensanskuelse, som den fører til. Men jeg vil ikke fragaae, at
jeg nærer nogen Tvivl, om den Grundanskuelse, til hvis Forsvar jeg
paa disse Blade giver et middelbart Bidrag, ved nogen Anstrængelse for nærværende Tid kan komme til sin Ret. Det er meget
muligt, at Negationen ei endnu har naaet det fulde Maal, som
den skal naae for at det kan blive ret øiensynligt, at det Øde,
den fører med sig, ei er den Sphære, hvori den menneskelige Aand
skal bevæge sig. Men det er ligemeget; de som ei kunne dele
Forstyrrelsens særegne Lyst, stræbe dog at bygge sig en Ark,
hvorpaa de kunne bjerge sig i Haab om bedre Tider.

% --- p. 5 ---
\opage{5}\setcounter{footnote}{0}%

\begin{center}II.\end{center}

Naar man fra et polemisk Standpunct søger at forklare den
% PRINTED AS IS: V020 pv1: crop, layer and tesseract agree -- a MISPRINT, printed as is
Verdensanskuelse, som hvis Modstander man fremstaaaer, er det
nutildags meget sædvanligt, at man, i Mangel af grundigt Begreb derom, søger en psychologisk Forklaringsgrund i sine Modstanderes moralske Charakteer. Saaledes er det nu en staaende
Tactik for Udødelighedslærens Angribere at forklare sig denne formeentlige Overtro af-en egoistisk Opholdelsesdrift, som ikke tillader Individet med ædel Selvopofrelse at finde sig i, at det som
endeligt Væsen er et forbigaaende Led i Rækken af det Absolutes
Tilsyneladelser. Men en Retorsion af samme Vaabenart ligger
da meget nær for Udødelighedslærens Forsvarere: de tilskrive nemlig Modstanderne en Bevidsthed om, at deres Individualitet, som
de saa redebon give Afkald paa, kun er af ringe Betydning og
fattigt Indhold. Deslige pragmatiske Forklaringsmaader, som
søge Grunden til enhver theoretisk Antagelse blot i menneskelig
Skrøbelighed, have kun ringe videnskabeligt Værd. De sande
Grunde saavel til Positionen som-Negationen maa søges ved et
dybere Blik i Menneskets Culturhistorie.

Foreløbig maa det tilstaaes Fornægternes Hob, at der til
Grund for deres Synsmaade til alle Tider hos Mange har ligget et sundt og kraftigt Liv i Momentet, om hvis Realitet de
have følt sig saa forvissede, at de høist ugjærne have henvendt
deres Opmærksamhed paa den abstracte Forestilling om et Tilkommende, der kun ved Hjelp af subjektive Phantasier kan vinde bestemt Indhold. Men hos dygtige Mennesker finder man tit
denne Uvillighed til at overlade sig til ørkesløse Phantasier forenet
med en fast Overbeviisning om personlig Udødelighed. Af Göthes Værker og de mange Mindeskrifter om hans Levnet seer man,
at han just forbandt disse to Stemninger, Tro paa Menneskets
% --- p. 6 ---
\opage{6}\setcounter{footnote}{0}%
Liv efter Døden og Ulyst til at gruble meget over dets Beskaffenhed.

For at forklare sig Negationens Almindelighed i vor Tid
maa man angive en særegen Grund dertil, foruden den almindelige
Ulyst til at vende Tanken fra den bekiendte Realitet til det Hiinsidige, der har saa faa sikkre Bestemmelser. Vi forudsætte det
nemlig her som en afgjort Kjendsgjerning, at udtrykkelig Benægtelse af Udødeligheden nutildags er meer almindelig end til nogen
foregaaende Tid i de christelige Aarhundreder. Ikke blot fra
Tydskland og Frankrig, men endog fra England, hører man jævnlige Beretninger derom, og flere europæiske Statistikere have heri
søgt en af de væsentligste Forklaringsgrunde til Nutidens talrige
Selvmord\footnote{I Paris udgjøre de i Maaneden September 1836 forefaldne Selvmord, forsaavidt de ere komne til offentlig Kundskab, 66, siger og
skriver sex og tredsindstyve.}. Hvis Aabenhjertighed ikke var noget hindret, deels
ved Statsindretninger, deels ved en agtværdig Sky for at rokke
ved en af „Folketroens“ Grundpiller, vilde det snart vise sig, hvor
Faa af det ringe Tal, der har nogen selvstændig Mening, da
vilde paatage sig den ældgamle Læres Forsvar. Endeel af Udødelighedslærens nyeste Forsvarere have heller ikke taget i Betænkning
at vedgaae, at de tilforn have næret en ganske modsat Overbeviisning.

Den dybere Grund til alle disse Phænomener maa søges i
Videnskabens Udviklingsgang. De to Erkjendelser, den af Eenheden og den af Mangfoldigheden i Tilværelsen, der i lang Tid
bestode fredelig ved Siden af hinanden, uden at man havde nogen
klar Bevidsthed om deres Modsætning, kom tilsidst i voldsomt
Sammenstød, hvilket Stød har frembragt alle Hovedforvandlinger
i den nyere Tids Philosophie; thi det er egentlig een og samme
Modsætning, der snart har viist sig som Strid mellem Idealisme
% --- p. 7 ---
\opage{7}\setcounter{footnote}{0}%
og Realisme, snart mellem Frihedslære og Nødvendighedslære,
snart mellem Rationalisme og Supranaturalisme o. s. v. Da
Spinoza med kraftig Eensidighed havde gjort Eenheden gjældende i den nyere Tids Videnskab, fremkaldtes derved stor Forargelse, og Leibnitz, der var gjennemtrængt af Overbeviisning om
Mangfoldighedens Realitet, opstillede derimod sin Monadelære.
Men da han tillige følte det Savn af Eenheden, som medførtes
ved denne Hypothes, tilføiede han sin Lære om den forudbestemte
Harmonie. Denne Theorie var rigtignok en udvortes Nødhjælp,
et Slags mechanisk Forbindelse af Eenhed og Mangfoldighed: det
var kort sagt et Palliativ, som ei i Tidens Længde kunde staae
sin Prøve; men allerede Leibnitz selv havde dog et klart Begreb
om, hvilken væsentlig Splid i Bevidstheden, han derved vilde
bringe til Forlig, siden han paa Titelen af flere Værker kalder sig
Opfinder af Læren om den forudbestemte Harmonie.

Udtrykket „Strid mellem Erkjendelsen af Eenheden og af
Mangfoldigheden,“ som jeg her har betjent mig af for at betegne
det væsentligste Stridspunct i den nyere Tids Philosophie, er ikke
forstaaeligt for den, som slet ikke har noget historisk Bekjendtskab
med de Bevægelser, som her i faa Linier sammenfattes under eet
Overblik. Ogsaa maa det som et altfor ubestemt Udtryk, der
blot paa en abstract Maade omfatter en stor Mangfoldighed af
Problemer, findes utilfredsstillende for den, som veed og forstaaer,
hvad her sigtes til. Derfor vil jeg med concretere Benævnelser
betegne det særegne Slags Eenhed og Mangfoldighed, hvorom her
er Tale. Her tænkes paa Modsætningen mellem en forkastelig
Pantheisme paa den ene Side og en cras Theisme paa den anden
Side, under hvis Yderlighed man danner sig en Gud med en
hedensk begrændset Personlighed, der blot staaer ved Siden af
andre endelige Personligheder.

Den Verdensanskuelse, som laae til Grund for Leibnitz's
Monadelære, gjentager sig i sine Grundtræk meget tit i Videnska%
% --- p. 8 ---
\opage{8}\setcounter{footnote}{0}%
bens Historie. Den har været stærkt udtalt i det nittende Aarhundrede, ja selv i det sidste Decennium. For ei at opfylde
disse Blade med en Overflødighed af literære Notitser, som om
det var en latinsk Disputats, vil jeg blot vise hen til Herbarts
System, der i Hovedsagen, som det her gjælder om, ei er forskjellig fra Leibnitz's Monadelære. I begge antages nemlig det sande
Reale i Tilværelsen for at bestaae i en Mangfoldighed af enkelte
Væsner eller Substantser, hver med sin særegne Beskaffenhed,
hvilke Væsner selv ei kunne blive Gjenstande for sandselig Anskuelse, men hvis Forbindelser for os antage Skinnet af en sand
Tilværelse. Efter denne mechaniske Theorie af Tilværelsen, i Følge
hvilken Intet i Sandhed bliver til eller forgaaer, men Alt bestaaer
af uforgiængelige enkelte Væsner, ved hvis Forbindelser og Adskillelser mangfoldige Skikkelser fremstaae, er vistnok en Slags Udødelighed sikkret disse enkelte Væsner, hvilke naturligviis ogsaa
kunne træde frem som menneskelige Sjæle.

Ved en saadan Construction af Tilværelsen, ved Antagelsen
af en Mangfoldighed af Substantser, som man for at være ganske consequent maa tillægge Evighed og saadan Selvstændighed, at
der i Systemets Aand, ikke kan antages noget uendeligt Væsen,
tabes den Eenhed, som er Fornuftens ufravigelige Fordring, aldeles af Sigte. Vil man da tilfredsstille den religiøse Bevidstheds
Fordringer, levnes der ingen anden Udvei end en saadan kunstig
Mechanisme, som Leibnitz's forudbestemte Harmonie. Dette bliver især tydeligt ved Herbarts Fremstilling, da han paalægger sig
mindre Tvang af Hensyn til den sædvanlige Religionslære, og de
enkelte Substantser i hans System derfor ingenlunde udgives for
Skabninger af Gud. Efter denne Synsmaade bliver som sagt
Eenheden fortrængt ved den kraftige Bevidsthed om Mangfoldighedens Selvstændighed. Med jo større Følgerigtighed man gjennemfører en saadan Eensidighed, desto mere bringes dens Utænkelighed til klar Bevidsthed; thi vil man fastholde den begyndte
% --- p. 9 ---
\opage{9}\setcounter{footnote}{0}%
Tankegang, da kan intet Forhold tænkes mellem de for sig existerende Væsner: de kunne ei engang i deres Selvstændighed udelukke hverandre uden et fælleds Medium, hvori de kunne udelukke
hverandre. For Enhver, der tænker, at de ere absolut adskilte
fra hverandre, foresvæver i det samme den dunkle Tanke om en
Sammenholdning, altsaa et Forhold og det Modsatte af absolut
Adskillelse.

I den endnu meest gangbare Form for den populære Religionslære udvikles Menneskesiælens Natur saaledes, at deri findes
stærke Spor af den gamle Lære, at Sjælen er en Substants, og
denne Lære har vistnok deels paa en umiddelbar, deels paa en
middelbar Maade faaet Bestyrkelse ved Leibnitzes philosophiske
Arbeider. Det er endnu langt fra, at Læren om Sjælen, som en
for sig existerende Ting, der er bragt ind i det menneskelige Legeme og maaskee har sit Ophold paa et bestemt Sted deri, fuldkommen er fortrængt af Tidsalderen. Der gives endnu dem, der
ikke finde det saa urimeligt at tænke sig den samme Sjæl indsluttet i en Organisme, der stod paa et andet Udviklingstrin, f.
Ex. et Dyrs eller en Plantes. I al Fald finde de den Sætning ret antagelig, at Sjælen kan existere adskilt fra den legemlige Skikkelse og netop i denne Tilstand føre den reneste og fuldkomneste Tilværelse. I slige Forestillingsmaader findes en Levning af den ældgamle metaphysiske Sætning, at Sjælen er en
fra Legemet forskjellig Substants. I de sædvanlige Religionsbøger ophæves rigtignok denne Sætning tildeels ved Læren, at Sjælene skabes og opholdes af Gud: de betragtes følgelig efter denne
Synsmaade ikke længer som evige Væsner, der bestaae ved sig
selv. Men hvo, der kjender noget til den Tankegang, hvorved
Folk uden nogen philosophisk Dannelse forvisse sig om Sjælens
Udødelighed, vil have overbeviist sig om, at Læren om den enkelte
Substants,, der ei kan opløses, endnu spøger i Manges Hoveder.
Det skjønnes ogsaa klart af den Uvillie, hvormed de modtage ad%
% --- p. 10 ---
\opage{10}\setcounter{footnote}{0}%
skillige Udtryk af den nyere Philosophie, f. Ex. den, at Sjælen
er Legemets Form. Ved en saadan Lære forekommer det dem,
at en physisk Mulighed af Sjælens Tilintetgjørelse paa en for
Troen farlig Maade faaer Indpas i Videnskaben. Selv en af
vore berømte Videnskabsmænd erklærer, (i et Brev, som findes
aftrykt i Jens Møllers theologiske Bibliothek), at han har
Prof. Sibbern mistænkt for en hemmelig Pantheisme, fordi
han i Dennes Psychologie ei finder Sjælen fremstillet som en
Substants for sig.

Uagtet Religionslærerne aldrig reent have fornægtet enten
Eenheden eller Mangfoldigheden i Tilværelsen, havde dog det formeentlig orthodoxe Partie, især af Skræk over Spinozas stærke
Pantheisme i lang Tid lagt for stor Vægt paa Individualiteten,
og Tilværelsens Eenhed blev derved tit paa en sørgelig Maade
tabt af Sigte. Men den fortrængte Bevidsthed om Eenheden i
Alt er med Kjæmpekraft indtraadt i sine Rettigheder og truer
stærkt med utilbørlig at tilbagetrænge Tanken om den endelige
Mangfoldigheds især Personlighedernes Realitet. Læren om Sjælene som selvstændige Ting eller Substantser, er forsvunden hos
Enhver, der blot har et lidet Anstrøg af Tidens philosophiske Cultur, men derved har ogsaa Udødelighedslæren tabt sin gamle Grundvold. Heri ligger den dybere Grund til, at ikke blot Tvivl men
udtrykkelig Fornægtelse af Læren om en personlig Udødelighed nuomstunder udbreder sig meer og meer. Reactionen mod den Splid
i Bevidstheden, under hvilken man havde glemt, at Gud ikke blot
omfatter, men ogsaa gjennemtrænger Alt, har ført til en raa,
forkastelig Pantheisme, med hvilken intet sandt Udødelighedsbegreb
kan bestaae. Men denne crasse Pantheisme er for Tiden ikke blot den
meest udbredte Synsmaade blandt Philosopher, men den er i en meget
eenfoldig og let fattelig Form trængt igjennem til de Fleste, som
have nogen exoterisk Kundskab om Philosophiens Tilstand. Den
herskende Pantheisme ligger nuomstunder enhver begyndende Tæn%
% --- p. 11 ---
\opage{11}\setcounter{footnote}{0}%
ker saa nær, at man ikke har mange scholastiske Kunstord behov
for at bringe den tillive hos ham, men ved nogle faa Formularer bliver den meget fatteligere for ham, end den christelige Troes
Artikler. Hos vore Digtere findes undertiden i faa Linier en
sammentrængt Fremstilling af denne Verdensanskuelse, som ingen
ærlig Pantheist vil undslaae sig for at erkjende for et kort Indbegreb af sin Læres Hovedsætninger. En gammel Videnskabsmand, som havde gjort et mangeaarigt Studium af Fortidens og
Nutidens philosophiske Værker har engang for nogle Aar-siden i
mit og Fleres Paahør erklæret om en saadan compendiøs Naturphilosophie, som etsteds i Ingemanns sorte Riddere lægges en ond
Aand i Munden, at dette egentlig var den sande Videnskabeligheds Resultater.

Den crasse Pantheisme, der nu er saa vidt udbredt, har
mange Theologer imod deres Villie ikke lidet fremmet ved deres
aldeles modsatte bornerede Theisme, som egentlig fortjener Navn
af Afguderi; thi de kunde med større Ret give deres begrændsede Gud Navn af Zeus eller Odin, end af Christendommen sande
Gud. Og det er især deres blinde Skræk for Pantheismen, der
har bragt dem paa saadan Afvei; det er de Orthodoxes egen ubevidste Rationalisme, der hindrer dem-i efter Tidens Tarv at oversætte sig de gamle Ord: „Du skal ikke gjøre dig noget Billede
eller nogen Lignelse.“ Det var værdt at vide, hvad mangen Theolog vilde svare, naar han paa sin Samvittighed blev spurgt, om
han ikke vilde ønske, at de bekjendte Ord, „I ham leve og ere
og røres vi,“ slet ikke forekom i den hellige Skrift. Han kan
umulig finde sig deri uden ved modificerende Omskrivninger at
betage dem deres pantheistiske Characteer. Ved at bestride Pantheismen med haardnakket Eensidighed, og gaae ud fra den Forudsætning, at den ikke indeholder Andet end reen Vildfarelse,
an man hverken give Troen paa en personlig Gud eller paa
Menneskets Udødelighed sin tabte Livlighed. Vil man blot be%
% --- p. 12 ---
\opage{12}\setcounter{footnote}{0}%
stride Pantheismen, og slet ikke indrømme den mindste Ret, da
føres man nødvendig til et System, efter hvilket Tilværelsen
smuldrer hen i en Sværm af hjælpeløse Atomer. Man bliver
tidlig eller sildig nødsaget til at ophøre med at bruge Ordet Pantheisme blot som et Øgenavn, thi det betegner en væsentlig Side
af den sande Erkjendelse. Pantheismen er en theologisk Forskole,
der maa gjennemgaaes af Enhver, der skal komme til en sand religiøs Indsigt. Det er alene paa Pantheismens Grund, at et
Tempel for den sande personlige Gud nuomstunder kan opreises.

Mangen fornuftig Theolog vil uden Tvivl gjøre den Bemærkning, at den Side af det theologiske Begreb, som jeg her betegner
med Navnet Pantheisme, har sin Plads i ethvert populært Religionsforedrag under Form af Begreberne Allestedsnærværelse, Almagt osv.
Det er heller ikke min Paastand, at dette Moment i nogen Religionslære, der har mindste Skin af Sandhed reentud kan fornægtes; men jeg paastaaer, at det kan være tilstede i en altfor
undertrykt Tilstand; at den Tænkemaade, som udtales i en Række
Sætninger, kan være grundfalsk, uden at en enkelt kan udhæves
for sig som aldeles urigtig. Man kan ogsaa med god Grund give
den fremførte Udsættelse den Vending, at Begreberne Almagt, Allestedsnærværelse o. s. v. i de betegnede Religionsforedrag et fremstilles i deres rette levende Betydning. Jeg paastaaer fremdeles,
at de der eensidig bestride Pantheismen, frakjende det modsatte
Partie deres gode Ret til at betjene sig af visse Talemaader, der
have ligesaa velgrundet Sandhed, som de gandske modsatte, der
blandt ældre Theologer have Gang og Gyldighed. Det høres f. Ex.
nødig af Mange, at man bruger saadanne Udtryk som, „at Gud
er Alt“; „at der er Intet udenfor Gud“ o. s. v.. Mange høre
heller de ligesaa rigtige Sætninger, at Verden er udenfor Gud;
Gud udenfor Verden o. s. v.. Saadanne eensidige Almeensætninger om Gud og Verden lade sig med lige Ret saavel bekræfte
som benægte; men Sandheden udtales kun da fuldstændig, naar
% --- p. 13 ---
\opage{13}\setcounter{footnote}{0}%
man baade viser Thesis og Antithesis i sin halve Sandhed.
Naar man da paatager sig den pantheistiske Sides Ret, strider
man for den ubetingede Frihed i den religiøse Sandheds Udtalelse
og gjør Indsigelse imod, at en vis Kreds af traditionelle Sætninger skulle betragtes som priviligerede fremfor andre, der netop have
samme Gyldighed, ja endogsaa ligesaa god Hjemmel i rettroende
Lærdomme. Det er heller ikke Meningen, at gamle Theologer
skulle vænne sig til en Talebrug, der var deres nuværende ganske
modsat: Characteren af deres Foredrag skal kun efterhaanden forandres, og det vil skee, uden at de selv blive det vaer. Denne
successive Forandring vil nødvendigviis frembringes af en Tidsaand,
% PRINTED AS IS: V039 q: crop, layer, tesseract -- a MISPRINT, as printed
der er Pantheismen saa overdreven gunstig. Snart ville kun enkelte udlevede eller halstarrige Charakteermennesker blandt Theologerne staae tilbage som forstenede Mindesmærker fra en undergaaet
Begrebsverden. De, hos hvem den opholdende og den fremskridende Kraft staae i den rette Ligevægt, ville mod deres Villie ved
Verdensdialectikens Magt henrives til den umærkelige Overgang,
som det populære Foredrag saa meget begunstiger. Som sædvanligt i lignende Tilfælde skulle Faa erkjende, Færre tilstaae, at de
have befundet sig paa en Afvei. Eenhedens sande Triumph skal
ikke bestaae deri, at den falske Theisme fældes ved et afgjørende
Slag, men dens Skikkelse skal langsomt svinde hen ligesom, Auroras
Ægtemand.

Det var altsaa, for kortelig at gjentage det, den Verdensanskuelse der lægger for stor Vægt paa Endeligheden eller Mangfoldigheden, der ovenfor blev betragtet som særdeles gunstig for
Udødelighedslæren. I Philosophien saae vi, at derved consequent
uddannede sig en Lære om mange evige Substantser, hvorved Tilværelsens Eenhed blev aldeles benægtet. I den sædvanlige Form
af Religionslæren er denne Retning ikke slet saa eensidig, fordi
man dog ikke der pleier at tillægge noget Slags Individer absolut Substantialitet. Men desuagtet bliver Tilværelsens Eenhed
% --- p. 14 ---
\opage{14}\setcounter{footnote}{0}%
siælden udtalt efter dens fulde Betydning. Man fremfører en
foregiven personlig Gud, som i Grunden tillige er et endeligt Væsen, hvis Realitet indskrænkes ved den Realitet, han sætter udenfor sig, saa at der fragaaer ham selv saa megen Realitet, som han
tillægger sine Skabninger. Men den Philosophie, der saa eensidig
gik ud paa at hævde Individernes Selvstændighed, at Tanken om
en Alt indbegribende Gud blev sat i Skygge, har nu overlevet sig
selv og givet Plads for en modsat Eensidighed. En af de nærmeste Følger deraf er, at Læren om Sjælenes mechaniske Forbindelse med Legemet og Evne til at existere for sig, er saa godt som
udryddet af Tidsalderens speculative Philosophie. Selv uden for
Philosophien er den næsten reent forsvunden hos den yngre Slægt,
for saavidt som den blot besidder nogen videnskabelig Dannelse.
Den gamle Tankebygning, hvori Læren om Sjælens Udødelighed
havde sin faste Plads, er saaledes styrtet sammen, og ved Forkastelsen af dens Hovedsætninger, med hvilke man har indbildt sig
at Udødelighedslæren maatte staae og falde, har denne Lære selv i
Manges Øine tabt sin Gyldighed. Heri findes den væsentlige
Grund til de vidt udbredte Tvivl om Menneskets Udødelighed:
det gamle System, hvormed dette Begreb var bragt i Harmonie,
er gaaet til Grunde, og det nye System er ikke saa fuldstændig
udført, at det deri har fundet sin rette Begrundelse. Endogsaa
de, der med fuldkommen Overbeviisning om, at Alt er Eet, bevare
deres Udødelighedstro urokket, bruge med en Slags Usikkerhed de
gamle Formularer. De finde f. Ex. af flere Grunde Betænkelighed ved at bruge Udtrykket Sjælens Udødelighed; men uagtet
deres mangfoldige Forsøg paa at bringe nye Betegnelser i Gang,
f. Ex. Personlighedens Idee, den individuelle Vedvaren o. s. v.
har hidtil ikke noget Udtryk kunnet opnaae det gamles almindelige
Gyldighed, og derved, at Nogle paany ere vendte tilbage til den
gamle Talebrug, er Tingenes Forvirring kun bleven forøget.

Fordi Udødelighedstroen har mistet sin gamle Rod, har den
% --- p. 15 ---
\opage{15}\setcounter{footnote}{0}%
ogsaa mistet sin Kraft til at frembringe den levende Veltalenhed,
som fødes af aandig Anskuelse og tiltaler Sindet. Saalænge som
Overbeviisningen om en oversandselig Sandhed tilveiebringes ved
en harmonisk Forening af Tænkning og Anskuelse, er den en uudtømmelig Kilde til friske Vendinger i Talen, til et levende Ord,
som er dens Affødning og en væsentlig Form for dens Liv og
Tilværelse. Naar denne Skabelse hører op, har den Overbeviisning, der satte den i Bevægelse, upaatvivlelig tabt sin Grundighed. Nuomstunder er ogsaa denne Form for Udødelighedslærens
Liv i Følge de alt udviklede Forhold i en meget sygelig Tilstand.
I geistlige Taler hører man tit en treven og, som man fristes
til at troe, halv skamfuld Opramsen af forslidte Formularer om
dette Kapitel. Selv Mænd, der over andre Gjenstande tale med
mærkelig Interesse, synes ikke uden et Slags Forlegenhed at gjentage de forbrugte Tirader om Udødeligheden. I Poesien, hvor
Anskuelsen skulde have Overvægt over Begrebet, staaer det endnu
slettere til med den lyriske Udtalelse af Haabet om en ny Tilværelse. Der har man saa tit tumlet med de gamle Riim paa
Taare --- Baare --- sees igjen --- Himmelen, at Liigvers netop
ved deres Trivialitet skaffe Publicum en kjær Underholdning. Naar
man een eller to Gange i Decenniet læser et hjerteligt Digt i
Anledning af et Dødsfald, da savner man gjerne al Berøring af
Udødelighedshaabet eller træffer et mat Vers deri, som ligner et
ufrivilligt Offer paa Anstændighedens Alter.\footnote{Saaledes saae man for nogle Aar siden et Digt af A. Ussing ved
hans Elskedes Grav, hvilket bar meer end almindeligt Præg af Sandhed og Alvor, men hvori intet Ord fandtes om Udødelighed.}

Det vil næppe lykkes ret mange for Tiden at giengive de
af Vedtægt hjemlede Talemaader om det Hiinsidige hele deres fordums Livlighed; thi den Verdensanskuelse, som var det Element,
hvori de hørte hjemme, har for Tiden tabt sin Betydning. Men
% --- p. 16 ---
\opage{16}\setcounter{footnote}{0}%
der vil komme den Tid, da fuld Overbeviisning vil gjenfødes paa
en anden Vei, og da vil den foryngede Anskuelse atter blive frugtbar paa livfulde, gribende Former i Digt og Tale.  Men Tidens
Gud vender sig langsomt og i den Mellemtid, da hans gamle Ansigt er bortvendt og det andet endnu ikke ret er kommet frem,
hersker der nødvendigviis et Tusmørke.  Under saadanne Katastropher er den gamle Form for Erkjendelsen betydningsløs og den,
der skal afløse den, er endnu ikke ret kommen til Verden, men
svæver for os som en halv usynlig Skikkelse med ubestemte Omrids.

Resultatet er da, for nu efter de forudskikkede Bemærkninger
til Slutning at gjentage det med større Bestemthed. og Korthed,
at man saavel i den nu fortrængte Metaphysik som i dens folkelige Efterklang kun dunkelt havde erkjendt Eenheden i Alt, og i
det en forandret Verdensanskuelse gjør sig gjældende, vakler den
gamle Udødelighedstro, som meest beroede derpaa, at man antog
det for en ontologisk eller physisk Umulighed, at de for sig existerende Sjæle kunde forgaae.  I den Tænkemaade, som afløser hiin, er
en modsat Eensidighed herskende, da Forvisningen om, at eet Liv
rører sig i Alt, er udartet til en raa Pantheisme, hvorefter den
endelige Personlighed næsten har tabt sin Realitet.  Under denne
Tænkemaade savner Udødelighedsbegrebet deels en videnskabelig
Plads i det nye Begrebssystem, deels almeengyldige Former, hvori
dets Indhold kan blive livlig udtalt.  Men dette Beviis og disse
Former kan ikke fuldføres, førend saavel Erkjendelsen af Eenheden,
som den af Mangfoldigheden, kommer til sin fulde Ret eller, for
at betegne det med et indholdsrigere Udtryk, førend baade Guds
og Menneskets Personlighed, erkjendes i sin fulde Betydning.

\begin{center}\rule{3em}{0.4pt}\end{center}

\begin{center}III.\end{center}

Følelsen af Udødelighedstroens almindelige Usikkerhed fremkalder den tit gjentagne Fordring til Videnskaben, at den skal til%
% --- p. 17 ---
\opage{17}\setcounter{footnote}{0}%
veiebringe et „strængt Beviis“ for Menneskets Udødelighed, men
en meget uklar Forestilling om Naturen af et saadant Beviis
ligger hos Mange til Grund for denne Fordring.  Der gives ikke
Faa, der synes at forudsætte, at en let overskuelig Tankegang kan
blive dem meddeelt udenfra, hvorved den Art af Overbeviisning,
som en mathematisk Demonstration medfører, kan blive dem bibragt om Udødelighedens Realitet.  Heri ligger for det Første
den besynderlige Forudsætning, at en saadan Lære kan begribes
for sig, udreven af sin organiske Sammenhæng med andre Hovedsætninger i et fuldstændigt System.  Fremdeles forudsættes, at
et tænkende Menneske kan paa en forunderlig Maade være fremmed i sin egen Begrebsverden.  En Anden skulde nemlig kunne
finde en uomtvistelig Antagelse deri, hvoraf Overbeviisningen om
Menneskets Udødelighed med stræng Nødvendighed lod sig udlede.  Den Vished, som Præmisserne til den Ende skulde have,
findes hos de Spørgendes Fleertal ei om en eneste væsentlig
Sandhed;  og om der endogsaa fandtes en saadan urokkelig Antagelse hos dem, vilde det være et stort Særsyn, om det var undgaaet deres egen Opmærksomhed, at Udødelighedslæren stod i en
saa nødvendig Sammenhæng dermed.

Hvad endeel Mennesker, der kræve hine strænge Beviser for
Udødeligheden egentlig ønske, have de sjælden selv ret tydeliggjort
sig: i Grunden begjære de Fortsættelsen af den samme Tilværelse,
som de have ført her paa Jorden.  Hiin Stræben efter Demonstrationer for Udødeligheden udspringer af en urimelig Lyst til at
beholde sit daglige Livs Prosa ubeskaaret, og midt i den at erholde
en sandselig-forstandig Vished om det Oversandseliges Realitet.
Ingen er modtagelig for Overbeviisningen om en virkelig tilværende oversandselig Verden, uden den som i forbigaaende Momenter er i Stand til at føle den sande Frihed-fra alle jordiske
% PRINTED AS IS: V064 q: crop and layer -- a MISPRINT, as printed
Interesser: ieg mener hermed ikke Ironiens Frihed, men den som
netop fører til det rette Liv i Verden.

% --- p. 18 ---
\opage{18}\setcounter{footnote}{0}%
Ved at giennemtænke det forunderlige Selvbedrag, hvori
Mange befinde sig, der spørge om et Beviis for Menneskets Udødelighed, faaer jeg Lyst til at fremkomme med en kort Fortælling,
i hvorvel jeg forudseer, at den efter kræsne Læseres Dom ikke
kan bestaae med den omhandlede Gjenstands Værdighed;  men jeg
troer, den er skikket til at anskueliggjøre den Stemning med Hensyn til Udødeligheden, som her tales om, og den kan med god
Grund fore Navn af Digtning og Sandhed.

\begin{center}\rule{3em}{0.4pt}\end{center}

Fortælleren besøgte for nogen Tid siden en ugift Bogholder
af sit Bekjendtskab, som han traf i en høirøstet Samtale med en
theologisk Candidat.  Strax ved Indtrædelsen modtog Bogholderen
ham med følgende Ord: Det er godt du kommer;  du skal nu
dømme mellem Ferdinand og mig.  Han har for et Øieblik siden
afhentet et herligt Skrift om Sjælens Udødelighed, som han bringer lige fra Bogbinderen.  Nu nægter han mig reentud at laane
mig det et Par Timers Tid.  Er det en venskabelig Adfærd? ---
Theologen svarede: Du veed selv Julius, du handler meget ilde
med Bøger, og jeg sætter stor Priis paa at samle mig et velconditioneret Bibliothek. --- Den Snak kommer du altid med, min
gode Ferdinand, fordi du engang saae mig skjære en Rulle hollandsk Tobak paa Bindet af en Nummerbog. -Jeg husker meget
godt, det var Morderen i Fjeldhulen eller det blodige Gjenfærd.
Men idag lovede jeg dig jo udtrykkelig, at jeg vilde behandle Bogen peent. --- „Jo, det skal jeg svare til.  Du har jo selv sagt,
du vilde læse den, mens du spiste dit Smørrebrød.“ --- Deri har
du Ret, svarede Bogholderen;  men om jeg ogsaa havde sat en
lille Plet paa den uskatteerlige Bog, mon da Skaden havde været
saa ubodelig?  Jeg tænker, hele Herligheden har kostet dig henved
tre Mark.  Men Grunden til det Venskabsstykke, at du afslog
mig Bogen, er oprigtig talt din Egenkjærlighed.  Det skulde al%
% --- p. 19 ---
\opage{19}\setcounter{footnote}{0}%
drig have fortrudt mig, naar det ikke havde været et Skrift af
denne Art, du nægtede mig.  Jeg har ikke tænkt synderlig paa
Sjælens Udødelighed, siden den Tid, jeg som Pog havde min Lectje
derom i Campes Ledetraad, og fik, som jeg endnu tydelig husker,
mit Ug.x derfor.  Men jeg er i Grunden et religiøst Menneske:
jeg erkjender fuldkommen, det er Umagen værdt, at faae den Sag
bragt paa det Rene, og jeg har i flere Aar foresat mig engang
ved Leilighed at læse et eller andet godt Skrift over denne Materie.  Idag havde jeg netop gode Stunder, fordi jeg sidder og
venter paa et Par gode Venner, med hvem jeg Kl. 1 skal kjøre
til Bellevue for at spise friske Torsk.  Hør Julius, laan mig saa
% PRINTED AS IS: V074 q: crop, layer, tesseract -- a MISPRINT, as printed
Bogen;  nu har ieg eengang sat mig i Hovedet at læse den idag.
Jeg er saa opsat derpaa, at jeg gjærne kunde falde paa at sende
Bud efter den i en Boglade, skjøndt du selv veed, jeg ikke er af
dem, der smider Penge bort. --- „Det bliver, som jeg har sagt,
jeg gjør det ikke.“

Deres Stemmer afløste hinanden saa hurtig, at Fortælleren,
der var opfordret til at være deres Voldgiftsmand, ikke kunde
komme til at lade sin Røst høre.  Omsider fik han dog Leilighed til at gjøre den Bemærkning, at Klokken alt var bleven halv
eet, og at Striden altsaa var unyttig, da Bogholderen umulig i
den Tid der var tilbage, kunde blive færdig med Bogens Læsning.
--- Hilledød, sagde Bogholderen, er Klokken saa mange og jeg har
endnu ikke begyndt min Paaklædning.  Jeg havde dog grumme
gjærne løbet den berømte Afhandling igjennem.  Jeg indseer fuldkommen Vigtigheden af, at man engang imellem beskjæftiger sig
med Religionens Hovedsandheder.  Saalænge Verden gaaer En
med, tænker man gjærne kun lidet derpaa, men --- man sige
hvad man vil --- Ingen kan vide sig sikker for at friste saa
haarde Slag af Skjæbnen, at det er godt, han med Fortrøstning
kan opløfte sig ved Tanken om et Hisset.  Jeg udgiver mig ikke
for en stærk Aand: det vilde være mig saare kjært, om Nogen
% --- p. 20 ---
\opage{20}\setcounter{footnote}{0}%
kunde føre mig et ret haandgribeligt Beviis for Sjælens Udødelighed.  Dersom jeg nu ikke en anden Gang faaer Leilighed til
at læse et grundigt Skrift derom, har jeg, reentud sagt, ikke noget
ret Fast at holde mig til i mørke Timer.  Men hvis er da
egentlig Skylden, min gode Ferdinand?  \textit{A propos}, siden vi nu
% PRINTED AS IS: V075 q: crop, layer, tesseract -- a MISPRINT for bragt, as printed
eengang have hragt denne Sag paa Bane, kunde du, som er
Theolog, gjærne i al Korthed foredrage mig de bedste Beviser for
Sjælens Udødelighed, mens jeg stryger min Ragekniv og tager
mit Skjæg af;  men kom mig ikke for nær: jeg skjærer mig let.
--- Det vil falde mig meget vanskeligt, sagde Theologen, at meddele dig Tanker, som virkelig vare dig nye og endda kunde tydeliggjøres for dig i et Øieblik.  --- Der har vi den gamle Snak,
Ferdinand.  Naar du vil betjene dig af et forstaaeligt Sprog,
hvorfor skulde jeg da ikke fatte det ligesaa godt som du?  Men vi
kjende nok den Laugsaand, som har besat baade dig og alle dine
Hrr. Colleger med den latinske og græske Dannelse.  I ville
bruge eders eget Kragemaal for at beholde Noget for Eder selv.
Men dette Kastevæsen vil neppe bestaae, naar de liberalere Begreber trænge ret igjennem: nu skal Videnskaben populariseres;  det
kræver Tidsaanden.  Du skal blot vogte dig for at bruge de
Kunstord, hvorunder de Lærde skjule deres Tanker for Lægfolk og
sige mig den simple Mening deraf paa godt bredt Dansk.  Men
skynd dig noget;  jeg er bange vi har Vognen her paa Øieblikket.
--- Ferdinand svarede noget forlegen: Der er dem, som beraabe
sig paa en Anelse, et Slags Forudfølelse af en tilkommende høiere
Tilværelse.  --- „Dette Punct kan du springe forbi;  jeg troer
hverken paa Anelser eller paa Sympathie: det er netop et strængt,
uimodsigeligt Beviis jeg forlanger af dig.“ --- „Desuden erindrer
du vel ogsaa, man pleier at beraabe sig paa det Misforhold, der
ofte finder Sted mellem de Godes og Slettes Livskaar.“ --- „Ja
det er godt nok: de Gode maae tit savne Løn for deres Selvopofrelse og Ædelmodighed: men hvo har egentlig lovet dem Be%
% --- p. 21 ---
\opage{21}\setcounter{footnote}{0}%
taling for deres Gjerninger? Hvo siger mig, at Dyden skal belønnes og Lasten straffes. Du maa da heller ikke troe, du kan
komme til mig med et saadant løst Raisonnement, som en Landsbypræst kan byde sine Confirmander.“ --- „Man har ogsaa omtrent samme Beviis, som vi sædvanlig i Skolen kalde det moralske, i en mere videnskabelig Form af Fichte.“ --- „Spring
Formen over og beret mig Indholdet i det Sprog, som din Moder har lært dig: det forstaaer baade du og jeg. Men der seer
jeg Wienervognen. Nu har du her spildt os Tiden med at gjentage gamle Sætninger af Balles Lærebog, i Stedet for at føre
mig et stringent Beviis. Farvel!“

\begin{center}\rule{3em}{0.4pt}\end{center}

Ovenstaaende Fortælling fremstiller blandt Andet den Erfaring, at Alle, der uden at kjende noget til Videnskabens Natur,
udbede sig en Demonstration af Udødelighedslærens Realitet, ikkeret vide hvad de begiære; først, fordi Sætninger om oversandselige Ting ikke kunne demonstreres efter den i Mathematiken brugelige Methode, da denne Videnskabs Læresætninger kun have
Gyldighed under bestemte Forudsætninger. Dernæst fordi det
undgaaer deres Opmærksomhed, at et evident Beviis for en isoleret Sætning om en oversandselig Realitet, ikke kan føres for
den, i hvis Bevidsthed der savnes hvilkensomhelst Antagelse om
oversandselige Ting, hvortil Beviset kunde knyttes. Det tilstrækkelige Beviis skulde da være et heelt philosophisk System. Naturligviis kan en levende Udvikling af en philosophisk Anskuelse
vække den samme eller endogsaa en klarere Anskuelse hos en
Anden. Men med den Tale, der ofte høres om Savnet af
strænge Beviser, f. Ex. for Udødeligheden, menes noget gandske
Andet. Denne Tale føres især af dem, der kun kjende den formelle Visheds Natur, og den beroer paa Mangel af klar Indsigt i den Sandhed, at deres saakaldte strænge Beviser blot udtale
% --- p. 22 ---
\opage{22}\setcounter{footnote}{0}%
hypothetiske Sætninger. Denne Overbeviisning, som eensidig Uddannelse af Forstanden fører til at kræve, hvor den ei kan opnaaes, er kun den formelle Overbevisning om, at en vis Forsætning medfører en vis Følge. Men af Mangel paa faste Forsætninger bringes aldrig en hypothetisk Slutning tilveie, som har
denne Art af Vished. Det er let at gjøre Dette klart ved Sætninger af Mathematiken eller den sædvanlige Logik, da det ved
nøiere Betragtning vil sees, at de allesammen i Grunden ere hypothetiske. Det er vist, at en Triangels Vinkler ere = 2 R.;
men dermed er jo ikke sagt, at der gives nogen Triangel, og man
veed kun da Noget om den virkelige Tilværelse, dersom man tillige veed, at Triangler komme for deri, men Dette læres jo ei i
Mathematiken. I Mathematiken bevæger man sig egentlig i en
ideal Sphære, hvori der føres strænge Beviser for blot hypothetiske Sandheder. Men, naar man foredrager Læresætninger om
det Oversandselige, da dreier man sig ikke blot i Tanker om Forhold, der passé paa en vis Art af Gjenstande under Forudsætning
af, at disse Gjenstande selv paa en anden Maade ere givne, men
man udgiver sig for at vide Noget om den oversandselige Realitet selv.

Sætte vi nu, at En i sin Bevidsthedskreds har Antagelser
om oversandselige Ting og Antagelser, der ledsages af fuld Overbeviisning, da bliver det desuagtet forunderligt, at han haaber, at
en Anden kan hjælpe ham til at finde det Punct deri, hvoraf
Udødelighedslæren med Evidents kan udledes. En saadan tiltroer
sig selv et ringe Bekjendtskab med sin egen Overbeviisning, da
han troer, at en Anden bedre end han selv kan gjøre Opdagelser
i hans egen Forestillingskreds.

Jeg troer, at Ingen, som har havt Leilighed til at erfare,
hvorledes Mange af hans Samtidige tænke om den omhandlede
Sag, vil paastaae, at den Vildfarelse, som jeg her taler imod, er
greben af Luften. Den forekommer ikke blot hos Mange, som
% --- p. 23 ---
\opage{23}\setcounter{footnote}{0}%
søge Beviis for Udødeligheden, og see sig om efter en kyndig Mand,
der i en Samtale kan give dem Beskeed om Sagen, men den
findes ogsaa jævnlig hos Videnskabsmænd, som undertiden virkelig
opfinde saadanne Beviser for Hvermand. Deres Methode er meget fattelig. De antage en Deel af Religionens Hovedsætninger
for afgjorte, men Læren om Udødelighed for usikker, og søge da
i det forudsatte Forraad af Sætninger, en eller anden, hvortil
Udødelighedslæren kan hæftes som en uimodsigelig Følgesætning;
ligesom en Mand i et Værelse seer sig om efter et Søm, hvorpaa han kan hænge sin Hat. Et træffende Exempel paa det
Slags Beviser, som jeg mener, haves i det af Moses Mendelsohn opfundne, der efter sin hele Betydning kan fremføres i et
Øieblik. „Mennesket er udødeligt; thi hvis det ikke var udødeligt, da var Gud ei almægtig, men Mennesket kunde unddrage
sig fra hans Magt ved at dræbe sig selv.“

Vi sammenfatte kortelig, hvad der i det sidst foregaaende er
sagt. Vi have paastaaet, at der ikke kan føres et apodictisk Beviis for Menneskets Udødelighed. Den uimodsigelige Vished, som
her menes, gjælder kun for hypothetiske Sandheder, som de der
udtales i Mathematiken og den formelle Logik. Den formelle
Vished, her tales om, kjendes meget vel af simple Folk. Selv
den meest Udannede refererer dette Begreb til Mathemathiken,
som Ingen er aldeles ubekiendt med. Deraf Forsikkringer som
den: „Det er saa vist, som at to og to er fire.“

Men kan der ikke ved den nys opdagede Methode i Philosophien, den speculativ dialectiske, føres et Beviis for Menneskets
Udødelighed, som kanskee ikke medfører en mathemathisk Vished,
men dog en fuldkommen Overbeviisning af en anden Art? ---
Det er just et af de væsentligste Stridspunkter i den nyere Philosophie, som det opkastede Spørgsmaal bringer os i Berøring
med. Det rette Svar paa dette Spørgsmaal kan ikke gives uden
ved en tilfredsstillende Udvikling af Forholdet imellem den aprio%
% --- p. 24 ---
\opage{24}\setcounter{footnote}{0}%
riske og den empiriske Erkjendelse. En saadan kan naturligviis
umulig meddeles i denne Afhandling, men desuagtet haaber jeg
ved nogle aphoristiske Vink at kaste noget Lys paa Problemet.
Men selv for at være istand hertil maa jeg bruge den Frihed,
som saa Mange nuomstunder med altfor liden Betænkelighed tage
sig, den Frihed at bruge et og andet af de philosophiske Kunstord,
som Sædvanen har bragt i Omløb selv i det friere Foredrag.
Det er i sig selv en mislig Fremgangsmaade, da saadanne Skoleord sjælden kunne føre til nogen tydelig Forstaaelse udenfor en
stræng videnskabelig Sammenhæng. Nogle Kunstord, som ikke
høre til Sprogets egentlige Stamme, have dog faaet en saadan
Gang og Gyldiged i det friere Foredrag, at de for de fleste dannede Læsere medføre en Slags umiddelbar Betydning, som i visse
Tankeforbindelser kan være tilstrækkelig. I ethvert Tilfælde skader det ikke at anstille et Forsøg, da her kun er levnet Valget
mellem enten aldeles at forbigaae det paa Bane bragte Spørgsmaal, eller at bruge en halv scholastisk Fremstilling, som mulig
forekommer nogle Læsere uklar. Jeg vil da gaae saaledes frem,
at jeg først med et Par Ord vil omtale Udødelighedslærens Tilstand i den speculativ-dialectiske eller hegelske Philosophie, og derpaa med Forbigaaelse af det udførlige Beviis, som jeg haaber at
kunne meddele ved en anden Leilighed, kortelig udvikle min Overbeviisning om Forholdet imellem Erkjendelsens aprioriske og aposterioriske Deel i Almindelighed; thi derved bliver Spørgsmaalet
om et apriorisk Beviis for Udødeligheden ogsaa besvaret.

For det Første tør jeg med den fuldkomneste Vished paastaae, at der i Hegels Værker ikke findes et eneste Sted, hvor et
saadant Begreb om Udødelighed forsvares, som Christendommen lærer, og som Nutidens dybsindigste Philosopher med saa stor Ret
have søgt at gjøre gjældende paany. Hermed paastaaes ei, at han
udtrykkelig benægter Udødelighedslæren, hvilket falskelig foregives i
flere tydske Skrifter, hvor man endogsaa tit har anført Tirader
% --- p. 25 ---
\opage{25}\setcounter{footnote}{0}%
af Hegels Logik for at godtgjøre denne Beskyldning. Man har
nemlig fremdraget Yttringer, hvor det abstracte Hiinsidiges Begreb omtales som en umoden metaphysisk Tanke, men det er ved
en grov Misforstaaelse, at man ved dette Udtryk har tænkt paa
de christelige Forestillinger om en en tilkommende høiere Tilværelse.
Der er ligesaa lidet paa saadanne Steder Tale derom, som der
overhovedet i Logiken efter hele dens Plan kan være Tale om saa
concrete Begreber som Himmel og Udødelighed. Men desuagtet
er det let for Enhver, som nogenlunde forstaaer at læse imellem
Linierne i Hegels Skrifter, at komme til fuldkommen Vished om,
at denne Philosoph antager Begrebet om den personlige Udødelighed for en Forestilling uden al Realitet. Altsaa kan der ikke
være Tale om noget Beviis for Udødeligheden ved speculativ Dialectik, dersom man spørger Methodens Opdager selv derom.
(Til dette Punct vender jeg igjen tilbage paa et senere Sted i
denne Afhandling.) Men uagtet Hegel selv har næret saadan
Overbeviisning, har dog En af hans Tilhængere paataget sig at
bevise paa en reen apriorisk Maade, af blot abstracte Begreber,
at Menneskets Udødelighed er uimodsigelig vis.

Inden jeg nu giver det belovede Svar paa det ovenfor opkastede Spørgsmaal om Forholdet mellem apriorisk og aposteriorisk
Erkjendelse, bemærker jeg først, at jeg her med Flid har brugt
det gamle Udtryk \textit{a priori}. Her var nemlig Valget givet mellem tre forskjellige Kunstord, det fra den kritiske Philosophie meest
bekjendte Udtryk „en apriorisk Deduction,“ Naturphilosophiens dermed nær beslægtede Yndlingsord, „en philosophisk Construction“
og Hegelianernes „immanente Tænkning,“ et Ord som mangen
ung tydsk Forfatter nuomstunder sætter en stor Ære i at kunne bruge
nogenlunde rigtig. Men jeg tager her ikke i Betænkning at bruge
det gamle Udtryk et apriorisk Beviis, som deels er det meest bekjendte, deels bedst udtrykker det, jeg her ønsker at betegne dermed. Kun anmodes Læseren om at glemme den Biforestilling
% --- p. 26 ---
\opage{26}\setcounter{footnote}{0}%
om Erkjendelsens Subjectivitet, som for en Kjender af den kritiske
Philosophie let klæber ved hiint Udtryk.

Kan man da (for at komme til Sagen) i Philosophien bevise Noget \textit{a priori}? --- Ja. Philosophien har et reent apriorisk Element, det blot ontologiske. Her har den af Hegel med
størst Energie gjennemførte, skjøndt allerede før hans Tid oftere
brugte, dialectiske Begrebsudvikling sin fulde Betydning. Tænkningens almindelige Bestemmelser lade sig virkelig udvikle i en
saadan dialectisk Sammenhæng, at den ene Kategorie fører til den
anden og et apriorisk System af Bestemmelser for al Tilværelse
lader sig ved denne Fremgangsmaade tilveiebringe. Om endogsaa
ingen hidtil given Fremstilling af denne Videnskab i alle Henseender svarer til Ideen af en apriorisk Ontologie, er der præsteret
saa Meget, at det maa være klart for enhver Sagkyndig, at en
saadan Plan lader sig realisere; thi derom kan man være temmelig overbeviist, naar man alt har seet, at de vigtigste aprioriske Begreber staae i et saadant Forhold til hinanden, at det ene gaaer
over i det andet; og Dette i det Mindste er paa Videnskabens
nærværende Standpunct viist med fuldkommen Klarhed. I den
saakaldte speculative Logik, hvis Idee Hegel har fattet, haves en
Gjenoplivelse af den ved den kritiske Philosophie fortrængte Metaphysik eller Ontologie; men med denne Videnskab er ogsaa Philosophiens reent aprioriske Partie tilendebragt. Dermed er et Indbegreb af almeengyldige og nødvendige Bestemmelse for alt Existerende udtalt, men ikke Nødvendigheden af den virkelig existerende
Verden med sin uendelige Bestemthed efterviist. Den virkelig existerende Verden kan man kun kjende ved Erfaring, og ingen Philosophie kan apriori bevise Nødvendigheden af den uudtømmelige
Rigdom af Bestemmelser, hvormed den træder frem. Men Ontologien indeholder ligesom Mathematiken en Sum af hypothetiske
Sætninger: den giver en apriorisk Udvikling af alle de Prædicater
der maa udsiges om Alt, hvad der skal kunne existere, men at
% --- p. 27 ---
\opage{27}\setcounter{footnote}{0}%
Noget virkelig existerer, maa erkjendes paa en anden Vei. Det
er un Existentsens ufravigelige Betingelser, ei dens hele factiske
Indhold, der af Videnskaben kan udvikles \textit{a priori}. Men naar
man altsaa ikke af blot abstracte Begreber kan godtgjøre Nødvendigheden af at Noget existerer, kan man naturligviis endnu mindre godtgjøre Nødvendigheden af, at „disse“ endelige Fornuftvæsner existere med Bestemmelsen Udødelighed.

Men maaskee man med en liden Modificaton af Fordringen
til et ontologisk Beviis for Udødeligheden, kunde fuldkommen tilfredsstille den. Ontologien lærer mig, hvad der nødvendig følger
af, at Noget sættes som tilværende; thi den udfolder mig egentlig blot Begrebet Tilværelse. Men maaskee jeg ved at forfølge
Begrebet Tilværelse befinder, at det ei er færdigt før jeg kommer
til Begrebet endelige Fornuftvæsner, som ere udødelige; naar jeg
da til denne aprioriske Erkjendelse lægger min ved Erfaring givne
Vished om, at Noget existerer, har jeg jo et, vel ei reent apriorisk, men dog aldeles tilfredsstillende Beviis for, at der maa gives
Mennesker eller endelige Fornuftvæsner, som ere udødelige. Mit
Beviis beroer altsaa deels paa den mig ved Erfaringen givne
Vished om, at Noget existerer, deels paa en Tankeudvikling, som
har uimodsigelig, logisk Vished. Uagtet Overbeviisningen om
Udødelighedens Realitet ved denne Vending af Beviset ikke kommer
til at beroe paa en reen apriorisk Tankegang, hvilket er en iøinefaldende Umulighed, er den dog fuldkommen fast og uimodsigelig.

Saaledes have vi med Føielighed indladt os paa en Tankegang, som er den eneste, der kunde tænkes at føre til et ontologisk Beviis for Menneskets Udødelighed. Men en opmærksom
Læser vil let have opdaget, at der er et stort Spring i Beviset.
Der er nemlig forudsat, at endelige Fornuftvæsner og Mennesker
ere Begreber, der uden Videre kunne sættes i hinandens Sted.
Hvis da Beviset skulde være aldeles tilfredsstillende, skulde jeg, ad
den ontologiske Vei, kunne komme til et saa fuldstændig bestemt
% --- p. 28 ---
\opage{28}\setcounter{footnote}{0}%
Begreb om det endelige Fornuftvæsen, at jeg med apodictisk Vished kunde ombytte det med Begrebet Menneske; men hvo vil troe
paa Muligheden af en saa lang apriorisk Tankegang? Indtil da
Nogen løser denne Opgave, som endnu Ingen har forsøgt paa,
nægte vi dristig Muligheden af at føre noget apriorisk Beviis for
Menneskets Udødelighed.

\begin{center}\rule{3em}{0.4pt}\end{center}

\begin{center}IV.\end{center}

I det Foregaaende ere Grændserne for den aprioriske Erkjendelse antydede, og det er kortelig viist, at hverken et mathematisk
eller et ontologisk Beviis lader sig føre for Menneskees Udødelighed. Det er med Hensyn til Videnskabens meest gjængse Talebrug, jeg har omtalt den Vished, som Mathematiken og den,
som Ontologien medfører, som to forskjellige Arter af Vished,
skjøndt jeg selv har den Overbeviisning, som jeg i det Mindste
ikke i den nærværende Sammenhæng seer mig i Stand til at
retfærdiggjøre, at man i begge Videnskaber bevæger sig i een og
samme aprioriske Region. Naar begge Videnskaber træde frem i
deres fuldendte Skikkelse, da maa det vistnok vise sig klart, at
Mathematiken er et i Detail gjennemført Capitel af Ontologien,
den er nemlig en eensidig Udvikling af Quantitetens Bestemmelser.
Men hvad ovenfor er sagt om aprioriske Beviser for Menneskets
Udødelighed, staaer fast, hvorledes man end dømmer om den sidste-Paastand.

Nu spørge vi, om der ikke kan gives noget andet videnskabeligt Beviis for en Bestemmelse ved individuelle Væsners Begreb, end et reent apriorisk. Vi træde her ind paa en anden
Erkjendelsens Mark, hvilken Ontologien gjennemtrænger som et
ordnende Princip, og hvori dens abstracte Begrebsbestemmelser
gjentage sig under mange concrete Skikkelser. Men her bliver
Erkjendelsens Fremstilling af meget meer organisk Natur, end i
% --- p. 29 ---
\opage{29}\setcounter{footnote}{0}%
dens aprioriske Territorium. Her bliver det meer øiensynligt, at
der gives flere mulige Anordninger af Begrebsfølgen, af hvilke
hver har sine Fordele og sine Mangler. Hovedsagen er her en
concret Verdensanskuelse, for hvis Gyldighed der ikke kan føres
andet Beviis end dens fuldstændige Fremstilling. Denne Verdensanskuelse maa som en Spire i Forveien være tilstede i Tænkerens Aand, om han skal være i Stand til at fremstille et
videnskabeligt System, og Systemets formelle Fuldkommenhed
tjener blot til at give det større Tydelighed for Philosophen selv
og at skaffe Meddelelsen større Bestemthed. Er den Verdensanskuelse, der ligger til Grund for den hele Fremstilling oprindelig skjæv eller usand, da hjælper al Energie, Skarpsind og Flid
i Udførelsen af de enkelte Dele ikke det Mindste til at skaffe
Systemet større Gyldighed. Den fuldkommen systematiske Udarbeidelse bringer blot Ophavsmanden til ret at forstaae sig selv og
tvinger Læseren til at tænke sig Systemet netop saaledes, som
Forfatteren har tænkt sig det; i det den meer og meer afskjærer
ham enhver Mulighed af at forfalske Systemets Verdensanskuelse
ved Misforstaaelser og Tilsætninger af en anden Aand.

Enhver fuldstændig Verdensanskuelse har to Sider: den har
deels sin Rod i Erfaringens Verden, deels taber den sig i det
Usandselige. Men ved denne i begge sine Yderligheder ret iøinefaldende Modsætning frembringes den store Misforstaaelse, at enhver af disse Regioner udgjøre et aldeles for sig værende Heelt.
Der gives Tænkere, som holde for, at selv den udvortes Erfaringsverden kan blive Gjenstand for en afsluttet Erkjendelse, der
ved sin afgjorte Paalidelighed har et Fortrin for Speculationen, der
efter deres Mening gaaer ud paa en Drømmeverden. Denne
Mening kommer af Mangel paa Eftertanke; thi en grundig Reflexion vilde vise den grovere Empiriker, at han idelig benytter sig
af Begreber vg Grundsætninger, som ikke kunne forklares af sandselige Fornemmelser, og som dog strække sig gjennem alle hans
% --- p. 30 ---
\opage{30}\setcounter{footnote}{0}%
Antagelser. Nu gives der vistnok Videnskabsmænd af en saadan
intellectuel Uskyldighed, (der snart kommer af Genie, snart af
Dumhed,) at de kunne arbeide med Held i en eller anden Videnskab, uden at forstyrres ved Reflexioner over Videnskabens usandselige Forudsætninger. Men denne Uskyldighed er blandt virksomme Hoveder meget sjælden: de Fleste føle her, at de maae
vogte sig for Grublen over deres egen Fremgangsmaade. Men
begynde de ret at forvikle sig i Reflexioner over deres Tænknings
aprioriske Forudsætninger, da gives der ingen anden Udvei for dem
af den tiltagende Forvirring, end at forfølge de usandselige
Traade --- de aprioriske Begreber --- som gjennemkrydse deres
Forestillingskreds, lige til den oversandselige Region.

Den anden Side af en fuldstændig Verdensanskuelse er Indbegrebet af dens Antagelser om det Oversandselige. Her er det
især, man i den nyere Tid har spurgt om saadanne Antagelsers
videnskabelige Grund. Her er det, man under mange forskjellige
Benævnelser af Problemet har spurgt om Muligheden af et veiledende Princip og en Methode, hvorved man kunde forvisse sig
om en transcendent eller transcendental Erkjendelses Realitet. Hvorledes Individet nuomstunder er kommet til sine bestemte Antagelser om det Oversandselige, kan ei være noget tvivlsomt Spørgsmaal; thi som aldeles eenligt veed det Intet om det Oversandselige. Voxer det op mellem Dyr, har jo endogsaa Erfaring lært,
at det ikke kommer til nogen Bevidsthed om Fornuftens Sandheder. Men i Staternes Skjød bliver det deelagtig i en Tradition om det Oversandselige, som endogsaa før Christendommen
har levet i ethvert udviklet Statssamfund. Blot ved sig selv,
som eenligt Individ, hæver Ingen sig til Bevidsthed om Religionens Begreber: den der troer, at han forkaster al Autoritet og
ved sin Tænknings frie Selvvirksomhed gjenvinder et nyt Resultat, der er hans blot selverhvervede Eiendom, har dog, forsaavidt
han paa en fremmende Maade griber ind i Slægtens Udvikling,
% --- p. 31 ---
\opage{31}\setcounter{footnote}{0}%
altid havt en stor Veiledning i Traditionen, uden hvilken hans
Tænkning havde forvildet sig i subjective, frugtesløse Indfald. De
fleste saakaldte Rationalister \textit{sensu strictissimo}, hvis Tænkning
blot er kritisk, ikke productiv, gjøre ikke Andet end reducere Traditionens Sætninger til en mindre Sum, en Rest, hvorved de
vilkaarlig blive staaende. Selv den absolute Fornægter, hvis aandelige Selvvirksomhed blot viser sig i Negationen af det Bestaaende, vilde jo uden Traditionens Modstand aldeles ikke finde
nogen Anledning til at gjøre sit rene Jeg gjældende. Hans
Tænknings Arbeide var endt af Mangel paa Stof, og han maatte
høilig savne „de Fordomme“, hvis Bestridelse dog skaffede hans
aandelige Liv et Slags Fylde og Bevægelse.

Den positive Selvtænkning, der videnskabelig udtaler, hvad
Tiden veed om det Oversandselige, veiledes, med eller uden Bevidsthed derom, af den i Samfundet levende og voxende Tradition,
der idelig fornyer sin Skikkelse. Selv naar han med den frieste
Selvtænkning fremstiller Erkjendelsen af det Oversandselige i en
selvskabt Form og med fuld Overbeviisning, da har dog den uafbrudt strømmende Overlevering tjent til at give hans Tænkning
sin Retning. En Grunderkjendelse, som han har tilfælleds med
sin aandigt levende Samtid, og som ikke blot møder ham under
Skikkelse af et Læreforedrag, men ligesaa vel under den tit mere
vækkende Form af skriftlige eller mundtlige Aphorismer, ja selv i
dristige Yttringer af hans egen Mund, som han ikke begriber,
hvorfra han faaer, giver ham de faste Puncter, hvorefter hans
Tænkning skrider frem: de tjene ham ligesom Landmaalerens brogede Stokke, der for denne betegne sammenhængende Linier.

Hvilken Kilde, der tilsidst maa tænkes for den religiøse Tradition, som er i en uafbrudt Vorden og Udvikling, og altsaa ikke
maa forestilles som en blot historisk Overlevering, skal egentlig
ikke omtales her, hvor vi tage denne Kjendsgjerning som given,
for dertil at knytte andre Betragtninger om den videnskabelige
% --- p. 32 ---
\opage{32}\setcounter{footnote}{0}%
Erkjendelses Natur. Dog gjøre vi i Forbigaaende den Bemærkning, at det kun er en foreløbig Forklaring, naar man udgiver
en saadan Traditions bestandige Vorden for en Frugt af en ubevidst Dialectik eller af Begrebets egen Selvbevægelse; om man
endogsaa tænker Noget derved, hvilket der just ikke altid er tilstrækkelig Grund til at antage om dem, der føre dette eller lignende Udtryk i Munden. Man kan ligesaa godt forklare sig de
tit ubevidste Fremskridt i Erkjendelsen, hvorved Tidens ideale Tradition beriges, af en høiere Tiltrækning, hvorved Tænkningen ledes
mod sin usynlige Pol; og derved er maaskee allerede Mere sagt.
Udrevet af al videnskabelig Sammenhæng har et Udtryk med videnskabelig Klang slet ikke større Grundighed, end et ret naturligt
Billede. Sammenligne vi Tænkningens uafbrudte Nærmelse til
et Maal, den ikke selv kjender, med Trækfuglens Flugt til et
Land, den aldrig har seet; da er det intet blændende Spil med
tilfældige Ligheder, man ved en saadan Tale driver, men man berører en gjennemgribende Naturanalogie, der har bogstavelig
Sandhed.

Vi have altsaa angivet to væsentlige Sider, der fremtræde
ved enhver nogenlunde fuldstændig Verdensanskuelse: paa den ene
Side den empiriske Opfattelse af Sandseverdnens Phænomener,
paa den anden Side den christelige Tradition om det Oversandselige, der altid beriges og berigtiges i Tidens Løb. Det Tredie
til disse to Extremer, der for os i den nærværende Tingenes Orden har den fuldkomneste Virkelighed og giver Menneskelivet sin
egentlige Fylde, er det Oversandseliges Tilstædeværelse i det Sandselige, naar det bliver Gjenstand for en Erfaring af høiere Art.
Men bestandig huske vi paa, at den sande Gjenstand for vor
Erkjendelse, derfor ikke er indskrænket til disse Aabenbaringsformer, men har en langt rigere Virkelighed, end den vi omfatte
med vore abstracte Begreber og Forestillinger.

Den christelige Tradition og den egentlige Empirie, saavel%
% --- p. 33 ---
\opage{33}\setcounter{footnote}{0}%
som den høiere Erfaring, hvori det Oversandselige møder os paa bestemt
Tid og Sted i Virkeligheden, afgive de discrete Puncter, som
maae have deres Plads i en rigtig Verdensanskuelse, og Philosophens systematiske Fremstilling udtaler med formel Fuldkommenhed kun den Erkjendelse, der først var tilstede hos ham
paa en umiddelbar Maade, og i en udannet Form. En Philosophie har tilsidst intet andet Beviis end Evidentsen af dens fuldstændig udtalte Verdensanskuelse. Hvis da den Verdensanskuelse,
som ligger til Grund derfor, enten strider mod den langt mægtigere Tradition, som er hele Menneskehedens Eiendom, eller mod
de ved lavere og høiere Erfaring fastsatte Puncter, der nødvendig
maae have deres Plads deri, da forsvinder den, saasnart som den
har fuldstændig udtalt sit Ord, og er tilfulde begreben. Det
hjælper da ikke, at Philosophen forsikkrer, at han har ført Beviis
for enhver af sine Sætninger; thi enhver Kjender af Systemet
mærker dog med Tiden, at den hele Cyclus af Sætninger tilsidst
bæres af en Verdensanskuelse, paa hvis Sandhed alle de enkelte
Deles Sandhed beroer.

Hidtil har ingen ført et. videnskabeligt Beviis for Ugyldigheden af den christelige Udødelighedslære; thi Opfinderen af et
saadant Beviis maatte nødvendigviis bruge een af to Fremgangsmaader: enten maatte han udhæve Kjendsgjerninger i Erfaringsverdnen, som ikke lode sig bringe i Harmonie med Læren om Udødelighed, hvilket Ingen hidtil har gjort; eller han maatte give en
overbevisende Fremstilling af en Verdensanskuelse, hvori Udødelighedslæren var aldeles tilovers. Det Sidste har været jævnlig
forsøgt, f. Ex. i den hegelske Philosophie, men Tiden vil vise og
har allerede saa temmelig viist, at den Verdensanskuelse, som fremstilles deri, ligesaa vel som enhver anden, der strider mod den
christelige Traditions Hovedsætninger, vil forsvinde som en Chimaire, naar den er bleven nogenlunde forstaaet af sin Samtid.
Saaledes har Fredrich Richter mod sin Villie bidraget meget til
% --- p. 34 ---
\opage{34}\setcounter{footnote}{0}%
Gjendrivelsen deraf, ved at give fattelige Forklaringer over flere
Puncter deri, der klart vise Udødelighedslærens Uovereensstemmelse
med Systemets Aand. Paa denne Maade fremkommer det indirecte Beviis for Menneskets Udødelighed, der jævnlig gjentager sig
i Videnskabens Historie, og som bedst føres af Udødelighedslærens
Modstandere. Det bestaaer i en consequent Udvikling af en Verdensanskuelse, der aldeles udelukker hiin Lære; thi det vil altid
befindes, at en saadan Verdensanskuelse kun behøver med Udførlighed at udbrede sig i alle sine Følgesætninger, for snart at finde
sin Undergang. Men da der kan gives utallige falske Verdensanskuelser, der udelukke Udødelighedslæren, ligger det i Sagens
Natur, at dette vidtløftige Beviis aldrig bliver fuldbragt.

Paa den anden Side lader der sig heller ikke føre noget
andet Beviis for Udødelighedsbegrebets Realitet, end den videnskabelige Fremstilling af en Verdensanskuelse, hvori det har Krav
paa sin Plads som et Led, der ikke kan undværes i Systemet.
Naar da den Sætning, at Mennesket er udødeligt, nærmest staaer
i Forbindelse med bestemte Sætninger i Systemet, kan derved frembringes det Skin, at der i visse Capitler deraf forefindes Beviser for Udødeligheden, som kan uddrages og for
sig tagne tilveiebringe fyldestgjørende Overbeviisning. Men, nærmere beseet, henter dog Beviset al sin Styrke fra den Verdensanskuelse, som ved hele Systemet faaer sin Fremstilling, og kan
ikke have nogen Betydning for Modtageren, med mindre man i
det Væsentlige hos ham kan forudsætte samme Verdensanskuelse.

Som en almindelig udbredt Verdensanskuelse, med hvilken
den rette Udødelighedslære ikke kan bestaae, have vi alt ovenfor
omtalt den eensidige Pantheisme, som bliver staaende ved det abstracte Begreb om een eneste Substants, der lader sig tilsyne i
en uendelig Række af forgjængelige Væsner. Ligesaa lidet harmonerer Udødelighedstroen med den logiske Pantheisme, som endeel af Hegels Efterfølgere bekjende sig til. Dette paastaae vi
% --- p. 35 ---
\opage{35}\setcounter{footnote}{0}%
med største Sikkerhed, uagtet nogle af Lærens Tilhængere indbilde
sig efterhaanden at kunne bringe dens kunstige Begrebsbygning i
Overeensstemmelse med Christendommens Hovedsætninger. Umuligheden heraf indlyser efter min Overbeviisning endog for en
løselig Betragtning. Efter den logiske Pantheisme kan det videnskabelige Begreb saaledes gjennemtrænge den hele Tilværelse, at
Viden og Væren falde aldeles sammen, uden at nogen ubegreben
Rest bliver tilbage. Der antages, at det er muligt ved en immanent eller, som vi her have foretrukket, at kalde det, en apriorisk Tænkning at udtemme alle Tilværelsens Bestemmelser, saa at
slet intet Residuum levnes som Gjenstand for blot Erfaring.
Deraf følger med den rimeligste Consequents, som en heel Deel
af Systemets ærlige Bekjendere uden Betænkning indrømmer, at
de talløse Fornuftvæsner ikke som „disse Personer“ udgjøre nogen
væsentlig Deel af Tilværelsen, men blot ere temporære Former
for Aandens med logisk Nødvendighed bestemmelige Proces. Ethvert endeligt Fornuftvæsen bliver efter den til Grund liggende
Hypothes kun et forsvindende Bølgeslag i Tænkningens Ocean,
hvis Undulation er bestemt ved en ufravigelig Nødvendighed.

Under Tidsalderens Bestræbelse for at opnaae den formelle
Forstaaelse af denne Philosophie, kom det ikke til den sidste Prøve,
som skulde vise, om nu omsider den Verdensanskuelse, som var
bleven udtalt deri med en ikke tilforn kjendt alsidig Bestemthed,
var i Stand til at frembringe fuld Overbeviisning. Nu, da
Systemet har udtalt sit Ord, kan man med Sikkerhed sige, at
det ikke har udholdt denne Prøve. Det strider nemlig mod væsentlige Puncter i den ældgamle Tradition, hvis Indhold er Menneskehedens sande Erkjendelse af det Oversandselige: thi ethvert
System, der udelukker væsentlige Partier af denne Erkjendelse,
kan aldrig i Længden hævde sine Fordringer paa almindelig Gyldighed i Videnskaben. Ikke mene vi, at det kan finde sin sande
Gjendrivelse ved en Henviisning til Traditionen under en af dens
% --- p. 36 ---
\opage{36}\setcounter{footnote}{0}%
umiddelbare Former, men det bliver nødvendigviis fortrængt ved
Sandhedens Magt i Videnskabens Rige. Det er ikke som udvortes Autoritet, at den religiøse Tradition betvinger et mangelfuldt System, men ved Tænkefrihedens Brug og Fornuftens
Kraft. Det er nemlig alene denne levende Tradition, der formaaer at fylde Sindet med fornuftig Overbeviisning.

Egentlig findes netop den her udtalte Synsmaade med andre
Ord hos Hegel selv, men hans Philosophie har ikke den foregivne, men kun en tilsyneladende Overeensstemmelse med den christelige Traditions Indhold. Et afgjørende Beviis herpaa er dens
fuldkomne Savn af en sand Udødelighedslære. Dette er nu almindelig følt af hiin Skoles selvstændige Tilhængere, og de søge
at raade Bod paa dette Savn, med hvilket de klart indsee, at
Systemet ikke kan opnaae det udbredte Herredømme, som de
havde haabet for det. Men, hvis en saadan ny Bestanddeel skal
% PRINTED AS IS: V142 q: crop, layer, tesseract -- a turned n, as printed
optages i Hegels System, følger deraf en gjenuemgribende Forvandling af alle Systemets Hovedpartier, efter hvilken dets Tilhængere ikke meer kunne beholde Hegels Navn til deres Løsen,
og ikke meer have noget Tilhold i dette Navns udvortes Autoritet.

Naar nu Hegels Disciple, hvergang et Fragment af et nyt
System kommer dem for Øie, møde det med den Erklæring, at
det ikke er beviist, da have de for saavidt Ret, som ingen Verdensanskuelse har været udtalt med formel Fuldstændighed, uden
deres store Mesters, og forsaavidt som der ikke kan føres noget
andet Beviis for en Philosophies Gyldighed, end dens bestemte og
fuldførte Fremstilling. Men der vækkes hos Ukyndige, der saa
let herved tænke paa den mathematiske Demonstration, en Forestilling om apodictisk Vished, som ei er tilveiebragt ved Hegels System, og som udenfor Ontologien ei kan tilveiebringes ved nogen
Philosophie. --- Naar da en Verdensanskuelse er bleven fuldstændig fremstillet, og saaledes det Beviis. for dens Gyldighed er ført,
som er det eneste mulige, da kan dette netop befindes at være et
% --- p. 37 ---
\opage{37}\setcounter{footnote}{0}%
Beviis for dens Ugyldighed, fordi Overbeviisningen derom muligen kan være vakt hos Enhver, hos hvem den sande Verdensanskuelse findes i en umiddelbar Skikkelse; thi en Saadan kommer let efter, at en urigtig Lærebygning, uagtet al sin architektoniske Fuldkommenhed, tilsidst dog kun er en omstændelig Paastand. I Sandhed gjendrevet kan da et System være ved sig
selv eller ved den Verdensanskuelse, som fremstilles deri; men
unægtelig savnes den udvortes Gjendrivelse, saa længe der ei i
dets Sted er sat en fuldstændig og bestemt Udvikling af den
sande Verdensanskuelse. --- Dog et andet Spørgsmaal, som vi
troe at maatte besvare benægtende, bliver det, om en formel fuldkommen Fremstilling af den sande Verdensanskuelse bør eller tør
fremtræde med en saa kunstig Articulation, og en saa stor Overflødighed af synonymiske Distinctioner, som det hegelske System har
gjort. Det vil sikkert allerede derved gjøres umuligt, at Tidsalderen i sin stigende Utaalmodighed, ei anden Gang vil følge en
saa gigantisk Lærebygning med sin Opmærksomhed. Men hvad
man end vil vente, der skal udvikle sig af Tidens Skjød, er det
især af Vigtighed at fastholde den Sandhed, at en Philosophie
ved sin blotte formelle Fuldkommenhed ikke vinder nogen af sin
Verdensanskuelse uafhængig Vished. Er dens Verdensanskuelse
skjæv, da afløses den foreløbig ved en mindre videnskabelig, umiddelbar Erkjendelse af Sandheden, der finder sin urolige, gjærende
Form i fragmentariske Fremstillinger.

Hvad der gjælder om den logiske Pantheisme, der med usædvanlig Kraft og Kunst er bleven udtalt i et fuldstændigt System, gjælder endnu meer om den vulgære Pantheisme, der har
havt sin Form i mangfoldige flygtig antydede Systemer og nu
tilsidst har den i det literære Conversationssprog. Mod den har
den sande eller alsidige Verdensanskuelse end større Ret til at fremtræde
i sin umiddelbare vaklende Skikkelse. Denne Pantheisme er nu bleven til en herskende Fordom. Hvad Betydning den har i Viden%
% --- p. 38 ---
\opage{38}\setcounter{footnote}{0}%
skabens Historie, er ovenfor omtalt, men den udvortes eller psychologiske Grund til dens nærværende Herredømme i den dannede
Verden er den vidt udbredte Frygt for at lægge borneret Tænkemaade for Dagen. Fordi det virkelig er Mangel paa Evne til
at holde Skridt med Tiden, der nuomstunder gjør, at Somme
reent savne det pantheistiske Element i deres Bevidsthed, frygter
man for at bekjende sig til Troen paa en personlig Gud. Man
holder nemlig for, at denne Tro ikke lader sig bringe i Harmonie med nogen Slags Pantheisme, og derfor kan paadrage sine
Tilhængere Ringeagt\footnote{Der gives intet literært Strategem, der i Tydsklands videnskabelige Polemik har gjort saa god Virkning og vakt saa megen Efterligning, som haanlig Omtale af Anderledestænkende. Ligefra
Fichte til Hegel har de berømte Systemers Efterfølgere saa nøiagtig copieret deres Mesteres Rhetorik i den haanlige Declamation, at man ret kan forestille sig, hvorledes deres unge Hjerter i
deres Læreaar have banket af Længsel efter at komme til at føre
et saadant overlegent Sprog. Saaledes overbyder Schad langt
Fichtes djærve Udfald paa Idealismens Modstandere og blandt
mange Andre har Hinrichs i Fortalen til sin Logik uden at
nævne sin Kilde udskrevet Sarcasmer af Hegel.} i den videnskabelige Verden.

Ligesom altsaa Overbeviisningen om Udødelighedsbegrebets
Realitet altid hviler i en Verdensanskuelse, kan der heller ikke
gives andet Beviis derfor, end den organisk sammenhængende
Fremstilling af en saadan Verdensanskuelse. Uagtet derfor Udødelighedslæren i Systemet kan fortrinsviis knyttes til visse Hovedstykker deraf, finder dog den opmærksomme Læser i alle dettes
Partier Spor til Forudsætning af Dødelighed eller Udødelighed.
Hvis Udødelighedsbegrebets Ugyldighed ei var forudsat i saa
mangfoldige Bearbeidelser af Philosophiens Discipliner, hvori det
% PRINTED AS IS: V155 q: crop, layer, tesseract -- a MISPRINT, as printed
ei med et Ord berøres, vilde de overalt have faaet et aldeles forandret Udseende. Med andre Ord en videnskabellg Fremstilling af
Philosophiens Læresætninger bliver kun for Læseren synthetisk, for
% --- p. 39 ---
\opage{39}\setcounter{footnote}{0}%
Philosophen selv er den altid analytisk, og ethvert Hovedpunct er
forudbestemt i den Verdensanskuelse, som successivt oprulles under
Fremstillingen og som i sin Totalitet er sit eget Forsvar.

\begin{center}\rule{3em}{0.4pt}\end{center}

\begin{center}V.\end{center}

Af Adskillige hører man den Erklæring, at de sætte liden
Priis paa videnskabelige Beviser for Menneskets Udødelighed, fordi
deres Overbeviisning derom beroer paa Udsagn af den hellige
Skrift, om hvis Guddommelighed de have fuldkommen Forvisning.
Men, dersom deres Forsikring er sand, dersom de føle virkelig
Overbeviisning om den christelige Læres Sandhed, kan denne
Overbeviisning dog umulig være fuldkommen uafhængig af Lærens Indhold. De maa dog først paa en eller anden Maade
være komne til Vished om Sandheden af Christendommens Hovedsætninger, inden de kunne bringes til at betragte Skriftens
Ord som deres Erkjendelses Grund. Da nu Udødelighedslæren
selv netop er en af Bibelens Grundlærdomme, hvis Gyldighed
man først maa antage, for at komme til fornuftig Tro paa Bibelen, kan man med bedre Grund opstille den omvendte Paastand, at fornuftig Overbeviisning om Skriftens Ufeilbarhed blandt
Andet beroer paa Overbeviisning om Menneskets Udødelighed.
Nuomstunder kan dog intet Menneske, der staaer paa et Dannelsestrin, der medfører Modtagelighed for fornuftig Overbeviisning, være saa ubekjendt med Skriftens Hovedlærdomme, at det
ikke veed, at Udødelighedslæren findes mellem dem. Hvis han da
paa en fornuftig Vei kommer til at troe paa Skriftens Autoritet,
maa denne Lære dog komme med i Betragtning, medens hans
Tro paa Bibelens Tilforladelighed uddanner sig. Saaledes forholder det sig ogsaa altid i Virkeligheden: al fornuftig Overbeviisning om Bibelens Guddommelighed, har sin Grund i en fore%
% --- p. 40 ---
\opage{40}\setcounter{footnote}{0}%
gaaende Forvisning om den christelige Verdensanskuelses Sandhed.
Nu er det jo endogsaa en stor Sjældenhed, at nogen Dannet
kommer til Bibelens Læsning, inden han er bleven bekjendt med
den Cyclus af religiøse Begreber, der udgjøre det væsentlige Indhold af Tidens Bevidsthed om det Oversandselige. Men hvilken
Combination af usædvanlige Omstændigheder man end her vil
digte sig, vilde aldrig noget Individ blot af Skriften være i
Stand til at udvikle sig den hele christelige Tradition. Et Erfaringsbeviis var her umuligt; thi den aldeles Eenlige kan aldrig
komme til fornuftig Selvbevidsthed, og naar han ved Samliv med
Andre vækkes dertil, da faaer han med det Samme Bekjendtskab
med Selskabets fælleds Forestillingskreds, der er saa gjennemtrængt
af det christelige Element, at dette ikke tit et Experiment lader sig
udsondre.

Paa dette Sted vil jeg forklare mig noget nærmere over
Udtrykket den christelige Tradition, som jeg paa disse Blade
saa tit har betjent mig af. --- Hos Enhver, der i christelige
Stater kommer til aandelig Selvbevidsthed, begynder et Anlæg
til en selvstændig Verdensanskuelse at udvikle sig, ved udenfra
kommende Anledninger. Den subjective Selvudvikling er paa ethvert Punct betinget ved et Stof, der med eller uden Plan meddeles udenfra. Individets religiøse Bevidsthedskreds, --- den være
nu meest udfyldt med Billeder eller med Begreber, --- drager sin
væsentlige Næring af den christelige Tradition, som under utallige
Metamorphoser bevæger sig i de cultiverede Stater. Det er den
samme Tradition, der snart viser sig under Form af populære
Forestillinger, snart reproduceres i philosophiske Begreber. Men
hvad der skal have væsentligt Hold, maa følge Strømmen af den
christelige Tradition. Hvad der vil danne sig aldeles udenfor
dette fælleds Medium, bygger i Luften uden Fundament, og opløser sig af sig selv som en Sum af subjective Indfald. Ingen sindriig Combination af subjective Indfald kan i længere
% --- p. 41 ---
\opage{41}\setcounter{footnote}{0}%
Tid friste sin falske Tilværelse uden tilsyneladende Overeensstemmelse med den christelige Tradition. Den høieste theoretiske Aandsfrihed kan kun opnaaes ved en selvstændig Reproduction af den
meddeelte fælleds Tradition; og deri bestaaer tildeels et Individs
Sundhed paa det aandelige Liv, at det er i Stand til af den
store Vrimmel af falske Forestillinger og Begreber, at uddrage
den rette Tradition til Næring for det aandelige Liv.

Men under hvilken Form Individet end besidder sin Erkjendelse, afledes al Vished for enkelte Sætninger fra deres Sammenhæng med den Verdensanskuelse, der bærer det Hele, og som i sin
Totalitet medfører Overbeviisning. Men den Verdensanskuelse,
der medfører sand Vished, er altid den fremskridende christelige
Aands Værk, og ingen blot Opdagelse af den individuelle Selvtænkning. Dog kan den saa aldeles træde frem under Selvtænkningens Form, at Individet kan føle fuldkommen fri Selvvirksomhed under den Tænkning, hvorved det reproducerer den fælleds
christelige Tradition. Det mærker maaskee ikke engang Traditionens Magt over det, men indbilder sig blot ved egen Kraft at
have opdaget den fælleds Sandhed, som under en ny Skikkelse
fremtræder hos det. --- Neppe vil Nogen, der veed hvad christelig Aand er, miskjende dens Indflydelse paa en heel Deel af
Schillers Arbeider, uagtet denne Digter ikke selv henførte dem
til dette Udspring, eller fremtraadte som en Bekjender af Christendommen. Den christelige Tradition kan fremtræde som rationalistisk, for saavidt som Individet er sig bevidst, at dets Verdensanskuelse, som stemmer dermed, udspringer af dets eget Indre, i
det den med fuldkommen Frihed reproduceres i et videnskabeligt
System. Men den er desuagtet paa en Maade supranaturalistisk,
fordi det kun er i sin Overeensstemmelse med deu objective Tradition at den frie Tænkning kan komme til Vished og Tilfredsstillelse. Naturligviis har den, hos hvem den christelige Tradition
kraftig udvikler sig, dog ikke noget rigtigt Begreb om dens Na%
% --- p. 42 ---
\opage{42}\setcounter{footnote}{0}%
tur, naar han anseer sin Erkjendelse for et Product af blot individuel Selvtænkning; men paa den anden Side giver man ikke
den christelige Aand den Ære, som tilkommer den, hvis man aldeles vil nægte det supranaturalistiske Elements Tilstædeværelse
hos Enhver, som ikke har saa udviklet Bevidsthed om sine Tankers
Genesis, at han selv antager dem for christelige og kan forfølge
dem tilbage til deres historiske Oprindelse. Det er vanskeligt at
indsee, at man mindre erkjender Christendommens Verdensmagt,
naar man anvender Prædikatet christelig paa den Bevidsthed om
det Oversandselige, der deriveres fra den, om den end fragaaer
sit Udspring, end naar man indskrænker denne Benævnelse til
dem, der udtrykkelig bekjende sig til et kort Symbols Formular.
--- Om man end tager i Betænkning at kalde Individerne
Christne, naar de ikke selv vedkjende sig dette Navn, kan man
dog med god Grund kalde deres Livsanskuelse christelig, naar dens
væsentlige Characteer bedst betegnes ved dette Prædicat.

Af det, som alt er sagt, er det nu for det Første klart, at
her ved Ordet Tradition ei tænkes paa en uforandret Overlevering af et bestemt Antal Ord. Den overleverede Sandhed udvikles og voxer hos mangen En, der ikke selv bliver det vaer; thi
Erkjendelsen kan ei være levende uden Fremskridt og Forandring
af Form. Men jeg har paa disse Blade med Flid brugt Ordet
Tradition, for dermed ret stærkt at betegne det ene Moment i al
Sandhedserkjendelse, som tit oversees, idet man udhæver det andet,
nemlig det, at al sand Erkjendelse træder frem som Selvtænkning. Det
eenlige Individs Selvtænkning er afhængig af Menneskeslægtens
fælleds Erkjendelse, saa at det aldrig ved sin blotte subjective
Aandskraft opdager Sandheden som et absolut nyt Indhold, men
den frieste Selvtænkning reproducerer og udvikler hele Slægtens
Arvelod. En saadan Traditions Realitet antager næsten enhver
fornuftig Tænker; men indrømmes en saadan Tradition, følger
deraf, at den maa have sit bestemte Indhold, og der bliver
% --- p. 43 ---
\opage{43}\setcounter{footnote}{0}%
ved enhver frisk Reproduction deraf med god Grund spurgt, om
dens Indhold er gjengivet med Fuldstændighed. Dernæst har jeg
kaldt denne Tradition christelig, et Prædicat som nærmest leder
Tanken til den oversandselige Side af den overleverede Verdensanskuelse. Thi med Hensyn hertil bestrider vel neppe nogen Indsigtsfuld den Kjendsgjerning, at Traditionen er væsentlig forandret i det christelige Tidsrum. Hvis nu Nogen frakjender den
christelige Revolution i Menneskehedens Bevidsthed en saadan gjennemgribende Indflydelse paa Erkjendelsens hele Organisation, at
Prædicatet christelig for den hele Tradition derved kan ansees for
at være retfærdiggjort, kan jeg ikke her føre nogen Strid derom:
jeg indskrænker mig blot til den Bemærkning, at det i ethvert
Tilfælde her nærmest er Bevidstheden om det Oversandselige der
kommer i Betragtning.

Efterat vi have berørt den Mening, at Udødelighedslæren
bedst kan bygges paa Skriftens Udsagn, føie vi nogle Bemærkninger til om en Theorie, som let kan bringes i Forbindelse dermed, nemlig den, at saavel Udødelighedslæren som Enhver af Religionslærens Hovedsætninger beroer paa Følelsens Vidnesbyrd.
Der gives neppe nogen anden Titel, hvorunder saamange Synsmaader samles i Flæng, men desuagtet have de mangfoldige Theorier, der gjøre Følelsen til Erkjendelsens Criterium, for den fælleds Benævnelses Skyld tilsammen været underkastede stærk literær Forfølgelse. Alligevel kommer ingen Overbeviisning om en
væsentlig Sandhed i Stand, med mindre Følelsen kan tilfredsstilles
derved, og det er altid fra Følelsen, at det første Stød til Reformer
i Erkjendelsen udgaaer. Saaledes kan der med god Grund paastaaes, at Følelsen er en Magt, som altid factisk regjerer Videnskabens Historie, enten man vil indrømme den Ret dertil eller
ikke. Endog den strænge Iver mod den eensidige Provocation til
Følelsen har selv sit Udspring fra en Følelse.

% --- p. 44 ---
\opage{44}\setcounter{footnote}{0}%
For at komme til Sagen forudsætte vi som fuldkommen afgjort,
at ethvert menneskeligt Individ har den Bestemmelse at udtrykke
det almindelige Fornuftliv under en eiendommelig Form, eller for
at udtrykke den samme Tanke med et Par andre Talemaader, at
dets virkelige Liv er dets Overeensstemmelse med sit Livs Idee;
at det under ethvert sandt Fremskridt i sin Udvikling retter sig
efter et høiere Forbillede. Naar Mennesket da i et Tidsmoment
befinder sig i fuldkommen Overeensstemmelse med sin Idee, og
har en umiddelbar Bevidsthed om denne Harmonie, er denne Bevidsthed en behagelig Følelse. Denne Sætning gjælder for alle
Menneskelivets Former, de lavere saavel som de høiere: følgelig
maa den Fremskriden i Sandheds Erkjendelse, som er alle Fornuftvæsners Bestemmelse, nødvendig medføre den behagelige Følelse af Individets Harmonie med Ideen.

Men ligesaa vist som det er, at den forskende Tænknings
Vendinger, naar den gaaer ud paa en fuldstændig Verdensanskuelse, altid bestemmes af Følelsen, saa lidet fremmer det Sandheds
Erkjendelse, om ethvert Individ vil indskrænke sig til sin Subjectivitet og gjøre den lydelig gjældende midt i Videnskabens Rige.
Det kommer da an paa, hvis Følelse tilsidst er den reen menneskelige; thi kun den, der kan godtgjøre, at han er den fuldkommen restituerede Adam kan forlange, at Andre skulle lade deres Tænkning styres efter hans eiendommelige Følelses Compas. Indtil
den Tid kan han ikke gjøre Andet end enten indslutte sig i sin
monadiske Selvtilfredshed, hvortil han har en ubestridelig Ret;
eller bestræbe sig for at sætte en Begrebsudvikling i Stedet for
en anden som oprører hans naturlige Følelse.

Vi have ovenfor sagt, at den Tænkning, der gaaer ud paa
Erkjendelse af concret Virkelighed, bestandig gaaer ud fra en Følelse. Vi maae her bestemt afvise den mulige Misforstaaelse af
disse Ord, at dermed skulle menes den upathologiske Følelse af
Overbeviisning, som man undertiden hører omtale, i Følge en an%
% --- p. 45 ---
\opage{45}\setcounter{footnote}{0}%
den Brug af Ordet Følelse. Om der gives nogen i Ordets
bogstavelige Betydning følelsesfri Tænkning lade vi her staae aldeles uafgjort; men det kan ikke nægtes, at der gives en Tænkning og en Overbeviisning, hvori den personlige Interesse træder
umærkeligt frem. En saadan uinteresseret Tænkning er den mathematiske og den ontologiske, men den har ogsaa kun at gjøre
med ahstracte Begreber, og medfører paa ingen Maade nogen
overmenneskelig Sindets Tilstand. Naar Tænkningen gaaer
ud paa en concret Verdensanskuelse, kan den kuns i Momenter
være fuldkommen ligegyldig og uden al personlig Bevægelse ved Resultatet. Det er en vidt udbredt Vildfarelse, at den Tænkning
% PRINTED AS IS: V182 q: crop and tesseract -- a MISPRINT for kunne, as printed
hvorved en reel Verdensanskuelse dannes, skulde knnne foretages
med Menneskets halve Sind, som Regning med Tal eller Betragtning af geometriske Forhold.

Men naar vi nu paa den ene Side forsvare den Sætning,
at Tænkningen i det Hele taget erholder sin Retning af Følelsen,
paa den anden Side indrømmer, at intet Individ er berettiget til
at opstille sin Følelse som den normale, synes det umuligt, at
den menneskelige Tænkning kan stride sikkert frem mod sit Maal,
mod Sandheds Erkjendelse. --- Men fordi man indseer, at det
enkelte Individs Følelse ikke afgiver nogen Garanti for dets Tænknings Rigtighed, kan man ved Siden deraf meget vel have den
Overbeviisning, at hele Slægtens Tænkning ved Følelsen drages
mod sin rette Pol. At antage dette bliver netop det Samme,
som at antage, at Fornuften bestandig er den stærkeste Magt i
Menneskeheden, og at Individernes Afvigelser, (der tildeels gjendrive hinanden), forsvinde for Sandhedens Aand, ved hvilken
hele Slægten føres afsted.

En Livsanskuelse, der oprører Menneskehedens Sandhedsfølelse, har ved at udtale sig selv med Tydelighed og Bestemthed
foreløbig gjendrevet sig selv, og derved vækkes Bestræbelsen for at
fremstille den christelige Tradition i et nyt videnskabeligt System,
% --- p. 46 ---
\opage{46}\setcounter{footnote}{0}%
der afløser det gamle, hvis Utilstrækkelighed er almindelig følt.
Men, at Mennesket, efterat Videnskabens Idee blot er anet,
skulde berolige sig ved en umiddelbar dunkel Bevidsthed om Religionens Hovedindhold, er utænkeligt.  Nuomstunder er det øiensynlig Menneskehedens Opgave, at bringe den christelige Traditions
Indhold til videnskabelig Bevidsthed, og derfor kan et mangelfuldt System kun fuldstændig gjendrives ved et rigtigere, eller et
der med større Alsidighed sammenfatter, hvad der hører til den
universelle Tradition.  Men da for Tiden enhver Religionsphilosophie, der har mindste Betydenhed, udgiver sig for at have den
christelige Tradition til sit Indhold, faaer dens udvortes Autoritet
et farligt Stød, naar dens Uovereensstemmelse med et eller andet
væsentligt Princip i Christendommen har tydelig viist sig.  Nu
gives der, ingen Sandhed, hvis Mangel det er saa vanskeligt at dække
ved Palliativer eller at udfylde ved Sætninger, der have fjern
Lighed dermed, som Udødelighedslæren, fordi Savn deraf gjennemtrænger alle Hovedstykker af Systemet.  Derfor vil Mangelen
af denne Lære altid være den hemmelige Brøst, der tidlig eller sildig bereder det forresten meest skarpsindige System sin Undergang.

\begin{center}\rule{3em}{0.4pt}\end{center}

\begin{center}VI.\end{center}

Ovenfor er den Sætning opstillet, at Udødelighedslæren bestandig vil gjøre sig gjældende som et nødvendigt Moment i Sandhedens Erkjendelse.  Ethvert System, der udtaler en Verdensanskuelse, hvori hiint væsentlige Moment mangler, vil nemlig, naar
det er gjennemført i sine Consequentser, først befindes uovereensstemmende med den sande Tradition i dens umiddelbare Skikkelse
og dernæst blive fortrængt af et fuldstændigere videnskabeligt System.  Vi have tilsyneladende talt om tre forskjellige Magter, der
nødvendig i Tidens Løb maa beseire ethvert vildfarende System,
den sande Tradition, den normale Følelse og den sande Videnskab%
% --- p. 47 ---
\opage{47}\setcounter{footnote}{0}%
men det er nærmere beseet een og samme universelle Overmagt,
som under disse forskjellige Titler gaaer af med Seiren.  Den
christelige Verdensanskuelse har ikke sin fuldstændige Form blot i
% PRINTED AS IS: V188: layer and tesseract -- a MISPRINT, as printed (the crop missed the word)
den hellige Skrift, men i mange forskjellige Aaabenbaringer af den
levende Verdensaand; thi hvad der tilfredsstiller den normale
Følelse er Christendommens væsentlige Indhold.  Men naar man
da fra et Systems Standpunct spørger, med hvad Ret en Verdensanskuelse giver sig Navn af sand, og med hvad Ret en Følelse roser sig af at være den normale, da kan intet andet almeengyldigt Svar gives derpaa end den systematiske Fremstilling deraf.
Over en videnskabelig Form, som er af den Betydenhed, at den
bør ændses, kan hverken ubestemt individuel Følelse eller Troeslæren i sin umiddelbare Skikkelse vinde grundig Seir.  For at beholde Marken maa Sandhedens Aand paatage sig samme Skikkelse som dens Modstandere.

Nærmest skal altsaa her den Sætning udvikles, at ethvert
System, der indeholder Udødelighedens Fornægtelse, gjendriver sig
selv ved sine Følgesætninger.  Det er ikke min Plan i en broget
Vrimmel af Phantasiebilleder, at fremstille de Livsforhold, hvori
Udødelighedshaabets Tab maa være meest oprørende for den naturlige Følelse.  Det har oftere været gjort, f. Ex. af Jean
Paul, og det er den eneste Maade, hvorpaa man kan anskueliggjøre den Sandhed, at Dødeligslæren ikke kan tilfredsstille den
menneskelige Følelse.  Her kan, som sagt, intet Beviis for Udødeligheden ventes udført; her anstilles blot Reflexioner over den Mulighed at føre Beviser derfor.  Vi betragte da nu den Tankegang,
som vi ovenfor have omtalt som et uendeligt indirecte Beviis for
Udødeligheden, nemlig den, at Tilintetgjørelseslæren i enhver Sammenhæng, ret udviklet, nøder den menneskelige Bevidsthed til en
Vending, hvorved den fra Fornægtelsens Tomhed kommer tilbage
til den sande Religions Indhold.

% --- p. 48 ---
\opage{48}\setcounter{footnote}{0}%
Naar Mennesket med Overbeviisning om sin individuelle
Forgjængelighed ret ordner sin Bevidsthedskreds, taber dets Liv enhver væsentlig Betydning.  Den menneskelige Bevidsthed kan da
i ingen Henseende beholde sin høiere Fylde, uden forsaavidt som
den med Selvmodsigelse bevarer et laant Gjenskin af Ideens Lys,
der hører forgangne Perioder til.  Erfaringen har lært, at den
menneskelige Bevidsthed vel kan udvikle sig til herlige Former,
hvor Begrebet om Udødeligheden kun er dunkelt forhaanden; men
under den dogmatiske Fornægtelse deraf taber det, naar Mennesket
har ret befindet sig, al sin Glands og Rigdom.  Denne Paastand ville vi tydeliggjøre ved kortelig at gjennemgaae de væsentlige Momenter af Ideen, som giver Menneskelivet sit Indhold.
Men dette skal efter Planen skee saaledes, at de Puncter meest
udhæves, ved hvilke Savnet af Udødelighedsbegrebet lettest indlyser.

For det Første kommer, efter Tilintetgjørelseslæren, den sande
Selvagtelse og overhovedet den rette Selvkjærlighed til at mangle
sin egentlige Gjenstand.  Her bliver den organiske Sammenhæng
af Ideens forskjellige Sider især i Dine faldende, fordi Selvkiærligheden gjennemtrænger alle Livets Former og det abstracte Selv,
strængt for sig taget, aldeles bliver til Intet.  Vi kunde derfor,
hvis her skulde gives en udførlig Udvikling af dette Moment efter Sagens Natur brtragte alle de følgende Momenter fra dette
særegne Standpunct.  Men vi indskrænke os til følgende Vink.

Naar Selvkjærligheden mangler sin Gjenstand kan Ideen
ikke komme til sit sande Liv i Menneskets Bevidsthed.  Mennesket erkjender nemlig ikke den egentlige Sandhed, naar det ikke
tillige føler den.  Den Erkjendelse, der giver sig selv det pralende
Navn af reen Tænkning og udelukker al Sympathie, den er kun
en abstract med Flid udstykket, Erkjendelse.  Forsaavidt Tænkningens Arbeide, ja Tænkningens Slid er en nødvendig Betingelse
for Opnaaelsen af concret Erkjendelse, er den følelsesfrie Tænkning, som finder Sted i Mathematik og Ontologie, af stor Be%
% --- p. 49 ---
\opage{49}\setcounter{footnote}{0}%
tydning og Vigtighed.  For saavidt fremdeles den fuldendte Erkjendelse i den nærværende Tingenes Orden aldrig kan opnaaes,
hører den abstracte Tænkning her aldrig op.  Men jo mere Begrebets sønderrevne Stykker begynde at samle sig til et Heelt, jo
mere Tænkeren begynder at skimte den rette Sammenhæng i Tilværelsen, desto mere blive Tænkningens Arbeider afbrudte ved Anskuelsens Hviledage.

Al Erkjendelse gaaer ud paa det altid Værende, men under
den sande Erkjendelse veed det erkjendende Subject, at det selv
hører med til det altid Værende.  Naar Tilintetgjørelseslæren ret
gjennemtrænger Menneskets Tankekreds, vil Erkjendelsen af det
Evige meer og meer fordunkles.  Ifølge denne Synsmaade hører
nemlig Individet som Individ ikke til det altid Værende, men
den, som ikke erkjender sig selv for at høre med til det Uforgjængelige i Tilværelsen, har ikke den levende Erkjendelse deraf, om
han end med større logisk Dygtighed kan fremstille sin mangelfulde Verdensanskuelse.  Kun, hvor den individuelle Udødeligheds
Begreb optages med som et Element i Menneskets Verdensbevidsthed, kan sand Selvkjærlighed finde sin Plads deri.  Og
Selvkjærlighed er en ufravigelig Betingelse for den levende Erkjendelse og den sympathetiske Tilknytning til det altid Bestaaende,
uden hvilken Bevidstheden ei kan svare til sin Idee.

De dogmatiske Fornægtere af Udødelighedstroen indrømme
naturligviis den Kjendsgjerning, at Tilintetgjørelseslæren vanskelig
kan blive ret almindelig udbredt i Menneskeheden, og de finde ligesom vi en vigtig Forklaringsgrund til dette Phænomen i den naturlige Selvkjærlighed.  Men de tage Selvkjærligheden i en anden Betydning end den vi her have taget den i.  De sætte den
ei i Forbindelse med den dybere Selvagtelse, i Følge hvilken Mennesket erkjender sig som Individ for et Led af den altid værende
oversandselige Verden; men de tænke ved Ordet Selvkjærlighed
blot paa den physiske Selvopholdelsesdrift, som Mennesket paa sit
% --- p. 50 ---
\opage{50}\setcounter{footnote}{0}%
abstract dyriske Standpunct ei kan overvinde.  Efter den rette
Verdensanskuelse er denne lavere Tilskyndelse kun et Moment i
den sande Selvbevaringsdrift, men et Moment, der er saa uvæsentligt, at denne, som Historien tit har viist, kan føre Individet
til at handle mod hiin elementære Drift.

Det er deels Praleri, deels Selvbedrag, naar saa Mange
paastaae, at en Verdensanskuelse, der tilfredsstiller den reen-menneskelige Selvfølelse kan danne sig uden Udødelighedsbegrebet.  I
denne Henseende lægger F. Richter i sit omtalte Skrift større
Forstand og Sandhedskjærlighed for Dagen, idet han uden Betænkning tilstaaer, at Visheden om Individernes Forgjængelighed
er den største af alle de Ulykker, som ere uadskillelig forbundne
med Menneskehedens Kaar, og at det er en Ulykke, mod hvilken
der ikke gives nogen tilfredsstillende Trøst.  Deri ville vist de
Fleste være enige med ham, men ikke i den Forudsætning, som
ligger til Grund for hans Philosophie, nemlig den, at en fuldkommen Erkjendelse af Sandhed skulde kunne slutte med en uopløselig Dissonants, eller med andre Ord, at Sandheden skulde
kunne være oprørende for den menneskelige Følelse.

Tilintetgjørelseslæren maa da for det Første af den Grund
forsvinde, at den oprører den sande Selvfølelse og Selvagtelse,
uden hvilken Individet ei med den Deeltagelse, som er et Moment i al sand Erkjendelse, kan slutte sig til det Evige.  Dernæst kan ingen sand selskabelig Interesse bestaae, naar Mennesker
antages blot for temporære Væsner, der uophørlig blive til og
for stedse forsvinde fra Tilværelsen.  Al Stræben for Menneskevel og
Indretning af Stater, maa nødvendig ophøre, dersom Menneskene
med Overbeviisning om hinandens absolute Forgjængelighed ret
befinde sig over det Intetsigende i deres Planer.  Kun den var
i Stand til under slige Vilkaar at bevare sin Livslyst, der ved
Travlhed i Livets Leeg kunde bringe det til en uafbrudt Selvforglemmelse.  Dertil ere Somme virkelig i Stand: de ansee Alt
% --- p. 51 ---
\opage{51}\setcounter{footnote}{0}%
% PRINTED AS IS: V207 q: crop, layer, tesseract -- a MISPRINT for at, as printed
for ar være Intet, men udføre med Tilfærdighed den ene Plan
efter den anden, for ved en frivillig Illusion at vænne sig til at
behandle det formeentlige Intet, som om det var Noget. Men
til et saadant tankeløst Selvbedrag kan hele Slægten aldrig bringe
det. Ingen alvorlig Stræben for Andres Vel, ingen Deeltagelse
i deres Vee kan bevares hos den eftertænksomme Mand, naar
han anseer deres Tilværelse for et kortvarigt Skyggespil. Man
har ofte hørt, at En i Drømme har svedt Angestens Sveed over
den Maade, hvorpaa han har behandlet Mennesker, der kun vare
hans egne Drømmesyner, men har pludselig følt sig lettet ved at
komme til Vished om, at Alt var en Drøm og følgelig ligegyldigt. Men efter Tilintetgjørelseslæren bliver Menneskets Handlemaade mod hverandre for den, der kjender Livets Korthed, omtrent saa ligegyldig som en længere Drøm. En Livsanskuelse,
der fører til practisk Nihilisme gaaer let over til en positiv Stræben efter at tilintetgjøre sig selv. Derpaa findes en heel Deel
Exempler i Historien. Et paafaldende Exempel have vi blandt
andre i en forchristelig Periode, da en consequent dialectisk Udvikling af den cyrenaiske og cyniske Philosophie netop førte til en
practisk Nihilisme, der ei var Product af Enkeltes Griller, men
havde sin dybe Rod i Tidsalderens Udvikling. Derved vaktes en
saa almindelig Livslede, at Regjeringen i Alexandria forbød Hegesias, med Tilnavn Peisithanatos, at foredrage sine Læresætninger, fordi de foranledigede en saa usædvanlig Mængde Selvmord. Det samme Phænomen gjentager sig i vor Tid, da den
om sig gribende nihilistiske Livsanskuelse føder idelige Selvmord af
sig, og de eftertænksomme Statsphilosopher have fuldkommen Ret,
naar de fra alle Kanter gjentage, at den egentlige Grund maa
søges i den dogmatiske Fornægtelse af Udødeligheden. Selvmorderne ere virkelig ogsaa de meest consequente Lærlinge af en saadan Skole: enhver tænksom Tilhænger deraf synes at maatte
ønske med Caligula, at Menneskeheden kun havde eet Hoved for
% --- p. 52 ---
\opage{52}\setcounter{footnote}{0}%
at han med eet Sværdslag kunde befrie den fra sin urimelige
Livsdrøm.

Den, der kun tillægger de menneskelige Individer Betydning
som forbigaaende Former for den ene Aands Selvaabenbaring,
kan umulig uden Inconsequents nære nogen sand Menneskekjærlighed. De som desuagtet, uden at sand Menneskeagtelse har
nogen Grund i deres Livsanskuelse, lade de velvillige Følelser indtage stort Rum i deres Bevidsthed, indrømme derved de naturlige
Følelser en Magt, som deres Forstand ikke kan billige. I Almindelighed er det ogsaa Trang til at erkjende bestemte menneskelige Individer som uforgjængelige Gjenstande for deres Kjærlighed,
der fører tvivlende eller vantroe Tænkere tilbage til den christelige
Verdensanskuelse. Derfor have ogsaa de Videnskabsmænd, som i
deres Ungdom have bidraget meest til at give Naturphilosophien
sin Form, i deres Alderdom opgivet den abstracte Pantheisme, som
ligger til Grund derfor, og om en Deel af dem er det bekjendt,
at Dødsfald iblandt Individer, til hvem de have været knyttede
ved stærke sympathetiske Baand, have aabnet deres Øine for Tomheden i deres Verdensanskuelse. At saadanne Kjendsgjerninger
kunne træffes i Videnskabernes Rige, viser, at den systematiske
Fremstilling af en Verdensanskuelse ei er den sidste Grund, hvorpaa Overbeviisningen om dens Gyldighed beroer, ikke engang for
dem, som bedst kjende Systemet i hele dets Detail. Den, der
er i Stand til at forsone sig med Tilintetgjørelseslæren, har enten
ikke ved nogen Erfaring lært den uendelige Menneskekjærlighed at
kjende, eller han savner den anskuelige Forestilling om Livets Korthed eller ogsaa han har sammensat sin Bevidsthedskreds af stridige Brudstykker. Den Kjærlighed, der betragter sin Gjenstand
som forgjængelig, er nødvendigviis af en anden Art, end den
der veed, at den hører til det evig Bestaaende.

Den høiere Drift til Selvbevaring og Trang til Kjærlighed
ere som oftest de Grunde, der i Virkeligheden fortrænge de man%
% --- p. 53 ---
\opage{53}\setcounter{footnote}{0}%
gelfulde Systemer, som drage Udødelighedens Fornægtelse efter
sig. Men ogsaa den ideale eller contemplative Drift taber, ret
forstaaet, efter disse Systemer sin Betydning. Vi betragte nu
ogsaa fra denne tredie Side den Sandhed, at den gjennemførte
Negation af Udødeligheden fører til Opgivelsen af den Verdensanskuelse, som drager den efter sig. Jeg gjentager, at vi her ved
Abstraction adskille, hvad der i sig selv er Eet og det Samme;
ret betragtede er allerede saavel den sande Selvbevaring som den
sande Menneskekjærlighed forskjellige Former, under hvilke Driften
til at leve for det Evige træder frem. Dette er meer iøinefaldende ved den tredie ideale Retning, som vi her ogsaa ville betragte for sig; jeg mener den, der fører til Videnskab, Kunst og
Religion. Ved enhver af disse Former skulle vi ligeledes see, at
de ved en consequent Tilintetgjørelseslære tabe deres egentlige
Betydning.

Med dogmatisk Negation af Udødeligheden kan ingen Kunst
bestaae. For at Mennesket skal kunne have Svingkraft til at
bevæge sig i den høiere Phantasies Regioner, maa det enten med
en Slags reflexionsfri Udødelighed bruge Livet som en Evighed,
uden nogen anskuelig Forestilling om dets Korthed, eller det maa
% PRINTED AS IS: V215 q: crop, layer, tesseract -- a MISPRINT for Kraft, as printed
have Overbeviisning om Udødelighedsbegrebets Realitet. Jeg opstiller her uden Betænkning den Paastand, at Tvivl om Individernes Udødelighed er Kunstens Kræft og Overbeviisning om
deres Forgjængelighed er Kunstens Grav. Denne Paastand kan
kun retfærdiggjøres ved en Theorie af Kunstens Natur, som her
ei lader sig udføre, men kun med et Par Ord antyde. Den
Synsmaade, jeg her vil fremsætte, vil for mangen En klinge
høist paradox, og vil maaskee af Somme opfattes som et dristigt
Udtryk, der ei skal forstaaes efter Ordet. Ja mangen En har
maaskee selv udtalt den samme Tirade, og dog kun betragtet den
som en rask rhetorisk Figur; men her menes bogstavelig, hvad her
bliver sagt. „Den sande Kunst er en Anticipation af det salige
% --- p. 54 ---
\opage{54}\setcounter{footnote}{0}%
Liv.“ Den fuldkomne Harmonie mellem Universalitet og Individualitet, som finder Sted i Kunstnerens Bevidsthed, naar han
producerer et ægte Kunstværk eller hos den, der med ægte Receptivitet gjentager den samme Production, er et Billed af de fuldkommen Forklaredes Salighed, hvori det individuelle Liv uden
Modstand bliver udfyldt af den Eviges Villie. I den nærværende
Tingenes Orden er det kun i enkelte Momenter, den personlige
Villie hæver sig til det hellige Standpunct, hvor Pligtloven slet
ikke føles som en strængt bydende Tvang; og Videnskaben nødes
efter sin Natur til at arbeide saa strængt i Abstractionens
Sphære, at Intuitionen, som styrer Tænkningen, tit træder meget
i Baggrunden. Kuns i de Momenter, da den personlige Villie
falder sammen med den almindelige, eller Begrebernes discrete
Puncter flyde sammen til Intuitionens Klarhed, føler Mennesket den samme Harmonie i Sindet, som den, hvori den
sande Kunst altid bevæger sig. I den sande kunstneriske Production maa den personlige Stræben saaledes falde sammen med
den høiere Bistand, at Anstrængelsens Følelse reent maa forsvinde. Men Kunsten er kun et Billed paa det salige Liv, fordi
den bevæger sig i en Verden, der baade fra et Standpunct, der
ligger over den og et Standpunct, der ligger under den, med
Sandhed betragtes som Leeg og Skin. Kunstnerens Verden beholder ligeover for sig en Verden, der med Sandhed sættes imod
den som den virkelige, og som med Rette gjør Fordring paa med
Tiden at gjøre Kunsten overflødig. Det virkelige Liv vil engang
blive saaledes forherliget, at Kunstnerne have fuldbragt deres prophetiske Kald, og Individernes Liv umiddelbart selv fremstiller den
Harmonie mellem det Menneskelige og Guddommelige, der af
Kunsten bebudes som tilkommende. Disharmonien mellem den
prosaiske og den poetiske Verden skal i Tidens Fylde forsones i
den sande Tilværelse. Dette er en Sandhed, som Kunstnerne
stærkt bestride, især de herlige, der føle deres Sind aldeles ud%
% --- p. 55 ---
\opage{55}\setcounter{footnote}{0}%
fyldt ved deres Kald. Men de gjøre det i en Tilstand af elskværdig Naivetet, ligesom den Vilde, der haaber at gjenfinde sine
kjære Vaaben i de Dødes Rige.

Naar da Kunsten er en Anticipation af det mere ætheriske
Liv, som hører en høiere Tilværelse til, maa der naturligviis findes uforenelige Bestanddele i Bevidstheden, idet man paa eengang
vil beholde Kunsten og fornægte Udødeligheden, og den Kunst, der
saa aldeles miskjender sin Betydning, vil efterhaanden som de
stridige Antagelser bringes ret i Berøring med hinanden, aldeles
fortrænges ud af Livet, ligesom et Træ mister sine Blomster, naar
dets Rod bliver afskaaret. Det er da sandsynligt, at den ret
gjennemførte Negation først reent vil fortrænge Kunsten, inden
Livets Tomhed og Øde ved dette Savn ret klart kan blive stillet
Tidsalderen for Øie.

Naar Udødelighedstroen aldeles har forladt en Tidsalder,
forsvinder den sande Kunst meer og meer og den svage kunstneriske Production, som bliver tilbage, slutter sig for Størstedelen
til Ruiner af andre Tiders Frembringelser. Digterne føle idelig,
at en Verdensanskuelse, i Følge hvilken den nærværende Sandseverden skal betragtes som et i sig afsluttet Heelt, ikke kan tilfredsstille Poesiens Fordringer. Derfor have de meest udmærkede
Digtere saa tit givet deres episke og dramatiske Digtninger en heterogen
Tilsætning af Syner fra en anden Verden, som ikke staaer i
Harmonie med det Hele. Hvoraf kommer denne Disharmonie,
uden deraf, at Brudstykker af to heterogene Tidsaldres Producter
ere mechanisk forbundne med hinanden? Dersom Bevidstheden om
en fra den nærværende Verden qualitativ forskjellig Verden, som
engang kan blive tilgjængelig for en aldeles ny Erfaring, ret
havde gjennemtrængt Tidsalderen, saaledes at den i den nærværende Tingenes Orden betragtede alle Ting som Fragmenter af et
organisk sammenhængende Universum, der her kun begynder at udvikle sig for os, da kunde den fremføre Phantasier om en hiin%
% --- p. 56 ---
\opage{56}\setcounter{footnote}{0}%
sidig Verden, der kunde harmonere med den bekjendte Verden; thi
den nærværende Verden fik da selv et andet Udseende.

Paa den meest barnagtige Maade har Savnet af Tro paa
det Oversandselige viist sig i de mange Romaner, der til en Tid
kom for Lyset, hvis hele Indhold var tilsyneladende Spøgelseshistorier, der til Slutning ved en naturlig Forklaring opløstes i Intet. Ogsaa de mange Digterphantasier, der bevæge sig
i en reen mythologisk Verden, skylde deres Oprindelse til samme
Trang, men de betragtes gjerne af Forfatteren selv som en betydningsløs Leeg, og den urigtige Forestilling om Kunsten maa
naturligviis have Indflydelse paa Kunstens Udøvelse. Thi saa
vist det er, at Kunstnere kunne producere høist geniale Værker
uden selv at forstaae deres Betydning, saa vist er det, at en
Misforstaaelse af Kunstens Natur, saadan som den nødvendigviis
følger af skjæv Cultur, maa virke fordærvelig tilbage paa den
kunstneriske Virksomhed. I Digte, der nuomstunder føre en Feeverden frem for Phantasien, er det ofte den vantroe Forfatters
og Læsers Mening, at man for et Øieblik skal lade sig hentrylle
i en oversandselig Verden og dog beholde den Bevidsthed, at der
ingen høiere Verden gives. Men Følgen deraf er, at Fortryllelsen bliver derefter. Dette forsøgte Selvbedrag er mere barnagtigt end Barnets, der lader Skammel og Alen gjælde for
Throne og Scepter, uagtet det tillige har Bevidsthed om, at det
er en Skammel og en Alen; thi Barnet veed dog, at der virkelig
gives Throne og Scepter, men den Voxne troer, at det hans
Feeverden betyder, er aldeles Intet. Et Trin af Kunsten er det,
paa hvilket Digteren udtaler mythologiske Traditioner, paa hvilke
han selv troer: et andet det, paa hvilket hans Phantasier, om de
endogsaa hente Motivet fra slige Sagn, betragtes af ham selv
som en betydningsløs Leeg: et tredie det, paa hvilket han betragter sin kunstneriske Production som Billedet paa en høierr Tilværelse. Paa det mellemste Udviklingstrin hører Kunsten efter%
% --- p. 57 ---
\opage{57}\setcounter{footnote}{0}%
haanden op at være det, som Digteren ikke længer antager den
for at være, et Forvarsel om en høiere Tilværelse. Men i en
saadan øde Tid, en prosaisk Tid i den allerværste Betydning af
Ordet, er det yderst mærkeligt, hvorledes de meest blasphemiske
Digtere mod deres Villie tvinges til at udtale en anden Aand,
end den, de troe paa. Deres formeentlige Bespottelse af alle positive Former i Tilværelsen bliver til en Klagesang over det bortflygtede høiere Liv og til et Længselssuk efter en poetisk Fremtid.
Især ere H. Heines Skrifter i denne Henseende høist mærkværdige, fordi hans foregivne Henrykkelse over den isolerede Sandseverdens Tilstrækkelighed for et sundt Menneskes Fordringer, idelig
afbrydes ved en hemmelig Veemods Udbrud. Ligeledes pynter han
idelig sine Spottetaler med Glimmerblomster fra en modsat Verdensanskuelses efterladte romantiske Former. Det gaaer overhovedet det unge Tydskland og det unge Frankrig saaledes, at al
deres Jubel over Kjødets Evangelium, som Heine kalder det, og
al deres paatagne Lystighed mod deres Villie bliver til smertelige
Klager. Det gaaer dem ligesom Penelopes Beilere, da de holdt
deres sidste muntre Aftensmaaltid: „De loe, som Homer siger,
med fremmede Ansigter; de aade Kjød blandet med Blod; deres
Øine fyldtes med Taarer og deres Sind tænkte paa Jammer.“
I saadanne Phænomener er det at den verdenshistoriske Dialectik
lægger sig for Dagen: Nihilismen optræder, derved at dens Tungemaal bliver forvendt, som sin egen ufrivillige Anklager og baner
Veien for den positive Verdensanskuelses Restitution.

Heller ikke Videnskaben kan i Tidens Længde beholde sin
Interesse med Tilintetgjørelseslæren eller den eensidige Pantheisme:
den kan hverken bestaae med Pantheismen i sin trivielle Form,
hvori de literære Plebeier opfatte den, eller i den fordringsfulde Skikkelse, som den antog i den logiske Pantheisme. Efter disse Synsmaader
antages det nemlig paa den ene Side for muligt, at Mennesket kan opnaae fuldkommen Erkjendelse, paa den anden Side antages det for en
afgjort Sag, at Menneskets Liv er indskrænket til den nærværende Tin%
% --- p. 58 ---
\opage{58}\setcounter{footnote}{0}%
genes Orden. Disse to stridige Sætninger bringes i Harmonie
derved, at endeel abstracte Sætninger udgives for den fuldkomne
Erkjendelse. For at Mennesket kan glæde sig ved den Overbeviisning at vide Alt, antages at det, der synes at falde udenfor
dets Videns Horizont, ikke har nogen Realitet.

At fuldendt Viden ikke kan opnaaes under Menneskehedens
nuværende Livskaar, kan indsees af følgende Betragtninger. De
empiriske Videnskaber kunne efter Sagens Natur ei tilendebringe
deres Bestræbelse for at sammenfatte de existerende Individer og
deres Tilværelsesmaader i et fuldstændigt Begrebssystem. Det tilstaaes vel af Enhver, der blot har et dunkelt Begreb om disse Videnskabers Fremgangsmaade. Man behøver jo desuden blot at
erindre sig, at disse Videnskabers Fordringer ei kunne tilfredsstilles
ved en fuldstændig Angivelse af alle telluriske Tilværelsesformer;
thi da bliver jo Resten af Himmelsystemet med sit respective Inventarium tilbage som en utilgjængelig Region. Men lader os
for et Øieblik endda blive staaende med vor Betragtning ved denne
Planet og antage, at Erfaringens Videnskaber kunne bringes saa
vidt, at de paa deres Viis kunne give fuldstændig Beskeed om
Jordens Organisation, Former og Love; da mangler jo dette System en væsentlig Side for at fortjene Navn af fuldendt Erkjendelse. Anskuelsen, der paa ethvert Punct er af uendelig Natur,
savnes aldeles i et saadant System af Begreber. Dette følger af
den almindelige Sætning, at enhver enkelt Anskuelse kun i abstract
Form kan gaae over til Forestilling, saa at Almeenbegreberne, selv
om Erindringen af den tilfældige Erfaringskreds, hvortil de referere sig, beholdes uden Fradrag, nødvendigviis maa være af abstract Natur. Uden her at opholde os ved det rige Indhold, der
gaaer tabt i Naturvidenskabens øvrige Sphærer, ville vi her tænke
paa den Deel af Tilværelsen, som har nærmest Sammenhæng
med vor Hovedundersøgelse: vi bringe i Erindring, hvilket uende%
% --- p. 59 ---
\opage{59}\setcounter{footnote}{0}%
ligt Indhold Anthropologien maa oversee, naar den taler om Arter og Slægter i Stedet for om Individer.

Vil Videnskabsmanden derimod vælge det Almindeligste til
sit Udgangspunct og ved en apriorisk Tænkning skride frem fra
det Abstracteste til det mindre Abstracte, uden at forlade den rene
Tænknings Region, da er han vistnok i Stand til stedse meer og
meer at orientere sig i hine dybe Grundbegreber, der i deres umiddelbare Form fremstille Fragmenter af en oprindelig Erkjendelse.
Men det er kun Tilværelsens Grundforhold, der ved denne Fremgangsmaade kunne bringes til tydelig Bevidsthed. Det hele Universum med sin Uendelighed af Bestemmelser kan aldrig sammensættes af aprioriske Tanker.

Antages Muligheden af, at Mennesket i den nærværende
Tingenes Orden kan komme til en fuldkommen Erkjendelse, da
ere de Videnskabsmænd meest conseqvente, der skildre den som en
mystisk Anskuelse af Eet og Alt, hvori Modsætningen mellem den
Anskuende og det Anskuede er aldeles forsvunden. At en saadan
Anskuelse af det Absolute skulde være en fuldkommen Erkjendelse
deraf, er en Paastand, hvis Urimelighed ved saa mange sindrige
Vendinger af Polemiken er godtgjort, at det neppe nuomstunder
falder en eneste grundig Videnskabsmand ind at paatage sig dens
Forsvar. Men enhver apriorisk eller om man vil immanent
Tænkning, hvorved Tilværelsens Grundforhold, (jeg mener Guds
Forhold til Universet) bliver tydeliggjort, maa dog gaae ud fra en
dunkel Anskuelse af disse Grundforhold og stræbe til en klarere
umiddelbar Erkjendelse deraf. Men for saavidt som en adæquat
Anskuelse af Eet og Alt i den nærværende Tingenes Orden ei
kan opnaaes af Nogen, kan den videnskabelige Tænkning under
de os bekjendte Forhold heller aldrig føre til Erkjendelsens Fuldendelse.

Kun ved een Vending kan det antages for muligt, at Mennesket i Sandseverdenen kan opnaae fuldkommen Erkjendelse: det
% --- p. 60 ---
\opage{60}\setcounter{footnote}{0}%
er nemlig ved den Tilsnigelse, at man frakjender enhver Ting
Prædicatet af Sandhed, naar den er utilgjængelig for den menneskelige Erkjendelse. Man beskjærer saaledes paa en vilkaarlig Maade
Tilværelsen for at læmpe den efter den menneskelige Erkjendelses
Leilighed, ligesom hiin Procrustes, der forkortede sine Gjæster efter
Maalet af sit Sengested. Sandheden af denne Paastand bliver
især indlysende ved følgende Betragtning. Hvad der skal blive
ganske gjennemsigtigt for det videnskabelige Begreb, kan efter Sagens Natur kun være en almindelig, uforanderlig Side af Tilværelsen. Vil man nu antage, at Erkjendelsen heraf ei er abstract,
da maa det Individuelle, forsaavidt det er individuelt, betragtes
som en Mængde tilfældige Former for det Absolutes Tilsyneladelse. Skal det Individuelle, hvorunder de eenlige menneskelige
Individer henhøre, der fødes og døe til bestemte Timer paa Jordkloden, skal dette Individuelle have sand Realitet, da kan ingen
af de os bekjendte videnskabelige Former fremstille fuldkommen Erkjendelse. Allerede Sprogets Natur gjør det umuligt paa en videnskabelig Maade at betegne disse bestemte Individer. Men der
maa gives en tredie høiere Erkjendelse, som hverken holder sig
blot i de rene Begrebers Region eller bliver staaende ved den empiriske
Tænknings Abstractioner, men hvori paa en Maade, som overstiger
menneskeligt Begreb, det Almindelige saaledes erkjendes, at det
Enkelte (Eenlige) ikke forsvinder deri. En saadan Erkjendelse
vilde være en intellectuel Anskuelse i Ordets egentlige Betydning,
men dette sande „\textit{tertium cognitionis genus}“ er uopnaaeligt
under dette Livs Vilkaar. Derfor er det jeg lidet ovenfor har
sagt, at Fornægtelsen af Menneskets Udødelighed blandt Andet af
den Grund vil høre op, at man seent eller tidlig vil komme til
en klar almindelig Overbeviisning om, at fuldendt Sandhedserkjendelse efter denne Synsmaade for bestandig er Mennesket umulig.

At ingen sand Religion kan bestaae med den eensidige Pantheisme, som vi overalt i denne Afhandling have betragtet som en
% --- p. 61 ---
\opage{61}\setcounter{footnote}{0}%
Tilintetgjørelseslære, er om muligt endnu mere indlysende, end at
Kunsten og Videnskaben fra dette Standpunct tabe deres Betydning. Det forekommer rigtignok En og Anden, at Den var kommen til et ophøiet Stade, som saaledes havde fordybet sig i Beskuelsen af det Almindelige, at han med Resignation rolig saa sin
ubetydelige Personlighed forsvinde som et Offer for det Almindeliges Aabenbaring. Men dette Almindelige kan ikke substitueres
for den personlige Gud, hvis Begreb bestandig bevares i den sande
religiøse Tradition. Det er et Væsen, som ikke har nogen fuldendt Erkjendelse og ingen Villie i Ordets rette Betydning. Det
gjennemløber i en instinctmæssig Stræben en Række Forvandlinger, indtil det i Menneskeformen opnaaer et Glimt af Selvbevidsthed. Den menneskelige Erkjendelse er den høieste Erkjendelse,
der gives i Universet; thi Menneskets Erkjendelse af Gud er Guds
eneste Erkjendelse af sig selv. Den guddommelige Erkjendelse
savner da i enhver Henseende Continuitet, thi deels er den absplittet i coexistende monadiske Selvbevidstheder, der ere ugjennemsigtige for hinanden, deels glimter den frem i successive Puncter,
som uophørlig slukkes igjen for bestandig. Continuitet har da dette
store bevidstløse Dyr, som vi kalde Gud, kun i sin Naturgrund,
hvor det med blind Almagt sammenholder den vrimlende Mangfoldighed af forgjængelige Væsner, i hvilke det deler sig selv.
Men Magten og Erkjendelsen lade sig aldrig forene; thi naar
der hist og her hos Væsnet i dets Menneskeform vaagner en dunkel Erindring om hvad det er, da har det i denne forgjængelige
Skikkelse aldeles tabt Almagtens Scepter af sin Haand. Dette
er Hovedindholdet af den store orientalske Verdensfabel, som Mennesket i det sidste halve Seculum har udarbeidet i saa mange
Consequentser. Men det halv levende Væsen, som saaledes i
Drømme opfører Verdens Drama for sig selv, kan naturligviis
ikke være Gienstand for den religiøse Contemplation efter Ordets
rette Betydning.

% --- p. 62 ---
\opage{62}\setcounter{footnote}{0}%
Den Verdensanskuelse, som fremstilles i denne Mythologie,
som nu udgjør det videnskabelige Europas herskende Tro, lader
sig paa mange Maader, uden stor Vanskelighed, dølge under kunstige Former, hvori selv alle Christendommens Hovedsætninger bevare en tilsyneladende Gyldighed. Men selv Tilhængerne af meget
forviklede Systemer, som gjennem flere Aar have maattet arbeide
sig gjennem en sindrig Væv af abstracte Begreber, røbe dog jævnlig, at de i den ovenomtalte Mythologie finde det naturligste Udtryk for deres Verdensanskuelse. Saaledes har nylig K. L. Michelet, der ingenlunde høre til de Mænd af den hegelske Skole, % NOTE: the top of the d is worn (crop b, machines v); Udtryk
der tænke frit paa egen Haand, i en Recension af et lidet Skrift,
der i usædvanlig simpel Form fremstiller det nittende Aarhundredes Verdensmythe, erklæret, at dette Skrift i de væsentlige Resultater var fuldkommen overeensstemmende med den hegelske Philosophie.

Efter den eensidige Pantheisme, der efter sin Natur er uadskillelig fra Tilintetgjørelseslæren, er fuldkommen Erkjendelse ikke
blot uopnaaelig for Menneskeslægten, men der gives ingen saadan
Erkjendelse, hverken i Himmelen eller paa Jorden. Verden bevæger sig i et uophørligt Kredsløb, efter en streng Nødvendigheds
Love fra Intet og til Intet. --- Denne nihilistiske Side ved den
moderne Pantheisme findes neppe nogensteds gjennemført med den
Kraft og Følgerigtighed, som i et Værk af Arthur Schopenhauer\footnote{Die Welt als Wille und Vorstellung.}, Professor i Berlin, som, skjendt det er udkommet i
Aaret 1819, endnu sjelden blot løselig bliver omtalt af andre
Philosopher. Forfatteren betegner selv i de meest frimodige Udtryk sin Philosophie som antichristelig og nihilistisk. Saaledes
sætter han flere indiske Myther høit over Christendommens Dogmer, og roser meget den braminske Lære om de Saliges Opholds%
% --- p. 63 ---
\opage{63}\setcounter{footnote}{0}%
sted efter Døden, fordi det betegnes med Ordet „Niban,“ som
egentlig betyder „Intet.“ Hvis dette Værk, som vist ikke tilfældigviis er kommet til at staae i Literaturens Baggrund, var meer
almindelig bekjendt, vilde det vist bidrage meget til at gjøre den
Nihilisme, som den eensidige Pantheisme fører til, ret anskuelig
for Tidsalderen.

Saaledes har jeg ved flere fragmentariske Vink aabnet en
Udsigt til den umenneskelige Livsanskuelse, hvortil den udbredte
Tilintetgjørelseslære ufeilbarlig maa føre sine Bekjendere. Det er
den Sinai-Ørken, hvori Menneskeslægten har fordybet sig næsten
i 40 Aar; men den maa gjennemvandres til alle Sider, før det
er muligt at komme til Bjergene, hvorfra man kan øine det forjættede Land. I Videnskaben er fra mange Sider begyndt en
kraftig Reaction mod Nihilismen, men der vil maaskee gaae lang
Tid med, inden det gamle Systems servile Tilhængere blive lede
ved deres egen Verdensanskuelse. Først maa et langt Erfaringsbeviis overtyde Tidsalderen om, at det er Mennesket umuligt,
naar det ret befinder sig, at at leve med hine Meninger; dernæst
maa den erindre sig, at den hele Synsmaade beroer paa gandske
vilkaarlige Forudsætninger. Næsten fra Barnsbeen af har den
yngre Slægt, saavel i Conversation som af Skrifter, optaget det
ene Tankefrø efter det andet, som indeholdt forskjellige Spirer til
hiin store Kreds af Fordomme. Derfor falder det Tidsalderen
vanskeligt, til Trods for store Autoriteter og formeentlige videnskabelige Beviser, som den endog ubeseete holder for fuldkommen
afgjørende, dristig at sige sig selv, at den egentlig slet ingen Borgen har for hine Fordommes Gyldighed, men at den kan tage
Mod til sig og forkaste dem, hvilket Øieblik den vil.

Ovenstaaende aphoristiske Tanker, hvormed jeg har søgt at
udpege de forskjellige Sider, hvorfra Nihilismen klart viser sig i
sin sande Skikkelse, ere dog ikke nedskrevne reent uden Plan. Jeg
har nemlig gjort opmærksom paa, at hverken Menneskets Kjærlig%
% --- p. 64 ---
\opage{64}\setcounter{footnote}{0}%
hed til sig selv eller til andre Mennesker efter denne Verdensanskuelse kan beholde sin Betydning; at baade Kunst og Videnskab og Religion paa dette Standpunct efterhaanden maa tabe
deres Interesse. Men under disse Begreber kan Alt indbefattes,
hvad der er i Stand til at udfylde den menneskelige Bevidsthed,
og naar Mennesket taber Interesse for alt Livets Indhold, kan
det ikke uden Selvmodsigelse fortsætte sin tomme Tilværelse. At mangfoldige Mennesker uden Tro paa Udødelighed, desuagtet have lagt stor
Livslyst for Dagen, maa for det Meste forklares af en Slags Tankeløshed, i Følge hvilken en heel Deel Mennesker kun maadelig ere orienterede i deres egen Forestillingskreds. Men nu har Menneskeslægten faaet
en melancholsk Retning, der hverken vil tillade den at lade sine
Selvmodsigelser ubemærkede, eller bestandig at glemme sig selv i
en betydningsløs Travlhed. At selv dybsindige Philosopher, som
ikke kunne undgaae at bringe hine Modsigelser til klar Bevidsthed,
desuagtet kunne bevare deres Levelyst (for ei at sige Livslyst) har
sin Grund i, at de idelig tabe sig i deres Anstrængelse for at
give den falske Verdensanskuelse en videnskabelig Form. Dette
Arbeid, der ofte blot er en anstrængende Leeg med vilkaarlige Former, fylder deres hele Sind og bringer dem til en Selvforglemmelse, der opholder deres Liv. Men egentlig ere de Tidsaandens
blinde Redskaber, der arbeide paa Negationens Fuldbringelse, for
at Nihilismen i alle sine Former kan blive bragt til Tidsalderens
Bevidsthed. Naar denne Videnskabens Retning er fuldbragt, naar
intet speculativt Talent længer føler sig tilfredsstillet ved hiint Arbeid, fordi de ikke længer her møde nogen væsentlig Opgave for
deres Kræfter, da vil hele Videnskabens Virksomhed gaae ud paa
at give den sande Verdensanskuelse en videnskabelig Form. Ved
den sande Verdensanskuelse forstaaer jeg her den gamle, som i
Videnskaben er ny, den i hvilken Læren om en personlig Gud og
Menneskets Udødelighed er et Hovedpunct.

At Videnskaben paa den Vei, den havde indslaaet, vilde ende
% --- p. 65 ---
\opage{65}\setcounter{footnote}{0}%
ved et nihilistisk Resultat havde Schiller med en mærkelig Forudfølelse udtalt i sit bekjendte Digt, det tilslørede Billed i Sais.
Dette Digt, som Samtiden og Digteren selv vist ikke tilfulde
forstod, har nu i Videnskabens Historie faaet en oplysende Commentar. Der fortælles, at en dristig Yngling løftede Sandhedens
Slør, men saae da et Syn, som betog ham al sin Livslyst, saa
at han efter den Tid gik omkring som en levende Død. Hvad
han saae, eller rettere hvad Schiller med sit prophetiske Blik
saae, fortælles ei, men det var den frembrydende Nihilisme, Braminernes Niban.

Inden jeg forlader den Betragtning, at den fuldstændige
Bevidsthed om Nihilismens Natur vil bringe Tidsalderen til at
forlade den, vil jeg gjøre et Par her uvedkommende Bemærkninger, blot for at forebygge Mistydning af, hvad her er sagt.
Det har aldrig været min Mening, at den Revolution, som har
taget sin Begyndelse i Videnskaben, skulde være radical i den Betydning, at Menneskehedens Aand med Eet skulde have tabt Udbyttet af mange Aarhundreders Arbeide, eller at en ny-Periode
skulde begynde uden historisk Sammenhæng med Fortiden. Videnskaben kan ikke forstaae sit Standpunct uden idelig at recapitulere
den Vandring, som har ført dertil. Ved den Betragtning, at
Forfølgelsen af den eensidige Pantheisme fører til Nihilisme, bliver ei enhver Pantheisme forkastet som blot Vildfarelse. Philosophiens nærmeste Opgave er for Tiden at forsone den erhvervede
Bevidsthed om et Væsen, der et Eet og Alt, med Bevidstheden
om dette Væsens Modsætning til uforgjængelige Individer. Men
begge Modsætningens Led komme ei til deres fulde Ret, før Læren
om en personlig Gud og udødelige Mennesker, saavel i videnskabelige som populære Former, vinder det Herredømme, som den
eensidige Pantheisme for Tiden er i Besiddelse af. Alt ovenfor
er det Samme bemærket, nemlig, at den sande Lære om en personlig Gud hviler paa Pantheismens Grundvold.

% --- p. 66 ---
\opage{66}\setcounter{footnote}{0}%
Dernæst maa den Bemærkning gjøres, at det theoretiske og
practiske Beviis for, at Menneskeslægten ei kan leve med den
gjennemførte Tilintetgjørelseslære, hvilket Tidsalderen i saa mange
Skikkelser har oplevet, langtfra er det fuldstændige Beviis for
Slægtens Udødelighed. Det var jo meget tænkeligt, at det netop
var Menneskets Bestemmelse --- hvis ellers teleologiske Tanker i
denne Sammenhæng kunne have nogen Betydning --- frivilligt at
give Afkald paa Livets Blændværk og at vende tilbage til det
rene Intet. Denne Opgivelse af „Villien til Liver“ anseer Schopenhauer i sit Hovedværk for Menneskeslægtens Befrielse og udtaler saaledes med større Sandhed den eensidige Pantheismes Følger, end de mange Philosopher, der bringe dem i en skuffende
Harmonie med Borgerlivets og Religionens Ideer. --- Dog kan
man ikke nægte, at den Overbeviisning, at Mennesket kun ved
Selvbedrag og Selvforglemmelse kan leve med Tro paa Nihilismen, er et temmeligt stærkt \textit{argumentum ad hominem} imod
dens Gyldighed. Men det fuldstændige Beviis for Udødeligheden
er den videnskabelige Udførelse af den sande Verdensanskuelse,
hvori Udødelighedslæren er et Hovedstykke. Dette Beviis er paa
flere Kanter i sin Udvikling, men det er kun i sin første Begyndelse, saalænge det kun er Enkeltes Værk. Det er endnu ikke
saaledes blevet Tiden en Hovedsag, at det snart kan fortrænge
den nihilistiske Pantheisme, som taler med tusinde Tungemaal.
Men Trangen til en videnskabelig Fremstilling af den gamle Tradition, som det sidste halve Seculum er afveget usædvanlig langt
fra, lader sig paa ingen Maade tilfredsstille ved Henviisning til
Christendommen i sin umiddelbare Skikkelse. Især er denne Vending frugtesløs, hvis man derved tænker paa et i faa Sætninger
affattet Symbol, som Udtryk for Christendommens Hovedindhold.
At troe paa isolerede Sætninger, som man har seet fuldstændiggjorte ved mangfoldige stridige Systemer, er ikke engang muligt
uden at erklære sig for den ene eller anden af hine Udviklinger.
% --- p. 67 ---
\opage{67}\setcounter{footnote}{0}%
Hvor liden Indskrænkning det medfører for Tankens Frihed, at
en Philosoph holder sig strængt til en Troesbekjendelses Bogstav,
har man nylig seet et mærkeligt Exempel paa. Marheinecke
har nemlig udarbeidet en Dogmatik i hegelske Former, som nøiagtig holder sig til den apostoliske Symbolums Ord, og deri savnes
desuagtet den christelige Udødelighedslære. Det er en svær Opgave for en Troende, der udtaler sin Tro i populære Formularer,
at blive uenig med en Tilhænger af det moderne System, naar
denne ret lægger an paa at holde den gode Forstaaelse vedlige.
Det er kun i Videnskabens Region, at videnskabelige Vildfarelser
kunne finde deres Gjendrivelse.

\begin{center}\rule{3em}{0.4pt}\end{center}

\begin{center}VII.\end{center}

I de mange Skrifter om den personlige Udødelighed, som i
den senere Tid have seet Lyset, have somme Forfattere vovet sig
dristig ud paa de løse Gisningers Mark og fremført en heel
Deel udførlige Bestemmelser for den nye Tingenes Orden, som
for Mennesket skal afløse Livet i de fem Sandsers' Verden.
Andre have derimod med saa stor Angstelighed frakjendt det tilkommende Liv ethvert bestemt Prædicat, at man, hvis det ikke
modsagdes ved deres hele Tankegang, kunne fristes til at troe, at
de vilde reducere alt det Hiinsidige til Braminernes omtalte Niban. Det meest almindelige Udtryk for denne Retning findes i
den bekjendte metaphysiske Fordom, at Tidens og Rummets Former ikke have nogen Anvendelighed paa Menneskets tilkommende
Bevidsthedskreds. Denne underlige Lære, som man jævnlig møder
selv hos Nutidens meest fordomsfrie Philosopher, er formodentlig
endnu en Levning af den bekjendte Sætning i den kritiske Philosophie, at de menneskelige Erkjendelsesformer blot have subjectiv
Betydning. Men den kritiske Philosophie gaaer i saa Henseende
meget mere consequent tilværks, idet den ikke blot betragter Tid
% --- p. 68 ---
\opage{68}\setcounter{footnote}{0}%
og Rum som subjective Former, men ligesaafuldt nægter Begreberne Virkelighed og Tilværelse deres objective Gyldighed. Nutidens halvidealistiske Forsvarere af Udødelighedslæren indvikle sig
derimod i Modsigelser, som klart kunne lægges for Dagen, ved
en dialectisk Udvikling af de almindelige Bestemmelser ved Begreberne Virkelighed og Tilværelse. En virkelig Verden uden Tid
og Rum er en fiirkantet Cirkel og en rund Triangel. Endogsaa
uden nogen metaphysisk Tydeliggjørelse kan det umiddelbart indsees;
at man forener uforenelige Begreber, naar man tænker sig en
Fleerhed af levende Individer, som ikke ere udenfor hinanden,
eller et virkeligt Liv, hvori eet Moment ei følger efter et andet.

Af disse Misgreb er det dog især det første, der er meest
til Hinder for den sande Gjenfødelse af Udødelighedslæren. Sindige Videnskabsmænd maae nødvendig med Ringeagt betragte hiin
Dristighed i at fremføre løse Formodninger, som i den senere
Tid igjen paa flere Kanter begynder at gjøre sig gjældende. Naar
det kommer an paa at give bestemte Topographier over det for
Mennesket under sine nærværende Livskaar aldeles utilgjængelige
Hiinsidige, har naturligviis enhver uvidende Sværmer ligesaa
megen Ret til at kundgjøre sin Mening, som den meest kundskabsrige Forsker. Saadanne Bestræbelser af ledige Hoveder for
at udfylde Tilværelsens ubekjendte Regioner med deres egne frie
Phantasier, bidrage meget til at fortrænge Bevidstheden om det
Oversandselige, idet alvorlige Mænd med Aversion vende sig bort
fra en Sphære, hvor man blot har med tomme Muligheder at
gjøre, uden nogensteds at finde fast Fod for Videnskab og rolig
Overbeviisning.

Aldeles uden Betydning kunne vistnok Phantasier om den
oversandselige Tilværelses Beskaffenhed paa ingen Maade siges at
være. Det vil altid blive Mennesket en Tvang med subjective
Phantasier at fylde hiin dunkle Region af sin Bevidsthed, for
derved ligesom at minde sig selv om, at det der finder en ledig
% --- p. 69 ---
\opage{69}\setcounter{footnote}{0}%
Plads, som upaatvivlelig maa betegne en virkelig Gjenstand for
tilkommende Erfaring. Det synes næsten umuligt at holde fast
paa den levende Vished om Realiteten af det store savnede Supplement til Menneskelivet, uden derved at man idelig fremstiller og
forkaster bestemte Billeder deraf. Disse Positionens og Negationens Roller blive snart fordeelte til flere Personer saaledes, at
den Ene godtgjør Ugyldigheden af den Andens Digtninger: snart
er det et eenligt Subject, som selv opbygger og selv nedriver de
Luftslotte, hvormed han udfylder sin ubestemte Erkjendelse af det
Hiinsidige. Saaledes opstiller Plato i sin Phædon en udførlig
Topographie af de Dødes Rige, men føier derefter den Erklæring
til, at det vilde være en Daarlighed om Nogen vilde paastaae,
at Sagen netop forholdt sig saaledes. Den maa, efter hans Mening, enten forholde sig paa hiin Maade eller paa en anden lignende. --- Forresten skal det ingenlunde være nægtet, at Phantasier over Aandeverdenen kunne indeholde dybe Genieblik, fordi Videnskaben ei kan forklare deres Indhold. At der i denne Region
ikke hersker blot Vilkaarlighed, følger allerede deraf, at ingen sindig
Kritiker vil erklære alle slige Phantasier for at have samme
Værdi. At de Feste ogsaa her med saa stor Enighed skjælne
mellem Rimelighed og Urimelighed, viser noksom, at de medbringe
Fordringer, som ikke blot have deres Udspring af den Enkeltes
Subjectivitet. Men en staaende Billedverden er her mindre end
noget andet Sted paa sin rette Plads. Hvad her i Sandhed skal
bringe en Tidsalder til almindelig Henrykkelse, maa være Tidsalderens egen Frembringelse. Derfor vil det for Tiden neppe lykkes nogen Digter eller Sværmer at fremføre Phantasier over det
Oversandselige, der møde vidt udbredt Sympathie hos deres Samtid. Det er kun i gamle Eventyr vi møde Pigebørn, som midt
i Vinteren bringe Roser fra Paradiis.

Den frodige Production af urimelige Fortællinger om Syner
% --- p. 70 ---
\opage{70}\setcounter{footnote}{0}%
fra en anden Verden, som for Tiden viser sig paa mange Steder,
er en forvirret Reaction mod Videnskabens Anmasselser, som paa
den anden Side med saa stor Frihed nægter Realiteten af Alt,
hvad der falder uden for dens Synskreds. Det er et mærkeligt
Særsyn, at selv det berlinske Tidsskrift, hvis nærmeste Formaal
er at virke for Udbredelsen af Hegels Philosophie, hvorefter alle
„Forestillinger“ om det Oversandselige betragtes som ufuldkomne
Former af Bevidstheden, i den senere Tid af og til har indrømmet Plads for alvorlige Undersøgelser om foregivne Visioner.

Men hvad man end dømmer om Betydningen af poetiske
Fictioner om Aandeverdenen, bliver det en afgjort Sag, at de
ikke have Noget at gjøre med det videnskabelige Begreb om de
tilkommende Forhold, hvorunder Mennesket skal fortsætte sit Liv
efter Dødens Forvandling. For Tiden kan Videnskaben ikke udtale synderlig andre Prædicater derom end dem, der følge af Begrebet \textbf{virkelig Tilværelse}, og det er ikke sandsynligt, at Mennesket i den nærværende Tingenes Orden nogensinde vil kunne
hæve sig til den aldeles nye Erkjendemaade, som gjælder for det
Hiinsidige. Tvertimod synes det at ligge i Sagens Natur, at
den fuldkomne Erkjendelses Gjennembrud netop maatte medføre en
saadan Forvandling af Tilværelsens Former, at det Hiinsidige var
opnaaet. (Her, som overalt i denne Afhandling, tænkes paa et
reelt Hiinsidigt, der er Gjenstand for en høiere Erfaring, som er
den sande Forening af Begreb og Anskuelse. Det er mig ikke
ubekjendt at endeel Philosopher holde for, at det netop er den her
betegnede Erkjendelse, der er faldet i deres Lod: Men Tiden har
opklaret deres Selvbedrag. De bevæge sig i en Kreds af abstracte Begreber, hvor Anskuelsen tabes af Sigte og Menneskelivets vigtigste Problemer betragtes som forbigaaende Illusioner,
der forsvinde for den rene Videnskab.)

Det er nødvendigt at skjælne Phantasiens Raaderum ved
% --- p. 71 ---
\opage{71}\setcounter{footnote}{0}%
skarpe Grændser fra Videnskabens Indhold, for at dette gandske
kan holdes frit for Skinnet af Vilkaarlighed. Hvad der kan
vides om den tilkommende Tingenes Orden, er ikke stort meer
% PRINTED AS IS: V298 q: crop, layer, tesseract -- a MISPRINT for det, as printed
end dens Tilværelse; thi Forestillingen om det ubestemte Hiinsidige skal ikke give Menneskelivet dets egentlige Rigdom af Indhold. Men Visheden om drt Tilkommende skal vise det Nærværende, som nærmest giver det sunde Menneskeliv sin Fylde, i
sit rigtige Lys. Naar det nærværende Liv betragtes som et
Fragment af et større Heelt, bliver det klart, at hverken den eensidige Tanke om Her eller den eensidige Tanke om Hisset medfører den sande Bevidsthed om Tilværelsens Uendelighed. Ved at
tage det nærværende Liv som en Deel af en Evighed lever man
alt her i een Betydning et uendeligt Liv, og kun ved at opleve
dette Livs høieste Momenter kan man danne sig en Forestilling
om en intensiv Uendelighed af et tilkommende Liv. Tanken om
det blot Hiinsidige kan da ikke byde den rolige Videnskab, som
forkaster Phantasiens Fostre, noget riigt Indhold til Bearbeidelse.
Den der lever blot for Betragtning af den for sig tagne Evighed uden at agte paa dens Reflex i den nærværende Tid, fører
et ligesaa tomt Liv, som den, der holder sig til den fra Evigheden
løsrevne Nutid som den egentlige Realitet.

Jeg gjentager, at Videnskaben ikke skal søge sit Indhold i
det afsondrede Hiinsidige; endnu mindre skal den eensidige Forestilling derom udfylde Menneskets Liv, men naar Bevidstheden
om det Hiinsidige tabes, og den Vildfarelse følgelig opstaaer, at
det nærværende Liv i sin isolerede Skikkelse er Indbegrebet af
det virkelige Liv, da maa Videnskaben bestræbe sig for at gjøre
den tabte Bevidsthed gjældende paany. Det skeer ikke for at aabne
Døren for ubegrændsede Phantasier og glemme det virkelige Livs
nærværende Indhold. Tvertimod skeer det for at Menneskets nærværende Liv kan erholde sin tabte Betydning tilbage og blive til et vir%
% --- p. 72 ---
\opage{72}\setcounter{footnote}{0}%
keligt Liv, som det ikke er, uden naar Mennesket erindrer dets
Sammenhæng med Evigheden.

\begin{flushright}(Fortsættes).\end{flushright}

% ===== ARTICLE II: MfL 17 pp. 422--453 (the journal's pp. 73--421 hold other pieces) =====
\bigskip
% --- p. 422 ---
\opage{422}\setcounter{footnote}{0}%

\begin{center}II.\end{center}

\begin{center}{\Large Tanker over Muligheden af Beviser for\\
% PRINTED AS IS: V301 q: crop, layer, tesseract -- a MISPRINT in the Art. II title, as printed
Menneneskets Udødelighed, med Hensyn\\
til den nyeste derhen hørende Literatur.}\end{center}

\begin{center}(Fortsat og sluttet.)\end{center}

\begin{center}\rule{3em}{0.4pt}\end{center}

\begin{center}VIII.\end{center}

I det Foregaaende have vi besvaret Spørgsmaalet, i hvilken Betydning af Ordet det er muligt at føre Beviser for Menneskets
Udødelighed. Herefter skulle vi meddele de mærkeligste Forsøg,
som vor Tids Philosopher have gjort paa videnskabelig at begrunde
deres Overbeviisning om Udødelighedsbegrebets Realitet. Udførligt
Hensyn vil dog kun blive taget til de Forfattere, som ere gaaedeen særegen Vei i Problemets Behandling.

En af de Første eller maaskee den Første, som optraadte
mod F. Richters ovennævnte Skrift, var C. H. Weiße. Hans første
Yttring i denne Anledning var en Recension i Berliner Jahrbücher; men ved endeel derimod stilede Angreb blev han bragt til
at udgive et nyt Skrift\footnote{Die philosophische Geheimlehre von der Unsterblichkeit des menschlichen Individuums von C. H. Weiße. Dresden 1834.}, som her især kommer i Betragtning.
Hine Angreb kom nemlig ikke blot fra F. Richter, men fra orthodoxe Theologer, som havde Recensionens Forfatter mistænkt for,
at han i Grunden var enig med F. Richter, men kun misbilligede
den Handling, at røbe den hegelske Philosophies Mysterier for
det uphilosophiske Publicum. At fralægge sig denne Beskyldning,
var i hiint Skrift Weißes nærmeste Formaal, men derved foranledigedes han til med det Samme at give Læseren nogle aphoristiske
Meddelelser af sin Udødelighedstheorie. --- Weiße hører til de Di%
% --- p. 423 ---
\opage{423}\setcounter{footnote}{0}%
sciple af Hegel, som holde for, at deres Lærers videnskabelige
Form og Methode lader sig bruge til Fremstillingen af den sande
Verdensanskuelse. Derfor betjener han sig i sin Afhandling tildeels af en Terminologie, som i en Reproduction af hans Tankegang ei ganske lader sig undgaae, med mindre man reent vil udslette det Eiendommelige deri. Den efterfølgende Fremstilling af
hans Beviis for Udødeligheden bliver given i directe Tale, skjøndt
jeg ingenlunde i alle Stykker kan vedkjende mig den deri udtalte
Overbeviisning.

„Hvis det menneskelige Individs Udødelighed skal lade sig bevise, maa Beviset føres i den Deel af Philosophien, som udvikler den absolute Aands Begreb. Hverken i Philosophiens abstracte
Region (ɔ: den speculative Logiks Kategorier) eller ved physiologiske
Undersøgelser, der snarest føre til Formodning om Individets Forgjængelighed, kan Udødeligheden finde sin videnskabelige Begrundelse. Men naar den absolute Aand fattes i sin concrete, levende
Skikkelse, naar dens høieste Energie ikke længer sættes i Tænkningen, der tænker sig selv, da kan man haabe fra dette Begreb
at finde Sikkerhed for de endelige Aanders uendelige Vedvaren i
Tiden. Spørgsmaalet bliver da ikke, om Udødelighed tilkommer
den mennelige Individualitet\footnote{Da disse Ord findes hos Forfatteren selv, har jeg naturligviis ikke
bekymret mig om, at give dem større Bestemthed. Den meest afgjorte Modstander af den christelige Udødelighedslære bringes forresten let til at indrømme, at menneskelig Individualitet tilkommer det
Udødelige. --- Det er mærkeligt, at intet af Partierne (saavidt jeg
veed) under Striden om Udødeligheden har benyttet sig af saadanne
Udtryks Tvetydighed til barnagtige dialectiske Drillerier. Uagtet
det er umuligt med almeengyldige Udtryk at betegne det Individuelle,
der menes, naar man taler om Individernes Vedvaren, have dog
de Stridende her bestandig taget hinandens Ord efter Meningen.
F. Richter bruger, for bestemt at betegne det aldeles Individuelle,
den Vending: Jeg F. Richter fra Magdeburg o. s. v.} og Personlighed, men om menne%
% --- p. 424 ---
\opage{424}\setcounter{footnote}{0}%
skelig Indivividualitet og Personlighed tilkommer det Udødelige.
Kun naar man tænker sig den tænkende Substantses Fremtræden under Personers Form som en Fordybelse af Substantsen i sig selv, hvorfra den vender tilbage --- ikke til abstract
Almindelighed, men til en rigere og dybere Eenhed med sig selv,
har man betraadt den Vei, som kan føre til Problemets videnskabelige Opløsning. Beviset for den menneskelige Aands Udødelighed har kun Betydning for den, hos hvem en høiere aandig
Erfaring kan forudsættes. Naar Beviset ikke skal opfattes som
en tom Formalisme, maa det referere sig til æsthetiske, sædelige
og religiøse Kjendsgjerninger, og ved en speculativ Behandling af
dette Indhold kan i Philosophiens System et videnskabeligt Beviis
for Udødeligheden blive ført. Her kan det ikke gives efter sin
hele Udstrækning; det faaer kun sin hele Betydning ved sin dialectiske Plads i den speculative Theologie. Her skal blot gives en
fragmentarisk Betegnelse af Bevisets Gang.

Det uundværlige Grundlæg for det philosophiske Beviis er
den æsthetiske Bevidsthed, det vil sige den Vished, at det, der i
visse Naturphænomener og i Kunstens Værker træder os imøde,
naar vi tillægge dem Skjønhed, er den absolute Aand selv, som
her træder frem i et særligt Phænomen og giver sig selv en aldeles individuel Begrændsning. Denne Overbeviisning hører til
Tidsalderens høiere Dannelse. Ogsaa de, der mangle Bekjendtskab med den speculative Philosophie, udtale jævnlig, fra deres
umiddelbare Standpunct, i forskjellige Udtryk den samme Vished.
Det er i hiin Erfaring blevet dem anskueligt, at det skjønne Phænomen er een af de utallige Skikkelser, hvori det Udødelige, Verdensaanden, som en uudtømmelig Proteus lader sig tilsyne. Det
skjønne Phænomen er den almindelige Typus for den Individuation, som maa fattes, om det Skabtes Udødelighed skal fattes.
Et levende Begreb om den absolute Aands særlige Fremtræden
i Tiden, saaledes som den erfares i hiint Phænomen, maa ligge
% --- p. 425 ---
\opage{425}\setcounter{footnote}{0}%
til Grund for den videnskabelige Tankegang, der fører til Overbeviisning om Menneskets Udødelighed.

Med denne Bevidsthed om det Skjønnes Natur er dog kun
givet det abstracte Almeenbegreb, hvorunder det i Sandhed Udødelige, hvis det skulde findes i de skabte Væsners Række, bliver at
subsumere. Det skjønne Phænomen er saa langt fra selv at
være det Udødelige, at dets Forgjængelighed med Grund er bleven
til et Ordsprog. Derfor var det langt fra, at den æsthetiske Bevidsthed, da den først vaagnede i den moderne Verden, havde den
Virkning, som den dybere Udvikling af dens Indhold maa medføre. Tvertimod førte den nærmest til et ganske modsat Resultat, idet den æsthetiske Verdensbetragtning vænnede Tidsalderen
til i Aandeverdenens individuelle Skikkelser „kun at see den ene,
den evige Grundkrafts forbigaaende Yttring.“ Men med Indsigten i det Skjønnes guddommelige Natur var Overbeviisningen om
guddommelige Kræfters Nærværelse og Virksomhed i Menneskeaanden nødvendig forbunden. Den æsthetiske Sands fører selv uden
tydelig Bevidsthed til en Reflexion, som med det Begrundede tillige fatter Grunden. De productive Kræfter i Mennesket, som
føde det Skjønne, og ved hvilke selv det forhaanden værende
Skjønne under Anskuelsen kommer til sin sande aandige Virkelighed, maatte saaledes erkjendes for Substantsen i det Skjønne
og for det sande Udødelige. Det er da et naturligt Udbrud af
den æsthetiske Bevidsthed, at der i Menneskeaanden er \textbf{noget}
\textbf{Udødeligt}, Noget, som hører til den guddommelige Substantses
evige Liv.

Men hverken den philosophiske Videnskab eller den sædeligreligiøse Drift kan blive staaende ved Forestillingen om noget Udødeligt, som hiin Reflexionsslutning førte til. Videnskaben har
fattet det Begreb om Aanden, at den bestaaer i Tænkning og
Bevidsthed: et aandeligt Noget, som ikke har Formen af Jeghed
og Selvbevidsthed, er for den en reen Uting. Den vil da op%
% --- p. 426 ---
\opage{426}\setcounter{footnote}{0}%
stille det Alternativ, at der enten slet ikke bliver Noget tilbage, af
hvad der hører til Individet som Phænomen i Tiden, eller, at
det med sin Selvbevidsthed vedbliver efter Døden. Skal den
æsthetiske Bevidsthed kunne forsones med den religiøse og den
philosophiske Verdensanskuelse, maa det skee saaledes, at man som
det sande Udødelige kun tænker sig det selvbevidste Subject, der
er deelagtigt i den æsthetiske Substantialitet og Forklarelse.

Her have vi fremhævet de Motiver til en sand philosophisk
Udødelighedstro, som ogsaa ere tilgiængelige for en mindre skarp
dialectisk Betragtning.\footnote{Weiße henviser paa dette Sted til sin Æsthetik, hvori han --- for
at følge hans egen Talebrug --- fremstiller den Dialectik, som fører
fra Phænomenets substants- og væsenløse Objectivitet til den selvbevidste Aands i og for sig værende baade subjective og objective
Substantialitet. Der skal efter hans Mening søges et Led af det
strænge, videnskabelige Beviis.} Alt løber derpaa ud, at man fatter
Begrebet om en absolut aandig, men dog creaturlig Individualitet og Personlighed, der har det samme Forhold til den blot endelige Personlighed, som det skjønne Phænomen til det naturlige.
Det skjønne Phænomen er det Naturliges høieste Sandhed; det
er Aabenbaringen af et Væsen, som er høiere end al Natur.
Netop saaledes er den aandige Substants, som vi erkjende for
den sande, for den udødelige, ikke et ubekjendt, spøgelseagtigt Noget, som svæver over den endelige Aand, men Menneskets egen
personlige Aand, som er gjenfødt i Aand og Sandhed. Det er
fuldkommen rigtigt, naar man blot erkjender Verdens guddommelige Substants for det Udødelige, men derimod erklærer Alt, hvad
der vel er af denne Substants, men ikke denne Substants selv,
for dødeligt og forbigaaende. Men det ved Guds Aand gjenfødte
Menneske er ikke blot af Substantsen, men Substantsen selv,
skjøndt ikke den hele Substants.

% --- p. 427 ---
\opage{427}\setcounter{footnote}{0}%
Den hele foregaaende Tankegang kan maaskee misforstaaes,
som om dermed skulde være meent, at den Gjenfødelse i Aanden,
som deri fremstilles som Betingelse for Menneskets Udødelighed,
kun var at opnaae ved en udtrykkelig æsthetisk Retning. Man
kunde troe der lærtes, at det kun var Beskjæftigelse med Kunsten
eller en uddannet Sands for skjønne Phænomener, der gjorde
Mennesket deelagtigt i det evige Liv. Men det er meget langt
fra, at herved tænkes paa at tilbagekalde det Afguderi med Kunsten, som den alvorlige Videnskab nu saa heldig har fortrængt.
Det har kun været Planen at betegne den Vei, som den videnskabelige Dialectik maa tilbagelægge fra et Standpunct til et andet. Dermed udhæves tillige Hovedmomenterne i den Ideegang,
som Tidsaanden, ikke det enkelte Individ i sin Udvikling gjennemløber. Det udødelige Jeg er ikke det poetiske Genie, som med
absolut Ironie seer ned paa det Forgjængeliges Verden, men det
religiøse Sind, som med Hengivenhed seer op til det Eviges og
Guddommeliges Verden. Med andre Ord: Overgangen fra Tanken om det Eviges Tilsyneladelse til Tanken om det Eviges Substants er en Overgang fra den æsthetiske Region til den sædeligreligiøse. --- Ved den æsthetiske Betragtning forberedes paa flere
Maader Overbeviisningen om Menneskets Udødelighed: blandt Andet vindes der den Erkjendelse, at jo Mere Menneskets Aand har
af det Absolutes og Eviges Substants, desto stærkere Individualitet og Personlighed har den. Det er saa langt fra, at hiin
Substants berøver Menneskets Aand sin Personlighed, at den
tvertimod gjør den i Sandhed personlig. Den vil ogsaa under
alle Betingelser vide at producere den et nyt Legeme, da den ikke
kan være det foruden.

Udviklingen af den æsthetiske Bevidsthed er altsaa en Forberedelse til den videnskabelige Fremstilling af Menneskeaandens Bestemmelse, men sin egentlige Opløsning finder Problemet først
i den speculative Theologie og Ethiken. Ogsaa her ligge høiere
% --- p. 428 ---
\opage{428}\setcounter{footnote}{0}%
Erfaringer til Grund for den dialectiske Tankegang, og uden
disse bliver den sande Forstaaelse af den philosophiske Fremstilling
aldeles umulig. En umiddelbar Bevidsthed maa nemlig give Individet Vished om, at det med sit sædelige Jeg lever i den udødelige Substants. Skal Individet være udødeligt, maa det som
gjenfødt i Aanden være af Substantsen og i Substantsen. Denne
Gjenfødelse --- thi saaledes er man saavel efter theologisk som
philosophisk Talebrug berettiget til at benævne den Kjendsgjerning,
her sigtes til --- skaffer først Mennesket Deelagtighed i det evige
Liv, hvilken det ikke har af Naturen. Den høiere Tilstand, som
Mennesket kommer i, naar det saaledes er optaget i Substantsen,
er Gjenstand for dets egen indre Anskuelse og giver Subjectet en
fast Tro paa dets Udødelighed, som er Grundlaget for den videnskabelige Tydeliggjørelse. Hvo der mangler en saadan Anskuelse,
kan kun bringe det til en overfladisk, bogstavelig Forstaaelse af
Philosophiens Beviser, men deres sande Betydning maa forblive
ham en Hemmelighed. Paa den anden Side kan Videnskaben
dog ikke uden Beviis forudsætte Gyldigheden af de her antydede
Erfaringer: den systematiske Fremstilling af Philosophien fører til
at fatte disse Kjendsgjerningers Begreber i deres dialectiske Sammenhæng.“

Saaledes betegner Weiße kortelig den Vei, som han mener
det videnskabelige Beviis for Menneskets Udødelighed maa gjennemløbe. For rigtig at sætte sig ind i Forfatterens Tankegang behøvede man vistnok en udførligere Forklaring af flere Momenter
deri, men i Skriftet selv er den fragmentariske Fremstilling af Beviset affærdiget paa nogle faa Blade, saa at jeg tvivler meget
paa, nogen Læser. der vil finde den nærmere Beskeed, som her
savnes om adskillige Punkter.

En Sætning, som Forfatteren \textit{udbreder} sig udførlig over,
har jeg blot løselig berørt: det er hans besynderlige Mening, at
Udødeligheden kun tilkommer en udvalgt Deel af Menneskeheden,
% --- p. 429 ---
\opage{429}\setcounter{footnote}{0}%
de udmærkede Individer, i hvis Lod den aandige Gjenfødelse,
Optagelsen i Substantsen, er falden. Han paastaaer endog, at
denne Lære har fuldkommen Overeensstemmelse med en rigtig fremstillet christelig Dogmatik. Christus skal --- saa mener han ---
intetsteds have lært, at Udødelighed var det naturlige Menneskes
Arvelod, eller var en nødvendig Følge af den menneskelige Naturs
Indretning. Udødeligheden eller det evige Liv skulde først tilfalde
Mennesket, som blev renset og forædlet ved Christi Lære. Udførelsen af denne Sætning, som hverken blandt Theologer eller Philosopher er splinterny, er jeg saa meget heller gledet over, som
Forfatteren i en senere Afhandling tildeels har taget denne sin
Paastand tilbage. Han er nemlig bleven overbeviist om sin exegetiske Vildfarelse ved en særdeles grundig Afhandling af Theologen Julius Müller, (hvilken findes indrykket i Umbreits theologische Studien und Kritiken).

I den nys fuldførte Fremstilling af Weißes Tanker findes,
som alt bemærket, et og andet Udtryk, som kun er forstaaeligt
for den, der har lidt Bekjendtskab med det hegelske Skolesprog.
Desuden er Forfatterens Tankegang noget ubestemt og fragmentarisk; ja det hele Skrift har overhovedet Characteren af et med
stor Ilfærdighed udført Leilighedsarbeide. Af begge disse Grunde
anseer jeg det for hensigtsmæssigt, med et Par Ord at udpege de
Punkter, som udgjøre Bevisets Kjærne, hvis der skulde findes nogen Kjærne deri. --- Det Særegne ved Beviset er dets Tilknytning til den æsthetiske Bevidsthed. Tidsalderen har (mener Forfatteren) en usædvanlig udviklet Sands for det Skjønnes Betydning: dens almindelige Cultur bestaaer for en stor Deel i den
klare Bevidsthed om, at Verdensaanden giver sig selv en Form i
det skjønne Kunstproduct. Dette har sin egentlige Virkelighed i
Menneskeaanden, der skuende producerer det. Ved Bevidstheden
om denne productive Kraft i sin Aand bringes Mennesket til Erkjendelsen af, at der lever noget Udødeligt i ham. I disse ube%
% --- p. 430 ---
\opage{430}\setcounter{footnote}{0}%
stemte Kategorier: „en guddommelig Kraft“, „noget Udødeligt“ o. s. v.
har Udødelighedstroen nuomstunder sin sædvanlige Form. Men
Videnskaben maa tydeliggjøre sig Betydningen af dette Noget og
kan, efter Forfatterens Paastand, kun finde Forklaringen i Begrebet selvbevidst Aand. Naar nu det gjenfødte Menneskes religiøse og sædelige Bevidsthed træder til hiint æsthetiske Klarsyn,
vinder Mennesket en umiddelbar Bevidsthed om, at det udødelige
Væsen er dets eget i den guddommelige Substants bevarede Selv.

Men hvad Værd maae vi tillægge det hele Beviis? ---
Det tager tilsidst sin Tilflugt til en en mystisk Anskuelse, hvis
Gyldighed skal forsvares i et tilkommende videnskabeligt Beviis.
For saavidt har Philophien ikke gjort noget Fremskridt derved, og
det skulde neppe have tildraget sig saa stor Opmærksomhed, hvis
Forfatteren ei i Forveien havde havt sit berømte Navn. Det var
et nyt Phænomen, at en Mand, der var saa fuldkommen deelagtig i Tidsalderens speculative Cultur, mod Skik og Brug talte
saa udtrykkeligt og uforbeholdent om Individets Udødelighed. Men,
for at jeg skal meddele min subjective Mening, røber den omtalte
Afhandling, at Forfatterens Udødelighedstro var temmelig nys
bleven vakt. Under denne Forudsætning kan man bedst forklare
sig en skeptisk Usikkerhed og Raadvildhed, der gaaer igjennem det
hele Skrift og stærkt contrasterer imod den faste Overbevisning,
som Forfatteren i snart derpaa følgende Skrifter har lagt for
Dagen. Hans unge Tro synes i kort Tid at have vundet meget
i Styrke, ligesom et Barn, der nys er kommet til Verden, i
sin første Livsperiode voxer meest.

Forøvrigt er det i Øine faldende, at den Vending, hvorved
den æsthetiske Bevidsthed gjøres til Udgangspunktet for Beviset,
er temmelig vilkaarlig. Herved menes ei, den er urigtig; thi
hvor Erkjendelsen virkelig hænger organisk sammen, maa der nødvendigviis kunne gjøres Overgang fra ethvert Moment deri til
ethvert andet, og denne Overgang maa tillige medføre Overbe%
% --- p. 431 ---
\opage{431}\setcounter{footnote}{0}%
viisning for Enhver, der ved Philosophens Ord kan vækkes til at
give hans Tanker den Integration, som er meent derved. Men,
naar Weiße mener at have betegnet den eneste rigtige, almeengyldige Tankegang, der maa følges i Beviisførelse for Udødeligheden,
naar han fremdeles holder for, at nogle fragmentariske Vink til
denne Tankegang ere af synderlig videnskabelig Vigtighed, da tager han sikkert meget feil.

Til Slutning blot een Bemærkning om Weißes Udødelighedslære.
Af den Yttring i hans Afhandling, som ovenfor er uddraget, at den absolute Aand nok vil vide at producere Menneskets Aand det nye Legeme,
uden hvilket den ei kan existere, synes man at kunne drage den Slutning,
at han til den Tid ei klart har tænkt sig Udødeligheden som en uafbrudt
Sammenhæng af Menneskets Selvbevidsthed. Han angiver i det
Mindste ikke her nogen Udvei for den, der med Nutidens Begreber om Sjælen vil tænke sig Menneskets Bevidsthed overført
fra den ene Corporisation til den anden. Desuagtet har han i
sine senere Skrifter forudsat som en afgjort Sag, at Udødeligheden bestod i en continuerlig Selvbevidsthed og fortsat Erindring.

\begin{center}IX.\end{center}

Kort efterat Weiße havde udgivet sin Afhandling, udkom et
Skrift af den yngre Fichte\footnote{Die Idee der Persönlichkeit oder der individuellen Fortdauer von
I. H. Fichte. Elberfeld 1834.}, hvilket for en stor Deel gaaer ud
paa at give videnskabeligt Beviis for Udødelighedstroens Realitet.
Først gjør han kortelig rede for sit metaphysiske Standpunct og
sit Forhold til endeel samtidige Philosopher, især til Hegel, Göschel og Günther. Efter denne polemiske Begyndelse, der næppe
tilfulde kan forstaaes af nogen Læser, der ei kjender Forfatterens
øvrige Værker, meddeler han nogle fragmentariske Sætninger om
Guds Forhold til Skabningen, i hvilke man under andre Vendinger og nye Tankeforbindelser møder endeel af de sædvanlige
% --- p. 432 ---
\opage{432}\setcounter{footnote}{0}%
Beviser for Udødeligheden. I dette Partie af Afhandlingen findes
en stor Tankefylde og en saa riig Veltalenhed, at dets Hovedinteresse nødvendig maa gaae tabt ved et sammentrængt Udtog.

„Ethvert Naturindivid --- saaledes omtrent kan Fichtes
Tankegang udtales --- har et oprindeligt og eiendommeligt Anlæg,
og dets hele Tilværelse bestaaer kun i, at dets Anlæg efterhaanden kommer til Virkelighed. Men ingen Natureenhed ophører
at være til, før dens Evne saaledes har udtømt sig i Yttring,
at ingen Rest bliver tilbage. Det, som vi her have kaldt Individernes Anlæg, og som efter forskjellige Hensyn ogsaa undertiden
kaldes deres Forbilleder eller deres Ideer, er det guddommelige
Grundlæg for deres Tilværelse, der ikke kan tilintetgjøres ved nogen udvortes Magt. Skal en tilværende Eenhed forsvinde af
den virkelige Tilværelse, da maa det være, fordi Kilden til dens
Liv er udtømt, ikke fordi nogen udvortes Vold har tilintetgjort
den. Det er tænkeligt, at en eller anden Art\footnote{Man seer let, at Forf., naar Talen er om de lavere Trin af det organiske Liv, bruger Udtrykket Art i samme Betydning som Naturindivid.} af organiske Væsner ophører med sin Tilværelse, fordi den har fuldendt den Bane,
der er betegnet ved dens Anlæg, men før dens Livskilde saaledes
er udtømt, ophører den ei at realisere sig. Denne aprioriske Betragtning og en gjennemgribende Naturanalogie berettiger os til
den Antagelse, at Menneskeindividet heller ikke kan gjøres til
Intet, med mindre det til en bestemt Tid virkelig har fuldkommen
realiseret sit Anlæg: den os bekjendte Død kan da ikke være dets
Tilværelses Ophør, med mindre det til den Tid kan have udtømt
sit Anlæg, sin guddommelige Livskilde.

Agte vi nu paa Menneskelivets særegne Form, da finde vi,
at det selvbevidste Liv idelig bryder frem af en dunkel Grund, der i
Sammenligning med den Deel, som efterhaanden kommer til Reflexionens Klarhed, viser sig som et uudtømmeligt Indhold. Hiin
bundløse Afgrund, hvoraf det reflecterede Liv strømmer frem,
% --- p. 433 ---
\opage{433}\setcounter{footnote}{0}%
er det man ved en Misforstaaelse undertiden har tillagt større
Værdi end det Livsindhold, der bringes til fuldkommen klar Selvbevidsthed. Denne Vildfarelse har forlængst fundet sin tilstrækkelige
Gjendrivelse. Men det er dog en uomtvistelig Kjendsgjerning,
at Menneskelivet har hine to Sider; det har en Nat og en Dag,
og disse to Sider kunne hos det virkelige Individ ei adskilles ved
skarpe Grændser. Man er i den nyeste Tid saa langt fra den
Mysticisme, der lægger for stor Vægt paa Menneskelivets natlige
Side, hvoraf det selvbevidste Liv idelig beriges, at man nu med
en modsat Eensidighed kun agter det Liv for menneskeligt, der
alt er bragt til Reflexionens fuldkomne Klarhed. Men hvad man
end dømmer herom, lærer en grundig Naturbetragtning af Menneskelivet, at hiin dunkle Baggrund, der som Indfald, Opfindelseskraft, Genialitet o. s. v. idelig trænger ind i de selvbevidste Livsyttringer, umulig kan udtømmes i Menneskets jordiske
Liv. Derfor kan det heller ikke antages, at Menneskets Forsvinden fra Sandseverdenen er Slutningen af dets virkelige Tilværelse som Individ, da det til den Tid ei kan have udtømt
sit Anlæg eller reent udlevet sit Livs Indhold.

Denne Tankegang fører ogsaa naturlig til den Sætning,
at en følgende Livsudvikling i den hele Natur idelig findes præformeret i en foregaaende Livstilstand. Dette Phænomen viser
sig i Menneskets Tilværelse under Form af Anelser, Forudfølelser,
hvorved vi vises hen til en Mangfoldighed af Kjendsgjerninger,
som endnu langt fra har fundet sin fuldstændige Forklaring i den
videnskabelige Psychologie. En saadan Forudfølelse ligger til
Grund for den hos alle Folkeslag gjenkommende Tro, at der
gives et Liv efter Døden, et Land for Sjæle. Denne Tro maa
betragtes som aldeles ubegribelig, hvis man ei her erkjender en
apriorisk Forudsætning, som i Følge den menneskelige Fornufts
Natur nødvendig udvikler sig som det endelige Livs Forklaring.
Ligesom Grundtanken om det Absolute trænger sig selv frem, hvor
% --- p. 434 ---
\opage{434}\setcounter{footnote}{0}%
kun det Tilældige bliver umiddelbar givet, saaledes bryder ogsaa det
aprioriske Begreb om Uforgjængelighed frem som et nødvendigt Supplement, hvor kun Forgjængelighed frembydes af Erfaringen. Hvorledes forklares dette idelig gjenkommende Begreb om Individets
Uforgjængelighed, hvis det ei antages at være præformeret i
Menneskeaanden?

Troen paa Udødelighed erholder dobbelt Styrke ved den sædelig-religiøse Bevidsthed. Thi den, der veed, at hans Villie falder sammen med den guddommelige Villie, til hvis Organ han
hører, har umiddelbar Vished om, at han er optaget i det udødelige Væsen. Udviklingen af denne Bevidsthed fører til et fuldkommen gyldigt ethisk-religiøst Beviis for Sjælens Udødelighed.

Men uagtet alle disse Vendinger føre til Overbeviisning om
Udødelighedsbegrebets Realitet, bliver endnu en væsentlig Fordring
tilbage, som nødvendig maa tilfredsstilles, før det fuldstændige
Beviis for Udødeligheden kan siges at være tilveiebragt. Vi have
seet, at en omfattende Naturanalogie saavel som aprioriske Betragtninger vise hen paa Menneskets Udødelighed; men nu møder
Videnskaben tillige den Opgave at godtgjøre Muligheden af Udødeligheden. Der kræves et Beviis for, at det i den givne Natursammenhæng er tænkeligt, at Menneskets Liv kan fortsættes efter
Døden. Dette physiologiske Beviis for Udødeligheden kan saaledes
udvikles og styrkes ved ægte Naturbetragtning, at det efter sin
Natur aldrig kan blive fuldkommen afsluttet. Til et saadant Beviis gives her et lidet Bidrag.

Uagtet da Udviklingen af Bevidsthedens Forudsætninger, i
Særdeleshed Tydeliggjørelsen af de sædelige og religiøse Begreber,
medføre Overbeviisning om Udødeligheden, er det dog først i Naturvidenskabernes Kreds, at det videnskabelige Beviis maa finde
sin egentlige Fuldførelse. Det maa her vises, hvorledes Mennesket kan være udødeligt, hvorledes det er physisk muligt, at Mennesket efter Døden fortsætter sit Liv. Især venter det Spørgs%
% --- p. 435 ---
\opage{435}\setcounter{footnote}{0}%
maal sin Besvarelse, hvorledes man skal tænke sig, at Sjælen
fortsætter eller gjenvinder sin Corporisation, efterat Legemet ved
Døden bliver opløst i sine chemiske Bestanddele. Her tilveiebringer den grundige Naturbetragtning et Beviis efter Analogie, som
berettiger til den Slutning, at Døden betegner en ny Station
af Menneskets organiske Livsudvikling. I et saadant Beviis er
Ideen i sin Totalitet det Lys, der leder fra de ved Erfaring givne
Stykker af dens Virkelighed til en sikker Forudsætning af de
Partier, der ei kunne blive Gjenstand for nogen Jagttagelse. En
saadan Beviisførelse kan vel aldrig ganske tabe sin hypothetiske
Characteer, men der gives ikke nogen anden sikker Vei, som fører til Erkjendelsen af den concrete Virkelighed, der er utilgjængelig for Erfaringen. Et saadant Beviis vil efter sin Natur bestandig vinde meer Evidents, eftersom Physiologien (en Videnskab, som endnu kun er i sin Begyndelse) antager en fuldkomnere
Skikkelse. Til Grund for et saadant Beviis maa der først og
fornemmelig lægges rigtigere Begreber om Legem og Sjæl, end
de, der nu ere meest i Gang i Physiologie og Psychologie. Først
derved bliver man i Stand til at forklare den sande Betydning
saavel af Avling som af Død, der her maa gjøres gjældende.

Vi laane her den omtalte Sætning af Ontologien, at der til
Grund for enhver individuel Tilværelse ligger et guddommeligt
Anlæg, som successivt realiseres deri. Dette give vi her Navn
af det organiske Forbillede: efter dette Forbillede kommer enhver
menneskelig Individualitet til Virkelighed i Tid og Rum. Det
er den overvirkelige Magt, det sande Livsprincip, efter hvilket
Virkeligheden her aldeles retter sig. Hverken Sjæl eller Legem
kan her sættes som Begyndelsespunktet, men begge maae i deres
Eenhed tænkes som liggende i Forbilledet, og under dettes Virkeliggjørelse træde begge frem som fuldkommen samtidige. Det lader sig ei paastaae, at Sjælen skaber Legemet, da derved forudsættes en Prioritet af Sjælen i Tiden, som ei er tænkelig; men
% --- p. 436 ---
\opage{436}\setcounter{footnote}{0}%
saavel Sjæl som Legeme, den udeelbare Personlighed, kommer i
sin Totalitet til Virkelighed ved Individers Forening til Generationen. I Conceptionens Moment er det nye Individ dog noget
aldeles Nyt, der fra Intet kommer til Virkelighed, idet her det
overvirkelige Forbillede, den guddommelige Idee, begynder sin virkelige Tilværelse.

Den Sætning, som her især er af Vigtighed, er den, som
af Nutidens Physiologie er bragt til saa høi Grad af Klarhed,
den, at Menneskets Legem ei er den palpable Skikkelse, som i
Almindelighed faaer dette Navn. Dette Legem vexler bestandig
sine materielle Partikler og er kun Tilsyneladelsen af det indre
Legem, den organiske Identitet, som holder sig under den uafbrudte Vexel. Legemet er den bestemte organiske Proces, hvorved
det individuelle Forbilled under Reproduction og Assimilation udpræges i Omverden. Hvis der da i Individualitetens Forbilled Intet indeholdtes, uden denne organiske Retning, kunde Menneskets
Liv være sluttet med Legemets Opløsning ved Døden. Men
forsaavidt som Forbilledet indeholder en anden Side, det selvbevidste personlige Liv, som ei kan være kommet til fuldendt Virkelighed, naar Døden indtræder, idet Individualitetens bevidstløse
Indhold til den Tid ikke heelt er optaget i Selvbevidsthedens
Kreds: er det guddommelige Anlæg her ei udlevet, og Menneskets
virkelige Tilværelse, (hvori det organiske Liv er et Moment) kan
ei være sluttet med Døden. Erkjendelsen heraf er det, som deels
i en umiddelbar Form anteciperes hos det naturlige Menneske,
som venter Udødelighed, uden at vide hvorfor, deels under den
sædelige og religiøse Bevidsthed bliver til fast Overbeviisning, idet
Mennesket paa Grund af hiin høiere Erfaring veed sit eget Liv
optaget i Eenhed med det guddommelige Liv.

Men da nu det høiere personlige Liv kræver det organiske
Liv, som sin ufravigelige Betingelse, bliver her det Spørgsmaal
forelagt Naturvidenskaben, hvorledes et nyt organisk Liv føier sig
% --- p. 437 ---
\opage{437}\setcounter{footnote}{0}%
til det af os bekjendte, der ophører med Døden: i hvad Forhold
Menneskets Tilstand efter Døden staaer til den, der gik forud.
Her vises for det Første hen til den Lov, som gjælder for enhver
Art af organisk Liv, at det udvikler sig under en afvexlende Contraction og Expansion. Exempler paa saadanne Contractioner haves i den sædvanlige periodiske Søvn, visse Dyrs Vinter- og
Sommerdvale, Insecters Puppetilstand o. s. v. Paa disse Contractioner vente vi kun af den Grund en ny Expansion, at Erfaringen har lært os, at en saadan følger paa. En saadan Contraction er Døden, som er Slutningen af een Periode af det
organiske Liv, der falder sammen med Begyndelsen af en ny.

Her finder den Sætning, at Corporisationen er en uafbrudt
Vexel af udvortes Livsbetingelser, en vigtig Anvendelse. Det sande
Legem, vide vi, er den organiske Proces, hvorved Individualiteten underkaster sig udvortes materielle Midler, og deri, at en
Sum af materielle Midler afløses, ligger ingen Grund for Antagelsen af Livets Ophør. Livet kunde kun ophøre, hvis det guddommelige Forbillede var kommet til fuld Virkelighed; thi da
først var Levekraften udtømt. Kun de for os synlige Livsmedier
ere traadte ud af det organiske Livs Tjeneste, men Individualiteten
hævder sig i en anden Corporisation, for hvilken Mennesket i den
Organisation, som er Gjenstand for vor Erfaring, ei har nogen
Receptivitet. Det organiske Liv er en selv ved Døden ei afbrudt
Proces: det Spørgsmaal, hvilken Deel af Mennesket der gaaer
tabt i Døden, har her ingen Betydning; thi der gaaer Intet
tabt; Individualiteten bevares heel og holden giennem Dødens Crisis; selve Corporisationen fortsættes i et andet Medium.

Ligesom man tit i Naturen seer organiske Væsner at undergaae flere Forvandlinger i deres Corporisation, saaledes er Døden
en organisk Metamorphose, hvorved en ny Periode begynder i
Menneskets Livsudvikling. Men der kan ei være Talen om en
immateriel Mellemtilstand, hvori den abstracte Sjæl befandt sig
% --- p. 438 ---
\opage{438}\setcounter{footnote}{0}%
i Overgangen fra een Corporisation til en anden. Den organiske
Proces, uden hvilken Livet var utænkeligt, er her aldeles continuerlig, idet Individualiteten paa eengang slipper et Livsmedium
og underkaster sig et nyt. Der gives ingen fornuftig Grund for
den udbredte Mening, at Døden er en ubegribelig Forflyttelse til
et Sted, der er absolut adskilt fra den nærværende Skueplads for
vort Liv, og slet ikke staaer i noget rumligt Forhold dertil. Denne
Mening strider aldeles mod det rigtige Begreb om Naturens
Eenhed og Sammenhæng. De Livsmedier, hvorved Menneskets
personlige Tilværelse fortsættes, kunne ei søges i fjerne Regioner,
men de omgive og gjennemtrænge Mennesket under dets nærværende
organiske Tilstand. De nye Livsbetingelser maae nødvendig tænkes som en finere Materie, der ei afficerer de Dødeliges Organer.
Det ligger netop i Sagens Natur, at Mennesket i sin nuværende Tilstand ei har Receptivitet for det Medium der hører til
den nye Livsform.

Ligesom Menneskets organiske Liv fortsættes gjennem Dødens Catastrophe, saaledes gaaer ogsaa den hele Grundretning af
det personlige Liv med over i dets nye Tilværelse, saa at dets
Tilstand, dets indre Harmonie eller Disharmonie, bliver et umiddelbart Resultat af dets jordiske Liv.

Til disse Slutninger berettiger os det physiologiske Beviis,
som grundes blot paa en analogisk Naturbetragtning: det fører,
strængt for sig taget, kun til den Sætning, at Døden ei er Livets
Ophør, men forslaaer ei til Antagelsen af egentlig Udødelighed.
I den hidtil udførte Tankegang ligger ingen Grund til at benægte den Mulighed, at Menneskets Liv engang kunde sluttes
med fuldkommen Tilintetgjørelse; hvilket var tænkeligt, hvis Mennesket engang fik udtømt sit Livs Kilde,--- det vil sige, hvis det
fik realiseret Alt, hvad der ligger i dets Anlæg og som her ei
kommer til Virkelighed. Fuldstændigt Beviis for Menneskets Udødelighed vindes derfor ei ved disse Naturbetragtninger alene, men
% --- p. 439 ---
\opage{439}\setcounter{footnote}{0}%
disse maae forbindes med en rigtig Udvikling af den sædelig-religiøse Bevidsthed, som har sin Plads i en anden Region af Philosophiens System. Her var det kun Planen at besvare det vigtige
Spørgsmaal, som nuomstunder sædvanlig forbigaaes, hvorledes
Fortsættelsen af Menneskets Liv efter Døden i den bekjendte Natursammenhæng var tænkelig. Idet vi blot have henvendt Opmærksomheden paa dette Punct, have vi her ikke berørt andre
Problemer, som høre til en fuldstændig Udødelighedslære, f. Ex.
Betydningen af Læren om Legemets Opstandelse. Saadanne Problemers videnskabelige Opløsning maa finde Sted i Religionsphilosophien.“

Her har jeg givet et frit Referat af det omhandlede Skrifts
Hovedindhold. Den, der læser Afhandlingen selv, finder deri
en Mangfoldighed af aandrige Tanker over forskjellige Dele saavel af Physiologien som Psychologien, hvilke her efter Planen
maae forbigaaes. Men selv i Gjengivelsen af de Tanker i
Afhandlingen, som angaae Hovedspørgsmaalet, tilstaaer jeg, at
jeg med Forsæt har viist nogen Tilbageholdenhed. De Læsere,
som ved at gjøre sig bekjendte med det Udtog, som her er meddeelt,
ville fælde den Dom, at deri allerede findes Noget af vilkaarligt
Phantasispil, kunde let ved Læsning af Afhandlingen selv finde flere
Incitamenter for deres Spottelyst, som jeg for min Deel ei vil lette
dem Adgangen til. --- Da jeg nemlig etsteds ovenfor skrev, at
Adskillige ved Tanken om, at et virkeligt Liv maatte være bundet
til Tid og Sted, lade sig forlede til at give altfor nøie Besked
om de Dødes Rige, tænkte jeg blandt Andet paa det her omhandlede Skrift af Fichte. Thi jeg kan ikke nægte, han forekommer mig selv at give saa bestemt Besked, at Samtiden, der
selv i Hovedsagen er temmelig sceptisk, sikkert med Uvillie vender
sig bort fra hans dristige Gisninger.

% --- p. 440 ---
\opage{440}\setcounter{footnote}{0}%

\begin{center}X.\end{center}

Et andet Skrift over Udødeligheden, der af flere Grunde
har tilbraget sig stor Opmærksomhed, er udgivet af C. F. Göschel\footnote{Von den Beweisen für die Unsterblichkeit der menschlichen Seele im
Lichte der speculativen Philosophie, von Carl Friederich Göschel.
Berlin 1835.}. Det har især gjort Opsigt, at en saa erklæret Tilhænger
af Hegels Philosophie som Forfatteren, her optræder som en ivrig
Forsvarer af den christelige Udødelighedslære. Men uagtet denne
og andre væsentlige Afvigelser, vedbliver dog Göschel bestandig
at udgive sig selv for en Philosoph af den hegelske Skole og vil
ikke give den Paastand Medhold, at han maa betragtes som en
Frafalden.

Han angiver det selv som sin Afhandlings Formaal først at
fremstille de traditionelle Beviser for Sjælens Udødelighed, saaledes
som de forefindes, og siden at udvikle dem paany af den speculative
Philosophies Principer og efter dens Methode. Idet jeg kortelig
gjennemgaaer hans Tankegang, vil jeg for Tydeligheds Skyld
vel betjene mig af nogle Omskrivninger, men dog beholde de meest
characteristiske af hans Talemaader. Det sidste bliver saa meget
meer nødvendigt, som Afhandlingen netop er et Forsøg paa at
bringe de bekjendte Argumentationer i Forbindelse med Hegels
Philosophie og følgelig at oversætte dem i hegelske Former.

„Man har for længe siden bragt det til tydelig Bevidsthed,
at der kun kan gives tre Hovedbeviser for Guds Tilværelse. Paa
en lignende Maade kan man orientere sig i den store Mangfoldighed af Beviser for Udødeligheden, som overleveres i Videnskabens
Historie, ved at agte paa den tilsvarende Triplicitet, som udtaler
sig deri.

Det første theologiske Beviis (det cosmologiske) slutter fra
Verden som umiddelbar og tilfældig til dens immanente Grund,
% --- p. 441 ---
\opage{441}\setcounter{footnote}{0}%
Guddommen. Det teleologiske slutter fra Verdens Hensigtsmæssighed til et fornuftigt Væsen, som bestemmer den efter Formaal.
Det ontologiske Beviis endelig slutter fra det allerfuldkomneste
Væsens Begreb til dette Begrebs Realitet.

I analogisk Sammenhæng med disse Beviser staaer en tilsvarende Kreds af Beviser for Sjælens Udødelighed. Først tages
da Sjælen, som den fremstiller sig, tilfældig eller umiddelbar.
Den søger da sin Begrundelse og det ikke blot udad til for at
komme til Gud, men ogsaa indad for at komme til sig selv.
Den viser sig umiddelbart som Bevidsthed eller som det Værende,
der reflecteres i sig selv, som en i sig selv afsluttet Eenhed, der
er sikkret for al Overgang i Andet og følgelig uforgjængelig.
Sjælen er enkelt; thi dens Væsen er Tænkning, og Tænkningen
bestaaer i det Mangfoldiges Eenhed. Sjælen er enkelt, immateriel, den er det sig selv Bevægende o. s. v.: i saadanne Vendinger har man fremstillet Forsætningen og deraf draget den Slutning, at den ingen Forandring er underkastet, at den ikke kan
gaae over i Andet. --- I den populære Philosophie har man i
Stedet for Enkelthed oftere brugt Udtrykket Udeelbarhed og saaledes betragtet Sjælen som en udvortes Gjenstand, uagtet denne
Betragtningsmaade netop er ophævet ved Prædicatet Enkelthed.
Aristoteles udtrykker bestemtere det, som her menes, naar han
kalder Sjælen enkelt, fordi den bevæger sig i sig selv. Dette
metaphysiske Beviis for Udødeligheden svarer til den cosmologiske
Tankegang i Theologien; thi det slutter fra Sjælens umiddelbare
Tilværelse til dens Udødelighed, ligesom det nævnte theologiske Beviis slutter fra den tilfældige Tilværelse til den evige Væren.

I nær Forbindelse med det teleologiske Beviis for den theologiske Idees Realitet staaer i den anden Tankekreds det ligeledes
teleologiske Beviis for Sjælens Uforgjængelighed. Ligesom man
nemlig ved Tanken om Tilværelsens Hensigtsmæssighed vises hen
til det absolute Princip, der bestemmer den efter sine Formaal,
% --- p. 442 ---
\opage{442}\setcounter{footnote}{0}%
vises man ogsaa paa den anden Side derved hen til Tanken om
Formaalene, de menneskelige Sjæle som Verdens Formaal.  Sjælen er Selvhensigt; den er som Guds egentlige Formaal i Skabningen, som Guds Billede, virkelig, uendelig.  Formaalet at
være selv er uendeligt: da dette ligger i Sjælen, er den skikket
dertil: Evnen borger for Realiteten, fordi Gud ellers ikke vilde naae
sit Diemed.  Til denne Tankegang høre de forskjellige moralske
og religiøse Beviser for Udødelighed, i hvilke man erkjender, at
Sjælen efter sin intensive Fylde er bestemt til en Uendelighed af
Formaal og følgelig maa tænkes som uendelig vedvarende.

I de to sidste Beviser, baade det, der holdt sig til Sjælens
Natur som enkelt eller tænkende, og det, der betragter den som
Verdens Formaal, gik man ud fra Sjælens Tilværelse.  Men ligesom i den theologiske Sphære har man ogsaa her en ontologisk
Beviisførelse, der fra Begrebet Uforgjængelighed slutter til Begrebets Realitet.  Da Begrebet Uforgjængelighed, Vedvaren, findes
i Sjælen, maa ogsaa real Uforgjængelighed tilkomme den.  Sjælen kan slet ikke tænke sig absolut Forgjængelighed: den tænker
Vedvaren og kan ei andet end tænke den, følgelig maa ogsaa
dette nødvendige Begreb have Realitet.  Ved Begrebet Udødelighed,
Uforgjængelighed, er det saa naturligt at tænke paa Selvbevidsthedens Udødelighed, at Ordet bruges absolut, uden nærmere forklarende Tillæg.  Ved at forfølge den nødvendige Tænkning af Vedvaren gjennem sine forskjellige Udviklingstrin, finder man den under Form af Følelse, Forestilling og endelig som Begreb, og det
som et Begreb, hvis Modsatte er utænkeligt.

Saaledes have vi fundet de tre dogmatiske Beviser for Udødeligheden, der løbe paarallel med de tre bekjendte theologiske Beviser: det bliver da Philosophiens nærmere Opgave at udvikle % NOTE: a smear on the paper at the mark; the colon from the reader and the sense
dem saaledes, som de gjentage sig i Speculationens Lys.  Den
her udhævede Triplicitet maa lægges til Grund for en tilkommende
philosophisk Historie af Udødelighedsbegrebet, thi vi have her fun%
% --- p. 443 ---
\opage{443}\setcounter{footnote}{0}%
det de tre Hovedstammer, hvoraf den hele Mangfoldighed af Beviser har sit Udspring.  Først gaaer man ud fra Selvbevidsthedens
Natur, dernæst fra Bevidstheden om Gud; og for det Tredie
forenes begge Begreber i den selvbevidste Bevidsthed om Gud.
Det endelige Væsen tænker sig Gud, og dette Forhold mellem
Tænkningen og Gud har evig Realitet.

Naar vi forfølge Sjælens Begreb, dets immanente Selvudvikling fra dets umiddelbare Skikkelse til dets fulde Sandhed, da
føres vi til dens Udødelighed.  Sjælen er først umiddelbar Eenhed; i dennes Diremtion kommer Bevidstheden frem, Adskillelsen
af Tænkning og Væren, og endelig forsones Modsætningen i
Aanden, som saaledes har bevæget sig frem til sin egen Sandhed.
1) Sjæl; 2) Bevidsthed; 3) Aand ere de forskjellige Stationer,
som Aandens Begreb gjennemløber.

Det er overalt den samme Aand, der skjult rører sig først i
Sjælens Tilværelse, dernæst i dens Bevidsthed, indtil den omsider
viser sig i sin Sandhed. --- Først træder den frem som Sjæl i
og med Legemet, følgelig som Individ: her have vi den anthropologiske Sphære.  I den anden Sphære, den phænomenologiske,
viser den sig som Subject, følgelig som Bevidsthed.  Det naturlige Individ er ligegyldigt mod sin Andet-Væren og vedligeholder kun Arten; men det derfra forskjellige Subject er selv, og
er kun, for saavidt det ei bliver Andet.  I den tredie Sphære,
den psychologiske, træder Aanden frem som Aand, idet den har overvundet sin Gjenstand, forsonet sig med den og optaget den i sig.
Den har fra sin første Periode medtaget sin Individualitet, fra
den anden sin Selvbevidsthed: den er som Aand Alt i Alt og er
bleven sig den Rigdom bevidst, som fra Begyndelsen af har ledsaget den.

Aandens Begyndelse var umiddelbar: den var kommen til
sig selv og til Andet uden at vide, hvorledes den var kommen
dertil og hvorfra den selv var kommen.  Den har i sig selv som
% --- p. 444 ---
\opage{444}\setcounter{footnote}{0}%
Individ fundet en Begyndelse, men denne tilfældige Begyndelse
viser hen til en Begyndelse udenfor den selv, som veed af den at
sige; thi da Begyndelsen førte til Tænkning, maa den ogsaa
være kommen af Tænkning; Bevidstheden kan ei være kommen
af det Bevidstløse.  Individet forefinder sin Begyndelse i sig selv
og har følgelig ikke selv gjort den; men denne Begyndelse viser
hen til en høiere Begyndelse. --- Den Slutning, hvorved Aanden fuldendte sig, var Forsoningen med Andet, dens Forening med alt
Andet, som derefter hører sammen.  I dette Andet maa den første
Begyndelse ligge, som er Maalet for vor Søgen.  Naar Aanden
virkelig bestaaer i Forening og virkelig har opnaaet denne Forening
maa den deels erkjende sig selv i den nærmere Bestemmelse af
endelig Aand, deels i det Andet, hvori den veed sig indlemmet,
erkjende sin Begyndelse, den absolute Aand.  Sig selv kan den ei
erkjende for Aand uden tillige at vide sig selv som endelig Aand:
men som endelig Aand bestaaer den i Foreningen med den absolute Aand, hvori den finder sin Sandhed og Betydning. --- Da
først er Siælen Aand, naar den har udviklet sig til bevidst Eenhed med Gud og Verden; men den er endelig Aand; thi den finder sin Begyndelse som given og har den i Tiden.  Men den
absolute Aand sætter sig selv fra Evighed af, og ved den er den
endelige Aand.

Hermed er Intet udenfra optaget i den dialectiske Tankegang; thi netop fordi den absolute Aand ligger til Grund for den
endelige Aand som dennes Begyndelse, fører ethvert vilkaarlig valgt
Udgangspunct tilbage til denne sande Begyndelse.  Sjælen udvikler sig ikke til et aldeles Andet, men gaaer frem fra sin endelige
Begyndelse til sin absolute Begyndelse.  I den endelige Aands
Begreb, hvortil Sjælen udvikler sig, ligger alt Foreningen eller
Identiteten med den absolute Aand.  Den endelige Aand er i den
absolute Aand og denne i hiin.  Den endelige Aands Grund er
den absolute Aand, og da denne har sin Grund i sig selv, har
% --- p. 445 ---
\opage{445}\setcounter{footnote}{0}%
den ligeledes den endelige Aands Grund i sig, og heraf følger den
endelige Aands Fælledskab med den absolute Aand.  At tænke
Gud som absolut Aand er da ikke noget vilkaarligt Paafund,
men Gud realiserer sig selv i enhver endelig Aand, i det dennes
Udvikling fører til, at den erkjender Gud for den absolute Aand,
under hvis Forudsætning den først selv er Aand.

Da nu den endelige Aand veed sig forbunden med Gud som
sin Grund og sit Væsen, tilskriver den sig nødvendigviis en Præexistens eller rettere sagt en Essents før Existentsen.  Efter sin Existents kan den skabte Aand ei siges at have været evig, men vel
efter sin Essents, fordi den før sin Tid efter sit Væsen har været
evig i Gud. --- Men naar Sjælen efter sin Existents har
havt sin Begyndelse i Tiden, maa den ligeledes efter sin Existents
i Tiden vente sin Ende.  Ligesom den sande Begyndelse for dens
Existents i Tiden maatte søges i Evigheden, saaledes løber ogsaa
Enden paa dens Existents fremad ud i Evigheden.  Dens Begyndelse i Evigheden var den rene Væren og Væsnet; derimod Enden paa Sjælens Existents maa være det Indhold, som har udviklet sig af dens Tilværelse og Væsen.

Er da det Indhold, som har udviklet sig, igjen den rene
Væren? --- Nei.  Den rene Væren er kun den uudviklede Væren: forsaavidt er Enden modsat Begyndelsen, thi Enden er Sjælens til Aand udviklede Væren.  Hvad Sjælen har bragt uudviklet fra Evigheden, det tager den udviklet med sig tilbage til Evigheden.  Heraf er det klart, at Sjælens Udødelighed afhænger af
det Indhold, som den udvikler og aabenbarer i Tiden.  Om dette
Indhold bliver da det næste Spørgsmaal.

Sjælen viser sig i sin Existents først som Individ.  Idet
dernæst Individet fra Udenverdenen reflecteres i sig, lægges det for
Dagen, om det virkelig har et Indre, ɔ: om det ogsaa reflekterer
sig i sig selv.  I saa Fald bliver Individet til Subject, medens
det skrider frem til Bevidsthed om sig selv og om Andet.  Men
% --- p. 446 ---
\opage{446}\setcounter{footnote}{0}%
Modsætningen gjentager sig fremdeles for det paa begge Sider;
det abskiller ikke blot sig selv i Legem og Sjæl, men ligeledes det
Andet i Object og Subject.  Subjectet finder ikke blot Objectet
i sig, men ogsaa Subjectet udenfor sig og sin Sandhed i det
høieste Subject.  Saaledes culminerer Selvbevidstheden i Bevidsthed om Gud, og Sjælen træder herved i Forhold til Gud; men
det staaer i Forhold til ham som til noget Fremmed; Modsætningen har endnu ei fundet sin Opløsning.

Men idet Bevidstheden finder det Andet i sig og sig i det
Andet og omsider sin egen Sandhed i den høieste Bevidsthed, bliver Modsætningen overvunden: Bevidstheden hæver sig til Personlighed, idet den som Aand forenes med den absolute Aand.
Aanden er som saadan baade Subject og Object: den har intet
andet Subject og intet andet Object end sig selv.  Hvad der ei
er Aand, er heller ikke virkeligt, men det bestaaer kun i Momenter af Virkelighed, i Brudstykker af Aanden.  Naturen er kun i
enkelte Dele, hvad Aanden er udeelt: Aanden er Eet og sammenfatter i sig den udvortes Natur, hvilket denne selv ei er i
Stand til.

Den Vei, som Sjælen tager i sin Udvikling, er bleven opfattet som en Tilbagegang i Gud; men den er dog efter sit Indhold væsentlig en Fremgang.  Sjælen vender nemlig ei tilbage til
Essentsen, efter hvilken den var evig i Gud, før den blev skabt;
men den vender tilbage, idet den opfylder sin Bestemmelse og bestandig udvikler den videre.  Derfor er Sjælens Progres ikke
Sluttet med dens Regres i Gud.  Den er og bliver i Kraft af
sin Individualitet skilt fra Gud og alle Skabningens Væsner:
den beholder i Følge sin Subjectivitet Bevidsthed om sig selv,
Gud og alle Væsner.  I Følge sin Personlighed har den endelig
med sin Erkjendelse gjennemtrængt Grændsen og, uden at opgive
sin Endelighed, vedvarer den i det Uendelige.

Det Andet, hvormed den individuelle Aand veed sig identisk, er ikke blot andre skabte Væsner og Gud: derved tænkes og%
% --- p. 447 ---
\opage{447}\setcounter{footnote}{0}%
saa paa det Andet, som hører til den individuelle Aand selv, eller
paa dens Legeme.  Aanden er i denne Henseende Identitet af
Sjæl og Legeme.  Men det udvortes Legeme, som Sjælen kan
skille sig fra, maa vel skjælnes fra det indre Legem, som er den
immanent.  Ved dette sidste adskiller den sig fra Andet og er for
sig, men den forener sig. ogsaa derved med Andet, idet den veed
sig identisk med sin Andetværen og saaledes viser sig som Aand.
Det Indhold, hvortil Sjælen har udviklet sig, er Aanden.  Den
bestaaer deels i Sjælens Identitet med sit Legeme, deels i Aandens Identitet med sit Andet.  Ved den første Bestemmelse bevarer den sin Individualitet og Subjectivitet, ved den sidste sin
personlige Frihed.

Med Personligheden er Frihedens Væsen identisk.  Aandens
Frihed bestaaer nemlig negativ i Negationen af en fremmed Magt,
der kunde bestemme Aanden, positiv i dens egen Magt til Selvbestemmelse.  Vel veed den endelige Aand sig bestemt af den absolute Aand, men den seer i Følge sin Personlighed ei længer i
den nogen fremmed og udvortes Magt.

Men i Friheden ligger en Garantie for Udødeligheden: dette
er Slutningen paa vor Tankegang.  Sjælens Individualitet og
Subjectets Bevidsthed bevares i den endelige Aands Personlighed
og Frihed.  Thi efter Frihedens negative Side forsvinder den udvortes Magt, der kunde true Sjælens Vedvaren med Farer; efter
Frihedens positive Side har den endelige Aand formedelst sin
Forening med den absolute Aand Livets Magt af sit sande Selv
eller af Aanden.“

\medskip

„Ved den endelige Aands immanente Bevægelse fra Sjælen
til Personligheden kastes et nyt Lys paa de tre ovenfor berørte
dogmatiske Beviser for Sjælens Udødelighed, hvorved de kunne
% PRINTED AS IS: V394 q: crop, layer, tesseract -- a MISPRINT, as printed
gjentages i en grundigere Skikkelse.  Det sees let, at den Enkelthed, som ligger til Grund for det metaphysike Beviis, falder sam%
% --- p. 448 ---
\opage{448}\setcounter{footnote}{0}%
men med Sjælens Individualitet, hvorfra den immanente Bevægelse gaaer ud.  Den uendelige Perfectibilitet, som Sjælen efter
det teleologiske Beviis er sig bevidst, svarer til Subjectets Bevidsthed, hvortil Sjælen vaagner.  Den nødvendige Tanke om
Vedvaren, hvorpaa det ontologiske Beviis beraaber sig, svarer ligeledes til Aanden i sin Personlighed, som er Sjælens fuldstændige Begreb.  Rigtignok er Bevidstheden allerede forudsat i det
første Beviis, idet Enkeltheden netop udledes af den; men den
bliver ogsaa kun forudsat, ei udviklet: dens Udvikling falder i
den Tankegang, som hører det andet Beviis til.  Udviklingen
bestaaer nemlig i den Diremtion af Bevidstheden, hvori Subjectets
Selv og Andet træde mod hinanden som adskilte, skjøndt begge
maae erkjendes for Bevidsthedens Indhold.  Ved den sidste Tanke
gjøres Overgang til det tredie Beviis.  For saavidt indeholder
den speculative Tankegang efter sit Indhold intet Nyt, kun Standpunctet er et andet, hvorfra de fragmentariske Beviser nu begribes i deres Sammenhæng.

Ethvert Beviis, der før fremtraadte som afsluttet i sig selv,
gaaer nu over i det næstfølgende.  Enhver funden Sætning viser
sig som usand og opløser sig igjen.  Denne Dialectik er Aarsag
til den Tvivl, som de enkelte Beviser efterlade: Negationen bliver
nemlig følt, men det positive Moment oversees.

Det første er, at Sjælens Enkelthed, hvorfra det metaphysiske Beviis gaaer ud, ikke kan holde sig for Tanken, men negeres i det andet Beviis.  Naar Sjælen først som enkelt ikke kan
skille sig fra sig selv og ikke gaaer ud af sig selv for ei at fortabes, bliver dernæst dens væsentlige Bestemmelse den, at gaae ind
i Andet, thi enten Sjælens Formaal sættes i Erkjendelse eller i
Virksomhed, kræver det bestandig, at Sjælen skal gaae ud af sig
selv og ind i Andet.  Indholdet af det andet Beviis negeres igjen
i det tredie.  Diremtionen i Subject og Object, i Tænkning og
Væren, bliver nemlig opløst derved, at Subjectet har Bevidsthed
% --- p. 449 ---
\opage{449}\setcounter{footnote}{0}%
saavel om Objectet som om sig selv.  Derved fremkommer det
tredie Beviis, i Følge hvilket Eet er i det Andet, og Væren (Vedvaren) tilskrives Tænkningen (Vedvaren).

Vi gjennemgaae Tripliciteten endnu engang med andre Vendinger: 1) Sjælen forbliver i sig selv; den bevæger sig kun i sig
selv og af sig selv.  2) Den bliver sig selv bevidst og bevæger
sig ikke i sig selv, men lever og røres i Gud.  3) Begge Sætninger forsones i Aanden, efterat det fremmede Forhold mellem
den og Gud er ophævet.

Men ligesom Negationen yttrer sig ved enhver Overgang,
saaledes bliver samtidig dermed det positive Moment erkjendt, det
vil sige den Form, hvori den negerede Sætnings Indhold dog
vedligeholder sig.  Saaledes er Sandheden af Enkeltheden udtalt
i Individualiteten, thi Individualiteten er ikke andet end Eenheden,
som i Diremtionen vedligeholder sin Integritet.  Ligeledes er
Sandheden af \emph{Tanken}: at være bestemt til Formaal, udtalt i
Subjectet, som har Bevidsthed saavel om sig selv som om Andet,
som føler sig paa eengang, baade activt og passivt.  Tilsidst bliver
ogsaa den forudsatte umiddelbare Eenhed af Tænkning og Væren
restitueret i Aandens Personlighed.“

Saaledes omtrent har Göschel stræbt at omdanne de sædvanlige
Beviser for Menneskets Udødelighed, for at de skulle kunne stemme
med den Philosophie, som han giver Navn af den nyere Tids.
Skulde nu Nogen, der kjender de her brugte Formler, efterat
have læst det her givne korte Grundrids, staae i Tvivl, om Forfatteren ikke taler om en anden Udødelighed, end den der antages
% PRINTED AS IS: V401 q: crop, layer, tesseract -- a MISPRINT, as printed
i det christelige Sanfund, maa jeg føie den Bemærkning til, at
Göschel med mange flere tildeels utvetydige Vendinger anvender
sit Udødelighedsbegreb paa concrete Livsforhold.  Han meddeler
f. Ex. endeel Aphorismer over Legemets Opstandelse, Erindring
af Tilfældigheder og Personers Gjenkjendelse.  Der bliver saaledes
ikke mindste Tvivl om, at han virkelig staaer i den Formening,
% --- p. 450 ---
\opage{450}\setcounter{footnote}{0}%
at han er i Stand til at bringe sin speculative Philosophies Resultater i fuldkommen Harmonie med den christelige Traditions
Indhold.  Men næppe vil nogen skoleret Hegelianer give ham
Medhold i den Tro, at et saadant Kirkespiir lader sig anbringe
paa den hegelske Bygning, uden at det underkastes en gjennemgribende Forvandling.  Der har ogsaa fra flere Sider været gjort
Indsigelser derimod af Hegels Disciple, og man har opkastet
det Spørgsmaal, paa hvad Sted i Systemet (eller Schemaet
dertil) denne Cyclus af Begreber skal have sin Plads.  Det er
overhovedet temmelig langt fra, at Göschels Fremstilling kan
kaldes stræng methodisk: den udmærker sig tvertimod ved en usædvanlig broget Forvirring.  Snart træffer man, i den Bog Talen
er om, et Brudstykke, der holder sig strængt i hegelske Former,
snart føres man ind i et overdreven populært Foredrag; men
hvad der især gjør Afhandlingens Læsning besværlig, er den trættende, gammeldags Samling af Citater eller Mottoer i allehaande
Sprog, saavel i bunden som i ubunden Stiil, hvorved Fremstillingen aldeles uden Skjønsomhed er afbrudt.  Forfatteren har saaledes ingenlunde havt den Fordeel af at gjennemgaae Hegels Skole,
som dennes Tilhængere jævnlig ønske sig tillykke med, nemlig den
Bevidsthedens Tugt, som forebygger en uordentlig omspringende
Tankegang.

Den, der har Vedholdenhed til at følge Forfatteren gjennem
alle hans indholdsrige Fragmenter, vil vistnok overraskes ved mange
dybsindige Bemærkninger saavel til Theologien som den høiere
Psychologie.  Men den Grundanskuelse, som udtales i Skriftet,
maa snarere betragtes som et af de mange Fragmenter til et nyt
System, som Nutiden seer, end som en Erobring, der uden at
afstedkomme Revolution kan indlemmes i den hegelske Philosophie. ---
Endeel af Göschels Grundtanker ere uden al Tvivl hentede fra
den yngre Fichtes Tankekreds, uagtet Forfatteren i flere Aar har
befundet sig i en polemisk Stilling til denne Philosoph.

% --- p. 451 ---
\opage{451}\setcounter{footnote}{0}%
Göschel har paa de sidste Sider af sit Skrift fremført nogle
Citater, som han kalder Ord af Hegel over Sjælens Udødelighed.
Men de faae alle tilsammen, naar de læses i den Forbindelse,
hvoraf de ere løsrevne, heel anden Betydning.  Mærkeligt er det,
at i det Mindste een af disse Tirader oftere har været udskreven
netop som haandgribeligt Beviis for Hegels Kjætteri i dette Punct\footnote{Inden jeg forlader Göschels Afhandling om Udødeligheden tillader
jeg mig deraf at afskrive hans Dom over et lidet indholdsrigt Program over samme Gjenstand, som Prof. Sibbern for nogle Aar
siden har udstedet ved Universitetet.
¶ „Besondere Beachtung verdient Freder. Christ. Sibbern de
\textit{præexistentia, genesi et immortalitate animæ. Havniæ 1823.}
Dieses Programm enthält mehr als manches Buch.  Der physiologischen Erörterung wird hier außer der metaphysischen und moralischen Beweisführung auch der Gedanke der Ewigkeit, \textit{persuasio}
\textit{immortalitatis}, als Grundlage aller Demonstration vorausgeschickt.
Beilaüfig wird der Begriff der Perfectibilität mit Beziehung auf
Fr. Baader: \textit{sur la notion du temps} von der abstracten Vorstellung der schlechten Unendlichkeit befreit, \textit{quæ hominem per}
\textit{secula seculorum jubet currere ad finem, quem nunquam}
\textit{sit consecuturus.}“}.

\begin{center}\rule{3em}{0.4pt}\end{center}

\begin{center}XI.\end{center}

De ovenomtalte Skrifter om Udødeligheden ere, saavidt jeg
veed, de mærkværdigste Frembringelser i vor Tids rige Literatur
over denne Materie; i al Fald er det dem der meest have været
Gjenstand for den literære Conversation.  De fleste Flyveskrifter
over dette Dagens Yndlingsthema, hvilke jeg har havt Adgang
til, indeholde meer eller mindre heldige Gjentagelser af de sædvanlige Beviser uden nogen mærkelig Variation og uden fremragende
Veltalenhed.  Af de phantastiske Skrifter --- for ikke at tale om
de rene Spøgelseshistorier, hvis Literatur er betydelig forøget i
denne Vinter, --- nævner jeg som en Synderlighed et lidet Skrift
% --- p. 452 ---
\opage{452}\setcounter{footnote}{0}%
af \textit{Dr.} Mises\footnote{Das Büchlein vom Leben nach dem Tode, von \textit{Dr.} Mises. Dresden 1836.}. Det dreier sig om et Punct, som Tanken og
Grubleriet idelig maa føres tilbage til, nemlig de endelige Individers Begrændsning mod hverandre og tilkommende mulige Gjennemsigtighed for hverandre. Han fremfører med stor Dristighed
en egen Hypothes, i Følge hvilken de afdøde Mennesker med
Selvbevidsthed fortsætte deres Existents i de paa Jorden levendes
Tænkning og Villie.

Den meest uvidenskabelige Maade, hvorpaa man i den senere
Tid har villet virke for Udødelighedslæren, er den, at man har
udgivet Excerptsamlinger af andre Tiders Skrifter over denne
Gjenstand. Jeg tænker herved især paa de Hefter, som \textit{Dr.}
Hubert-Becker\footnote{Denne Hubert-Becker, som er en Tilhænger af Schellings nye System, har ogsaa udgivet et særeget Skrift om Udødeligheden, hvilket
indeholder en Kritik over Göschels Beviisførelse. Hans Mening
% PRINTED AS IS: V412 q: reader, crop, tesseract -- a turned n, as printed
om Sagen er den, at Nutidens Philosopher gjøre bedst i med Taalmodighed at oppebie Bekjendtgjørelseu af Schellings senere philosophiske Arbeider, hvori der skal blive given tilfredsstillende Besked
angaaende dette Problem. Han anseer det for en Indiscretion at
røbe hvad Schelling har betroet sine Tilhørere deraf: dog meddeler
han efter dertil indhentet Tilladelse nogle faa Linier derom i en
Anmærkning. Men, som man let kan forestille sig, bliver derved
ikke nogen klar Tanke fremstillet.} har udgivet med Uddrag af Skrifter fra det 16de
og 17de Aarhundrede „om Sjælens Tilstand efter Døden.“ Overhovedet ere deslige Reductioner af philosophiske Værker efter en
formindsket Maalestok et Slags Arbeider, som meget sjælden svare
til deres Bestemmelse. Især kan Referenten let bringe sin Læser
i en ubehagelig Svimmel, naar han byder ham flere paa hinanden følgende Extracter af heterogene Systemer. Forskjellige Tankegange sættes nemlig derved saaledes i Bevægelse hos Læseren
fra modsatte Synspunkter, at de forstyrre hinandens Udvikling
% --- p. 453 ---
\opage{453}\setcounter{footnote}{0}%
ligesom de excentriske Cirkler, der dannes i stille Vand, naar
flere Stene kastes deri ved Siden af hinanden. --- Men min Frygt
for at gjøre mig selv skyldig i et saadant Misgreb bringer mig til at
opgive de Meddelelser af mindre betydelige Skrifters Indhold, som
% PRINTED AS IS: R11 z: reader, crop, tesseract -- a MISPRINT for endnu, as printed
jeg endnn saae mig i Stand til at fremkomme med.

Det gjælder om ethvert isoleret Beviis for Menneskets Udødelighed, at det muligen kan aabne Prospectet til en dyb og grundig Verdensanskuelse, men det er ofte tilfældigt, om Læseren føier
den Integration til et saadant Tankefragment, som virkelig har
ligget i Forfatterens Mening. Dette gjælder især om de Skrifter, der referere sig til et System, med hvis Aand de ikke kunne
bestaae; thi deres strænge Videnskabelighed er da kun et Skin og
tilveiebringer ikke virkelig den dybere Forstaaelse, som de gjøre
Fordring paa. Det egentlig strænge Beviis kan kun føres i et
fuldstændig udtalt System, men da findes i et saadant System
en ubestemt Mangfoldighed af Beviser for Udødelighed. Thi, naar
en saadan følgerig Erkjendelse oprindelig ligger i Anskuelsen, der
udfoldes i Systemet, maae dens Conseqventser vise sig overalt i
Systemets Organisation.

\begin{flushright}Poul Møller.\end{flushright}

\end{document}
